% Lesson 8 — Grammar: Prepositions and prepositional phrases (place, time, manner and direction) — Anglais 1ère A — Cours signé OnBuch+
% Programme officiel MINESEC. Compiler avec : tectonic ENA08-grammar-prepositions-and-prepositional-phrases-place-ti.tex
\newcommand{\DOCMATIERE}{Anglais}
\newcommand{\DOCNIVEAU}{1ère A}
\newcommand{\DOCMODULE}{Chapter 1 : Family and Social Life}
\newcommand{\DOCLECON}{ENA08}
\newcommand{\DOCTITRE}{Lesson 8 — Grammar: Prepositions and prepositional phrases (place, time, manner and direction)}
% OnBuch+ — préambule « cours signés » ENRICHI (compilé avec tectonic / XeLaTeX)
% Système visuel « Pop » d'OnBuch+ : fond crème (jamais blanc), bordures encre
% épaisses, ombres « dures » décalées, titres Archivo Black en capitales.
% Version MATHÉMATIQUES : graphes (pgfplots/tikz), boîtes pédagogiques
% (définition, propriété, méthode, attention, le savais-tu, exercice résolu,
% exercice, TP, bilan), page de garde et signature OnBuch+ ancrée en bas.
%
% Macros attendues dans main.tex AVANT \input{preamble} :
%   \DOCMATIERE  \DOCNIVEAU  \DOCMODULE  \DOCLECON (numéro)  \DOCTITRE
\documentclass[11pt]{article}

\usepackage{fontspec}
\usepackage[english,french]{babel} % français = langue principale ; l'anglais via \eng{..} / textebox
\usepackage[a4paper,margin=2cm,top=3.4cm,bottom=2.8cm,headheight=1.4cm,footskip=1.3cm]{geometry}
\usepackage{amsmath,amssymb,amsfonts}
\usepackage{mathtools} % \xmapsto, \coloneqq... souvent utilisés par les modèles
\usepackage{cancel} % simplifications barrées (bilans de forces)
% \abs{z} (module, valeur absolue) et \norm{u} (norme), utilisés par les modèles
\providecommand{\abs}{}\let\abs\relax\DeclarePairedDelimiter{\abs}{\lvert}{\rvert}
\providecommand{\norm}{}\let\norm\relax\DeclarePairedDelimiter{\norm}{\lVert}{\rVert}
\usepackage[table]{xcolor} % [table] : fournit aussi \rowcolors (alternance de lignes)
\usepackage{pagecolor}
\usepackage{tikz}
\usetikzlibrary{arrows.meta,positioning,calc,shapes.geometric,shapes.misc,shapes.arrows,decorations.pathmorphing,decorations.pathreplacing,decorations.markings,patterns,shadows,mindmap,trees,fit,backgrounds,matrix,intersections,angles,quotes}
\usepackage{pgfplots}
\usepackage[version=4]{mhchem} % équations nucléaires \ce{^{235}_{92}U}
\usepackage[european,siunitx]{circuitikz} % schémas électriques
\pgfplotsset{compat=1.18}
\usepgfplotslibrary{fillbetween,groupplots} % aires entre courbes (calcul intégral)
\usepackage{enumitem}
\usepackage{fancyhdr}
% [most] charge toutes les bibliothèques tcolorbox (listings, minted, raster,
% documentation...) et gonfle inutilement la mémoire TeX sur un document long
% (~300 boîtes) : on ne charge que ce qui est réellement utilisé.
\usepackage[skins,breakable]{tcolorbox}
\usepackage{graphicx}
\usepackage{colortbl}
\usepackage{booktabs}
\usepackage{tabularx}
\usepackage{diagbox} % case d'en-tête diagonale des tableaux à double entrée
% Colonnes extensibles centrée / gauche / droite (C, L, R), souvent utilisées par les modèles
\newcolumntype{C}{>{\centering\arraybackslash}X}
\newcolumntype{L}{>{\raggedright\arraybackslash}X}
\newcolumntype{R}{>{\raggedleft\arraybackslash}X}
\usepackage{array}
\usepackage{multicol}
\usepackage{multirow}
\usepackage{siunitx}
% parse-numbers=false : accepte aussi \SI{2,5 \times 10^{-3}}{...}
\sisetup{locale=FR,output-decimal-marker={,},group-minimum-digits=4,parse-numbers=false}
\DeclareSIUnit\ohm{\ensuremath{\symup{\Omega}}} % U+2126 absent de la police
\DeclareSIUnit\division{div} % réglages d'oscilloscope (ms/div, V/div)
\DeclareSIUnit\tr{tr} % tours (vitesse de rotation, ex. \SI{1500}{\tr\per\minute})
\DeclareSIUnit\molar{M} % concentration molaire (ex. \SI{3}{\molar}, \SI{1}{\micro\molar})
\DeclareSIUnit\an{an} % année (le modèle confond parfois avec \year, un compteur TeX) — ex. \SI{5}{\kilo\meter\per\an}
\DeclareSIUnit\FCFA{FCFA} % franc CFA (ex. \SI{500}{\FCFA\per\kilo\gram})
\DeclareSIUnit\degreeBrix{\textdegree Brix} % teneur en sucre d'un sirop/jus (réfractométrie)
\DeclareSIUnit\month{mois} % le modèle utilise parfois \month, un compteur TeX, comme unité de durée
\usepackage{qrcode}
% \degree : commande fréquemment utilisée par les modèles pour un angle ou une
% température en dehors de \SI{}{} (non fournie par les paquets chargés).
\providecommand{\degree}{^{\circ}}
% \qed (fin de démonstration), utilisé par les modèles sans amsthm
\providecommand{\qed}{\hfill$\square$}
% \barre : double barre des valeurs interdites dans un tableau de variations (tkz-tab)
\providecommand{\barre}{\ensuremath{\|}}
% environnement proof (amsthm non chargé) utilisé par les modèles
\ifdefined\proof\else\newenvironment{proof}[1][Démonstration]{\par\noindent\textit{#1.}\ }{\qed\par}\fi
\usepackage{lastpage}
\usepackage{hyperref}
\hypersetup{hidelinks}

% ============================================================
% Palette « Pop » OnBuch+ — mêmes hex que la classe P (plus_ui.dart)
% ============================================================
\definecolor{poporange}{HTML}{F59321}
\definecolor{popdark}{HTML}{E07A0C}
\definecolor{popcream}{HTML}{FFF6E8}
\definecolor{popink}{HTML}{1C1714}
\definecolor{poppurple}{HTML}{7A5AE0}
\definecolor{popgreen}{HTML}{2E9E6A}
\definecolor{poppink}{HTML}{E0457B}
\definecolor{popblue}{HTML}{2F7DE1}
\definecolor{popgold}{HTML}{C98A12}
\definecolor{popmuted}{HTML}{8A7F76}
\definecolor{popline}{HTML}{EADFD2}
\definecolor{popgray}{HTML}{8A7F76}
\colorlet{popgrey}{popgray}
% Alias tolérants (le modèle nomme parfois les couleurs différemment)
\colorlet{popwhite}{white}
\colorlet{popblack}{popink}
\colorlet{popred}{poppink}
\colorlet{popyellow}{popgold}
% Teintes claires (fonds de tableaux/encadrés) — le modèle nomme parfois
% les couleurs avec un suffixe L (ex. popblueL) sans les avoir définies.
\colorlet{poporangeL}{poporange!20}
\colorlet{popdarkL}{popdark!20}
\colorlet{poppurpleL}{poppurple!20}
\colorlet{popgreenL}{popgreen!20}
\colorlet{poppinkL}{poppink!20}
\colorlet{popblueL}{popblue!20}
\colorlet{popgoldL}{popgold!20}
\colorlet{popredL}{popred!20}
\colorlet{popgrayL}{popgray!20}
% Teintes claires pour fonds de tableaux / zones
\colorlet{popblueL}{popblue!10!white}
\colorlet{poporangeL}{poporange!14!white}
\colorlet{popgreenL}{popgreen!12!white}
\colorlet{poppurpleL}{poppurple!12!white}
\colorlet{poppinkL}{poppink!12!white}
\colorlet{popgoldL}{popgold!14!white}

\pagecolor{popcream}

% ============================================================
% Typographie
% ============================================================
\defaultfontfeatures{Ligatures=TeX}
\setmainfont{PlusJakartaSans-Medium.ttf}[Path=../../../fonts/,BoldFont=PlusJakartaSans-Bold.ttf]
% Maths en sans-serif (Fira Math) avec chiffres et lettres droites de Plus
% Jakarta, pour que les formules se fondent dans le texte.
\usepackage[math-style=ISO,bold-style=ISO]{unicode-math}
\setmathfont{FiraMath-Regular.otf}[Path=../../../fonts/]
\setmathfont{PlusJakartaSans-Medium.ttf}[Path=../../../fonts/,range={up/num,up/{Latin,latin},\mathup}]
\usepackage{icomma} % virgule décimale sans espace en mode math
% Caractères mathématiques Unicode absents des polices de texte (Plus Jakarta,
% Archivo Black) : rendus par leur équivalent mathématique, en texte comme en
% formule — sinon ils s'affichent en carré vide (ex. « Arithmétique dans ℤ »).
\usepackage{newunicodechar}
\newunicodechar{ℝ}{\ensuremath{\mathbb{R}}}
\newunicodechar{ℤ}{\ensuremath{\mathbb{Z}}}
\newunicodechar{ℂ}{\ensuremath{\mathbb{C}}}
\newunicodechar{ℕ}{\ensuremath{\mathbb{N}}}
\newunicodechar{ℚ}{\ensuremath{\mathbb{Q}}}
\newunicodechar{𝔻}{\ensuremath{\mathbb{D}}}
\newunicodechar{α}{\ensuremath{\alpha}}
\newunicodechar{β}{\ensuremath{\beta}}
\newunicodechar{γ}{\ensuremath{\gamma}}
\newunicodechar{δ}{\ensuremath{\delta}}
\newunicodechar{ε}{\ensuremath{\varepsilon}}
\newunicodechar{θ}{\ensuremath{\theta}}
\newunicodechar{λ}{\ensuremath{\lambda}}
\newunicodechar{μ}{\ensuremath{\mu}}
\newunicodechar{σ}{\ensuremath{\sigma}}
\newunicodechar{τ}{\ensuremath{\tau}}
\newunicodechar{φ}{\ensuremath{\varphi}}
\newunicodechar{ψ}{\ensuremath{\psi}}
\newunicodechar{ω}{\ensuremath{\omega}}
\newunicodechar{Σ}{\ensuremath{\Sigma}}
% majuscules produites par \MakeUppercase dans les titres (α-aminés -> Α)
\newunicodechar{Α}{\ensuremath{\alpha}}
\newunicodechar{Β}{\ensuremath{\beta}}
\newunicodechar{Γ}{\ensuremath{\Gamma}}
\newunicodechar{Δ}{\ensuremath{\Delta}}
\newunicodechar{Ε}{\ensuremath{\varepsilon}}
\newunicodechar{Θ}{\ensuremath{\Theta}}
\newunicodechar{Λ}{\ensuremath{\Lambda}}
\newunicodechar{Μ}{\ensuremath{\mu}}
\newunicodechar{Τ}{\ensuremath{\tau}}
\newunicodechar{Φ}{\ensuremath{\Phi}}
\newunicodechar{Ψ}{\ensuremath{\Psi}}
\newunicodechar{Ω}{\ensuremath{\Omega}}
\newunicodechar{↦}{\ensuremath{\mapsto}}
\newunicodechar{∈}{\ensuremath{\in}}
\newunicodechar{∉}{\ensuremath{\notin}}
\newunicodechar{∧}{\ensuremath{\wedge}}
\newunicodechar{⇔}{\ensuremath{\Leftrightarrow}}
\newunicodechar{⟺}{\ensuremath{\Longleftrightarrow}}
\newunicodechar{⇒}{\ensuremath{\Rightarrow}}
\newunicodechar{⟹}{\ensuremath{\Longrightarrow}}
\newunicodechar{∘}{\ensuremath{\circ}}
\newunicodechar{∪}{\ensuremath{\cup}}
\newunicodechar{∩}{\ensuremath{\cap}}
\newunicodechar{≡}{\ensuremath{\equiv}}
\newunicodechar{⊂}{\ensuremath{\subset}}
\newunicodechar{⊥}{\ensuremath{\perp}}
\newunicodechar{∃}{\ensuremath{\exists}}
\newunicodechar{∀}{\ensuremath{\forall}}
\newunicodechar{∛}{\ensuremath{\sqrt[3]{\,}}}
\newunicodechar{∜}{\ensuremath{\sqrt[4]{\,}}}
\newunicodechar{⃗}{\ensuremath{{}^{\to}}}
\newunicodechar{ⁿ}{\ifmmode^{n}\else\textsuperscript{\ensuremath{n}}\fi}
\newunicodechar{ˣ}{\ifmmode^{x}\else\textsuperscript{\ensuremath{x}}\fi}
\newunicodechar{ᵇ}{\ifmmode^{b}\else\textsuperscript{\ensuremath{b}}\fi}
\newunicodechar{ʳ}{\ifmmode^{r}\else\textsuperscript{\ensuremath{r}}\fi}
\newunicodechar{ᵘ}{\ifmmode^{u}\else\textsuperscript{\ensuremath{u}}\fi}
\newunicodechar{ʸ}{\ifmmode^{y}\else\textsuperscript{\ensuremath{y}}\fi}
\newunicodechar{ᵃ}{\ifmmode^{a}\else\textsuperscript{\ensuremath{a}}\fi}
\newunicodechar{ᵉ}{\ifmmode^{e}\else\textsuperscript{\ensuremath{e}}\fi}
\newunicodechar{ᵏ}{\ifmmode^{k}\else\textsuperscript{\ensuremath{k}}\fi}
\newunicodechar{ᵖ}{\ifmmode^{p}\else\textsuperscript{\ensuremath{p}}\fi}
\newunicodechar{ᴺ}{\ifmmode^{N}\else\textsuperscript{\ensuremath{N}}\fi}
\newunicodechar{ᶜ}{\ifmmode^{c}\else\textsuperscript{\ensuremath{c}}\fi}
\newunicodechar{ʰ}{\ifmmode^{h}\else\textsuperscript{\ensuremath{h}}\fi}
\newunicodechar{ᵠ}{\ifmmode^{q}\else\textsuperscript{\ensuremath{q}}\fi}
\newunicodechar{ₙ}{\ifmmode_{n}\else\textsubscript{\ensuremath{n}}\fi}
\newunicodechar{ₐ}{\ifmmode_{a}\else\textsubscript{\ensuremath{a}}\fi}
\newunicodechar{ᵢ}{\ifmmode_{i}\else\textsubscript{\ensuremath{i}}\fi}
\newunicodechar{ᵣ}{\ifmmode_{r}\else\textsubscript{\ensuremath{r}}\fi}
\newunicodechar{ₖ}{\ifmmode_{k}\else\textsubscript{\ensuremath{k}}\fi}
\newunicodechar{ᵦ}{\ifmmode_{\beta}\else\textsubscript{\ensuremath{\beta}}\fi}
\newunicodechar{ₓ}{\ifmmode_{x}\else\textsubscript{\ensuremath{x}}\fi}
\newfontfamily\popdisplay{ArchivoBlack-Regular.ttf}[Path=../../../fonts/]
\newfontfamily\poplogo{SpaceGrotesk-Bold.ttf}[Path=../../../fonts/]
\newfontfamily\popsemi{PlusJakartaSans-SemiBold.ttf}[Path=../../../fonts/]
\newfontfamily\popxbold{PlusJakartaSans-ExtraBold.ttf}[Path=../../../fonts/]
\newfontfamily\popxboldls{PlusJakartaSans-ExtraBold.ttf}[Path=../../../fonts/,LetterSpace=3.0]

\setlist[enumerate]{leftmargin=1.7em,itemsep=5pt,topsep=4pt}
\setlist[itemize]{leftmargin=1.7em,itemsep=4pt,topsep=4pt,label={\color{poporange}\textbullet}}
\setlength{\parindent}{0pt}
\setlength{\parskip}{5pt}
\renewcommand{\arraystretch}{1.25}

% pgfplots : style Pop par défaut
\pgfplotsset{
  every axis/.append style={axis line style={popink,line width=1pt},tick style={popink},
    grid style={popline,line width=0.5pt},label style={font=\small},tick label style={font=\footnotesize}},
  popcurve/.style={poporange,line width=1.8pt,no markers,smooth},
}
% Même style disponible dans un \draw TikZ simple (hors pgfplots)
\tikzset{popcurve/.style={poporange,line width=1.8pt}}
\tikzset{popgrid/.style={popline,very thin}}
\colorlet{popcurve}{poporange} % le nom du style est parfois utilisé comme couleur

% ============================================================
% Pill / squiggle
% ============================================================
\newcommand{\poppill}[3]{%
  \tikz[baseline=-0.55ex]\node[draw=popink,line width=1.2pt,rounded corners=2.6pt,%
    fill=#2,inner xsep=6.5pt,inner ysep=3pt]{%
    \popxboldls\fontsize{7}{7}\selectfont\color{#3}\MakeUppercase{#1}};%
}
\newcommand{\popsquiggle}[1]{%
  \tikz[baseline=-0.4ex]{
    \draw[#1,line width=2.6pt,line cap=round]
      plot [smooth] coordinates {(0,0.09) (0.34,0.01) (0.68,0.21) (1.02,0.01) (1.36,0.10)};
  }%
}
% Mot-clé surligné (marqueur orange sous le texte)
\newcommand{\cle}[1]{{\popxbold\color{popdark}#1}}

% ============================================================
% En-tête / pied
% ============================================================
\pagestyle{fancy}
\fancyhf{}
\renewcommand{\headrulewidth}{0pt}
\renewcommand{\footrulewidth}{0pt}
\lhead{\raisebox{-0.15ex}{\poplogo\fontsize{13}{13}\selectfont\color{popink}OnBuch\color{poporange}+}}
\rhead{\poppill{\DOCMATIERE\ \textbullet\ \DOCNIVEAU\ \textbullet\ Leçon \DOCLECON}{white}{popink}}
\lfoot{\popxbold\fontsize{7.6}{7.6}\selectfont\color{popmuted}\MakeUppercase{Cours signé OnBuch+}\ \textbullet\ {\popsemi\DOCTITRE}}
\rfoot{\tikz[baseline=-0.6ex]\node[rounded corners=8.5pt,fill=popink,minimum height=17pt,minimum width=17pt,inner xsep=5pt,inner sep=0]{\popxbold\fontsize{8}{8}\selectfont\color{popcream}\thepage\,/\,\pageref*{LastPage}};}

% ============================================================
% Titres
% ============================================================
\newcommand{\chaptitre}[2]{%
  \begin{tcolorbox}[enhanced,colback=poporange,colframe=popink,boxrule=2.6pt,arc=11pt,
      drop shadow=popink,left=15pt,right=15pt,top=12pt,bottom=11pt,before skip=3pt,after skip=18pt]
    \poppill{#2}{popink}{popcream}\par\vspace{9pt}
    {\raggedright\hyphenpenalty=10000\exhyphenpenalty=10000\popdisplay\fontsize{22}{24}\selectfont\color{popcream}\MakeUppercase{#1}\par}
  \end{tcolorbox}
}
% Section numérotée : pastille orange avec le numéro + titre Archivo
\newcounter{popsec}
\newcommand{\coursec}[1]{%
  \stepcounter{popsec}\setcounter{popsub}{0}%
  \par\vspace{16pt}\noindent
  \begin{minipage}{\linewidth}
  \tikz[baseline=-0.7ex]\node[draw=popink,line width=1.6pt,fill=poporange,rounded corners=4pt,minimum size=22pt,inner sep=0]{\popdisplay\fontsize{12}{12}\selectfont\color{popcream}\thepopsec};%
  \hspace{8pt}{\popdisplay\fontsize{15}{17}\selectfont\color{popink}\MakeUppercase{#1}}\par
  \vspace{4pt}\noindent{\color{poporange}\rule{36pt}{3.6pt}}
  \end{minipage}\par\nopagebreak\vspace{8pt}
}
\newcounter{popsub}
\newcommand{\courssub}[1]{%
  \stepcounter{popsub}%
  \par\vspace{9pt}\noindent{\popxbold\fontsize{11.5}{13}\selectfont\color{popink}{\color{poporange}\thepopsec.\thepopsub}\hspace{6pt}#1}\par\nopagebreak\vspace{4pt}
}

% ============================================================
% Boîtes pédagogiques — même recette : fond blanc, bordure épaisse colorée,
% ombre dure encre, badge plein. Titre optionnel [#1].
% ============================================================
\tcbset{popbase/.style={breakable,enhanced,colback=white,boxrule=2.3pt,arc=9pt,
  shadow={3pt}{-3pt}{0pt}{popink},left=14pt,right=14pt,top=11pt,bottom=11pt,before skip=11pt,after skip=15pt,
  fontupper=\popsemi\fontsize{10.6}{14.5}\selectfont\color{popink},
  fontlower=\popsemi\fontsize{10.6}{14.5}\selectfont\color{popink}}}
% Coche dessinée (le glyphe ✓ n'existe pas dans les polices du document)
\newcommand{\popcheck}[1]{\tikz[baseline=-0.5ex]\draw[#1,line width=1.6pt,line cap=round,line join=round](-0.13,0.0)--(-0.04,-0.09)--(0.13,0.1);}
\providecommand{\coche}{\popcheck{popgreen}}
\providecommand{\resultat}[1]{\par\smallskip\noindent\fcolorbox{popgreen}{popgreenL}{\parbox{\dimexpr\linewidth-2\fboxsep-2\fboxrule\relax}{\textbf{Résultat.} #1}}\par\smallskip}
\newcommand{\popboxhead}[3]{\poppill{#1}{#2}{white}\ifx\relax#3\relax\else\hspace{6pt}{\popxbold\fontsize{10.6}{12}\selectfont\color{popink}#3}\fi\par\vspace{7pt}}

\newtcolorbox{aretenir}[1][]{popbase,colframe=popgreen,before upper={\popboxhead{À retenir}{popgreen}{#1}}}
% Texte en anglais (document de lecture, dialogue, extrait) : cadre
% d'étude. Un titre optionnel (source, auteur) s'affiche dans l'en-tête.
\newtcolorbox{textebox}[1][]{popbase,colframe=popblue,before upper={\popboxhead{Text}{popblue}{#1}\begin{otherlanguage}{english}},after upper={\end{otherlanguage}}}
% Marques ✓ / ✗ (forme correcte / fautive) souvent utilisées par les modèles
\providecommand{\cross}{\textcolor{poppink}{\ensuremath{\times}}}
\providecommand{\xmark}{\cross}\providecommand{\crossmark}{\cross}\providecommand{\cmark}{\ensuremath{\checkmark}}
% Ligne de réponse / justification à compléter (inventée par les modèles)
\providecommand{\justification}[1]{\par\noindent\textit{Justification~:} \dotfill\par}
% \textipa (package tipa non chargé) : la police ne couvre pas l'API -> simple passage
\providecommand{\textipa}[1]{#1}
% Soulignement (ulem non chargé) : \uline, \ul
\providecommand{\uline}[1]{\underline{#1}}\providecommand{\ul}[1]{\underline{#1}}
% \glossaire{\item ...} inventée par les modèles : liste de glossaire
\providecommand{\glossaire}[1]{\begin{itemize}#1\end{itemize}}
% \alert (beamer) inventée par les modèles : texte mis en évidence
\providecommand{\alert}[1]{\textbf{\textcolor{poppink}{#1}}}
% variantes ASCII d'ordinaux écrites en macro
\providecommand{\iere}{\textsuperscript{re}}\providecommand{\ieme}{\textsuperscript{e}}\providecommand{\ere}{\textsuperscript{re}}\providecommand{\eme}{\textsuperscript{e}}\providecommand{\er}{\textsuperscript{er}}\providecommand{\ie}{\textsuperscript{e}}
% Ordinaux français écrits en macro par les modèles : 1\ière, 3\ième
\newcommand{\ière}{\textsuperscript{re}}\newcommand{\ième}{\textsuperscript{e}}
% \exemple (inventée par les modèles) : étiquette « Exemple : »
\providecommand{\exemple}{\textbf{Exemple~:}\ }
% \square utilisable aussi hors mode math (cases à cocher)
\AtBeginDocument{\let\squareMath\square\renewcommand{\square}{\ensuremath{\squareMath}}}
% Mot ou phrase en anglais à l'intérieur d'un paragraphe français
\newcommand{\eng}[1]{\ifmmode\text{\textit{#1}}\else\begingroup\selectlanguage{english}\itshape #1\endgroup\fi}
\let\esp\eng % alias toléré
% flèches d'intonation inventées par les modèles
\providecommand{\textsearrow}{}\renewcommand{\textsearrow}{\ensuremath{\searrow}}\providecommand{\textnearrow}{}\renewcommand{\textnearrow}{\ensuremath{\nearrow}}
% \fr{..} : mot français (inventé par les modèles)
\providecommand{\fr}[1]{\textit{#1}}
\newtcolorbox{exemplebox}[1][]{popbase,colframe=poporange,before upper={\popboxhead{Exemple}{poporange}{#1}}}
\newtcolorbox{definition}[1][]{popbase,colframe=popblue,before upper={\popboxhead{Définition}{popblue}{#1}}}
\newtcolorbox{propriete}[1][]{popbase,colframe=poppurple,before upper={\popboxhead{Propriété}{poppurple}{#1}}}
\newtcolorbox{methode}[1][]{popbase,colframe=popgold,before upper={\popboxhead{Méthode}{popgold}{#1}}}
\newtcolorbox{attention}[1][]{popbase,colframe=poppink,before upper={\popboxhead{Attention}{poppink}{#1}}}
\newtcolorbox{experience}[1][]{popbase,colframe=popblue,colback=popblueL,before upper={\popboxhead{Expérience}{popblue}{#1}}}
% « Le savais-tu ? » : carte violette pleine (contexte, culture, Cameroun)
\newtcolorbox{savaistu}[1][]{popbase,colback=poppurple,colframe=popink,
  fontupper=\popsemi\fontsize{10.4}{14.2}\selectfont\color{white},
  before upper={\poppill{Le savais-tu\,?}{white}{poppurple}\ifx\relax#1\relax\else\hspace{6pt}{\popxbold\color{white}#1}\fi\par\vspace{7pt}}}
% Exercice résolu : énoncé en haut, \tcblower puis la solution sur fond crème
\newcounter{popexo}\newcounter{popexr}
\newtcolorbox{exoresolu}[1][]{popbase,colframe=poporange,colbacklower=popcream,
  segmentation style={popink,line width=1.2pt,dashed},
  before upper={\stepcounter{popexr}\popboxhead{Exercice résolu \thepopexr}{poporange}{#1}},
  before lower={\poppill{Solution}{popink}{popcream}\par\vspace{6pt}}}
% Exercice (non résolu en place) — la correction est dans la partie Corrigés
\newtcolorbox{exercice}[1][]{popbase,colframe=popink,
  before upper={\stepcounter{popexo}\popboxhead{Exercice \thepopexo}{popink}{#1}}}
% Corrigé
\newtcolorbox{corrige}[1][]{popbase,colframe=popgreen,colback=popgreenL,
  before upper={\popboxhead{Corrigé}{popgreen}{#1}}}
% Objectifs / prérequis / bilan
\newtcolorbox{objectifs}{popbase,colframe=popink,colback=poporangeL,
  before upper={\popboxhead{Objectifs de la leçon}{popink}{}}}
\newtcolorbox{prerequis}{popbase,colframe=popink,colback=popblueL,
  before upper={\popboxhead{Prérequis}{popblue}{}}}
\newtcolorbox{bilan}{popbase,colframe=popink,colback=white,boxrule=2.8pt,
  before upper={\popboxhead{Fiche bilan}{poporange}{L'essentiel à connaître pour l'examen}}}
% Figure encadrée avec légende
\newtcolorbox{popfigure}[1][]{enhanced,breakable=false,colback=white,colframe=popline,boxrule=1.4pt,arc=8pt,
  left=10pt,right=10pt,top=10pt,bottom=8pt,before skip=10pt,after skip=12pt,halign=center,
  lower separated=false,halign lower=center,
  fontlower=\popsemi\fontsize{9}{11.5}\selectfont\color{popmuted},
  #1}
\newcommand{\legende}[1]{\par\vspace{6pt}{\popsemi\fontsize{9}{11.5}\selectfont\color{popmuted}#1}}

% ============================================================
% Signature OnBuch+
% ============================================================
\newcommand{\poplogoinline}[1]{{\poplogo\fontsize{#1}{#1}\selectfont\color{popink}OnBuch\color{poporange}+}}
\newcommand{\signatureonbuch}{%
  \begin{tcolorbox}[enhanced,colback=popink,colframe=popink,boxrule=2.2pt,arc=10pt,
      drop shadow=poporange,left=14pt,right=14pt,top=10pt,bottom=10pt,before skip=0pt,after skip=0pt]
    \tikz[baseline=-0.7ex]\node[circle,fill=popgreen,draw=popcream,line width=1.3pt,minimum size=22pt,inner sep=0]{\popcheck{white}};%
    \hspace{9pt}\begin{minipage}[c]{0.62\linewidth}
      {\popxbold\fontsize{12}{13}\selectfont\color{popcream}Signé\ \textbullet\ {\poplogo OnBuch\color{poporange}+}}\par\vspace{2pt}
      {\raggedright\popsemi\fontsize{8}{10.5}\selectfont\color{popline}\MakeUppercase{Équipe pédagogique OnBuch+}\\\MakeUppercase{Conforme au programme officiel MINESEC}\par}
    \end{minipage}\hfill
    {\poplogo\fontsize{22}{22}\selectfont\color{popcream}OnBuch\color{poporange}+}
  \end{tcolorbox}%
}

% ============================================================
% Page de garde
% #1 titre · #2 matière · #3 niveau · #4 module · #5 n° leçon · #6 durée ·
% #7 description / objectifs (texte ou itemize)
% ============================================================
\newcommand{\pagegarde}[7]{%
  \thispagestyle{empty}
  \newgeometry{a4paper,margin=1.7cm,top=1.6cm,bottom=1.4cm}%
  \begin{tikzpicture}[remember picture,overlay]
    \fill[poporange!18] ([xshift=-3.2cm,yshift=-3.6cm]current page.north east) circle (4.6cm);
    \fill[poppurple!14] ([xshift=2.4cm,yshift=7.6cm]current page.south west) circle (3.8cm);
  \end{tikzpicture}%
  {\poplogo\fontsize{30}{30}\selectfont\color{popink}OnBuch\color{poporange}+}\hfill
  \raisebox{0.6ex}{\poppill{Cours signé \textbullet\ Premium}{popink}{popcream}}\par
  \vspace{1pt}\popsquiggle{poppurple}\par
  \vspace{8pt}
  \begin{tcolorbox}[enhanced,colback=poporange,colframe=popink,boxrule=2.8pt,arc=13pt,
      drop shadow=popink,left=18pt,right=18pt,top=12pt,bottom=13pt,after skip=10pt]
    \poppill{#2 \textbullet\ #3}{popink}{popcream}\hspace{5pt}\poppill{#4}{white}{popink}\par\vspace{8pt}
    {\popdisplay\fontsize{14}{14}\selectfont\color{popink}LEÇON #5}\par\vspace{4pt}
    {\raggedright\hyphenpenalty=10000\exhyphenpenalty=10000\popdisplay\fontsize{25}{27}\selectfont\color{popcream}\MakeUppercase{#1}\par}
    \vspace{8pt}
    {\popxbold\fontsize{9.5}{9.5}\selectfont\color{popink}\MakeUppercase{Durée conseillée : #6}}
  \end{tcolorbox}
  \begin{tcolorbox}[enhanced,colback=white,colframe=popink,boxrule=2.2pt,arc=10pt,
      drop shadow=popink,left=16pt,right=16pt,top=10pt,bottom=10pt,after skip=8pt]
    \poppill{Dans cette leçon}{poppurple}{white}\par\vspace{5pt}
    \popsemi\fontsize{9.3}{12.6}\selectfont\color{popink}#7
  \end{tcolorbox}
  \noindent\begin{minipage}[t]{0.487\linewidth}
    \begin{tcolorbox}[enhanced,colback=popgreen,colframe=popink,boxrule=2.2pt,arc=10pt,
        drop shadow=popink,left=11pt,right=11pt,top=10pt,bottom=10pt,height=3.1cm,valign=center]
      \tcbox[colback=white,colframe=popink,boxrule=1.3pt,arc=5pt,boxsep=0pt,nobeforeafter,
        left=3pt,right=3pt,top=3pt,bottom=3pt,valign=center]{\qrcode[height=1.9cm]{https://play.google.com/store/apps/details?id=com.luvvix.onbuch&hl=fr}}%
      \hspace{8pt}\begin{minipage}[c]{0.5\linewidth}
        {\popdisplay\fontsize{11}{12.5}\selectfont\color{white}\MakeUppercase{Télécharge OnBuch}}\par\vspace{4pt}
        {\raggedright\popsemi\fontsize{8.4}{11}\selectfont\color{white}Gratuit \textbullet\ Google Play\\Tuteur IA, annales, concours, résultats\par}
      \end{minipage}
    \end{tcolorbox}
  \end{minipage}\hfill
  \begin{minipage}[t]{0.487\linewidth}
    \begin{tcolorbox}[enhanced,colback=poppurple,colframe=popink,boxrule=2.2pt,arc=10pt,
        drop shadow=popink,left=11pt,right=11pt,top=10pt,bottom=10pt,height=3.1cm,valign=center]
      \tikz[baseline=-0.7ex]\node[circle,fill=white,draw=popink,line width=1.3pt,minimum size=1.6cm,inner sep=0]{\popdisplay\fontsize{24}{24}\selectfont\color{poppurple}+};%
      \hspace{8pt}\begin{minipage}[c]{0.6\linewidth}
        {\popdisplay\fontsize{11}{12.5}\selectfont\color{white}\MakeUppercase{Passe à OnBuch+}}\par\vspace{4pt}
        {\raggedright\popsemi\fontsize{8.4}{11}\selectfont\color{white}Tous les cours signés, Tuteur IA illimité, fiches PDF — onglet OnBuch+ de l'app\par}
      \end{minipage}
    \end{tcolorbox}
  \end{minipage}
  \vspace{8pt}
  \popcontenu
  \vfill
  \signatureonbuch
  \restoregeometry
  \newpage
}

% Rangée « Ce cours contient » (page de garde)
\newcommand{\poptile}[3]{% #1 couleur #2 gros texte #3 légende
  \begin{tcolorbox}[enhanced,colback=#1,colframe=popink,boxrule=2pt,arc=9pt,shadow={2.5pt}{-2.5pt}{0pt}{popink},
      left=6pt,right=6pt,top=9pt,bottom=9pt,halign=center,valign=center,height=1.75cm,width=\linewidth,nobeforeafter]
    {\popdisplay\fontsize{12.5}{13}\selectfont\color{white}\MakeUppercase{#2}}\par\vspace{3pt}
    {\centering\popsemi\fontsize{7.8}{9.5}\selectfont\color{white}#3\par}
  \end{tcolorbox}}
\newcommand{\popcontenu}{%
  \par\noindent{\popxboldls\fontsize{7.5}{7.5}\selectfont\color{popmuted}CE COURS SIGNÉ CONTIENT}\par\vspace{7pt}
  \noindent\begin{minipage}[t]{0.235\linewidth}\poptile{popblue}{Cours}{complet, illustré, pas à pas}\end{minipage}\hfill
  \begin{minipage}[t]{0.235\linewidth}\poptile{poporange}{Méthodes}{et exercices résolus}\end{minipage}\hfill
  \begin{minipage}[t]{0.235\linewidth}\poptile{poppink}{Exercices}{type examen + corrigés}\end{minipage}\hfill
  \begin{minipage}[t]{0.235\linewidth}\poptile{popgreen}{Fiche bilan}{l'essentiel à retenir}\end{minipage}\par}

% Sommaire
\newtcolorbox{sommaire}{enhanced,colback=white,colframe=popink,boxrule=2.2pt,arc=10pt,
  drop shadow=popink,left=18pt,right=18pt,top=13pt,bottom=13pt,before skip=6pt,after skip=20pt,
  before upper={\poppill{Sommaire}{poppurple}{white}\par\vspace{9pt}\popsemi\fontsize{10.6}{14.5}\selectfont\color{popink}}}

% Signature de fin de cours — poussée tout en bas de la dernière page
\newcommand{\findecours}{%
  \par\vfill\nopagebreak
  \begin{center}\popsquiggle{poporange}\end{center}\vspace{4pt}
  \signatureonbuch
}
\AtBeginDocument{\def\llbracket{\mathopen{[\mkern-3.5mu[}}\def\rrbracket{\mathclose{]\mkern-3.5mu]}}} % intervalles d'entiers (glyphe ⟦ absent de la police)
% Glyphes absents de Fira Math : redéfinis à partir de glyphes disponibles
\AtBeginDocument{%
  \def\setminus{\mathbin{\backslash}}\let\smallsetminus\setminus
  \def\vdots{\mathord{\vbox{\baselineskip3.2pt\lineskiplimit0pt\kern4pt\hbox{.}\hbox{.}\hbox{.}}}}%
  \def\ddots{\mathinner{\mkern1mu\raise6.5pt\hbox{.}\mkern2mu\raise3.6pt\hbox{.}\mkern2mu\raise0.7pt\hbox{.}\mkern1mu}}%
  \def\triangle{\mathord{\scalebox{0.9}{$\symup{\Delta}$}}}%
  \def\nabla{\mathord{\raisebox{\depth}{\scalebox{1}[-1]{$\symup{\Delta}$}}}}%
  \def\checkmark{\popcheck{popdark}}%
  \def\star{\mathbin{\ast}}\def\bigstar{\mathord{\ast}}%
  \def\diamond{\mathbin{\scalebox{0.6}{$\Diamond$}}}%
  \def\bigcup{\mathop{\mathchoice{\raisebox{-0.3em}{\scalebox{1.7}{$\cup$}}}{\scalebox{1.25}{$\cup$}}{\cup}{\cup}}\displaylimits}%
  \def\bigcap{\mathop{\mathchoice{\raisebox{-0.3em}{\scalebox{1.7}{$\cap$}}}{\scalebox{1.25}{$\cap$}}{\cap}{\cap}}\displaylimits}%
  \def\lesseqqgtr{\mathrel{\lessgtr}}\def\bigoplus{\mathop{\oplus}\displaylimits}%
}
\providecommand{\ssi}{si et seulement si} % abréviation fréquente
\AtBeginDocument{\providecommand{\wideparen}{\overparen}} % arc de cercle

%% FIGFIX-BEGIN v4 — correctifs globaux des figures (docs/onbuch-generation/tools/figfix)
\usepackage{adjustbox}\usepackage{etoolbox}
% toute figure TikZ/pgfplots plus large que la ligne est réduite à la largeur disponible
\BeforeBeginEnvironment{tikzpicture}{\begin{adjustbox}{max width=\linewidth}}
\AfterEndEnvironment{tikzpicture}{\end{adjustbox}}
% légendes pgfplots lisibles par-dessus les courbes
\pgfplotsset{every axis/.append style={legend style={fill=white,fill opacity=0.92,text opacity=1,draw=black!15,inner sep=3pt}}}
% étiquettes d'arêtes (syntaxe edge["..."]) avec fond blanc
\tikzset{every edge quotes/.append style={fill=white,inner sep=1.2pt}}
%% FIGFIX-END
\begin{document}

\pagegarde{Lesson 8 — Grammar: Prepositions and prepositional phrases (place, time, manner and direction)}{Anglais}{1ère A}{Chapter 1 : Family and Social Life}{ENA08}{2 h}{Cette leçon de grammaire présente les prépositions et les syntagmes prépositionnels de lieu, de temps, de manière et de direction en anglais. À travers l'observation d'exemples authentiques, des tableaux contrastifs avec le français et des exercices progressifs, les élèves apprennent à choisir et à employer correctement ces prépositions dans des phrases et de courts textes.
\vspace{4pt}
\begin{multicols}{2}\small
\begin{itemize}
\item À la fin de cette leçon, je sais identifier une préposition et un syntagme prépositionnel dans une phrase anglaise.
\item À la fin de cette leçon, je sais employer correctement les prépositions de lieu (in, on, at, under, behind, between, next to...).
\item À la fin de cette leçon, je sais employer correctement les prépositions de temps (in, on, at, for, since, during, by, until...).
\item À la fin de cette leçon, je sais exprimer la direction et le mouvement (to, into, out of, towards, across, through, along...).
\item À la fin de cette leçon, je sais exprimer la manière et le moyen (by, with, in, like, without...).
\item À la fin de cette leçon, je sais construire un syntagme prépositionnel complet (préposition + complément) et le placer correctement dans la phrase.
\end{itemize}\end{multicols}}

\begin{sommaire}
\begin{multicols}{2}
\begin{enumerate}
\item Observation et découverte
\item Prépositions de lieu (place)
\item Prépositions de temps (time)
\item Direction, mouvement et manière
\item Comparaison français/anglais et pièges des francophones
\item Entraînement guidé et production
\item Activité d'intégration
\item Méthodes et exercices résolus
\item Exercices
\item Corrigés des exercices
\item Fiche bilan
\end{enumerate}
\end{multicols}
\end{sommaire}

\begin{objectifs}
\begin{itemize}
\item À la fin de cette leçon, je sais identifier une préposition et un syntagme prépositionnel dans une phrase anglaise.
\item À la fin de cette leçon, je sais employer correctement les prépositions de lieu (in, on, at, under, behind, between, next to...).
\item À la fin de cette leçon, je sais employer correctement les prépositions de temps (in, on, at, for, since, during, by, until...).
\item À la fin de cette leçon, je sais exprimer la direction et le mouvement (to, into, out of, towards, across, through, along...).
\item À la fin de cette leçon, je sais exprimer la manière et le moyen (by, with, in, like, without...).
\item À la fin de cette leçon, je sais construire un syntagme prépositionnel complet (préposition + complément) et le placer correctement dans la phrase.
\item À la fin de cette leçon, je sais éviter les erreurs fréquentes des francophones (arrive to, depend of, listen the radio...).
\item À la fin de cette leçon, je sais utiliser les prépositions dans un court texte descriptif ou narratif.
\end{itemize}
\end{objectifs}

\begin{prerequis}
\begin{itemize}
\item La structure de base de la phrase anglaise (sujet + verbe + complément) vue en Seconde
\item Les principaux auxiliaires et temps simples (present simple, past simple)
\item Le vocabulaire de base de la famille, de la maison et de la ville (Chapter 1 : Family and Social Life)
\item La notion de complément de phrase (complément de lieu, de temps) vue en classe de Seconde
\end{itemize}
\end{prerequis}

\newpage

\coursec{Observation et découverte}

\courssub{Situation d'accroche et problématique}

Imagine : ton ami te demande où se trouve la nouvelle librairie de Yaoundé. Tu réponds : \eng{It is in front of the post office, on the left of the bank.} « Elle est en face de la poste, à gauche de la banque. » Le lendemain, tu lui donnes rendez-vous : \eng{Let's meet at 4 p.m. on Saturday, in front of the school gate.} « Retrouvons-nous à 16 h samedi, devant le portail de l'école. » Sans le savoir, tu viens d'utiliser des \cle{prépositions} de lieu, de temps et de direction !

En anglais comme en français, ces petits mots (\eng{in}, \eng{on}, \eng{at}, \eng{by}, \eng{to}...) sont partout : dans les conversations, les indications routières à Bamenda ou à Lagos, et les épreuves du Bac. Mais attention : ils ne se traduisent presque jamais mot à mot. Cette leçon t'apprend à les maîtriser pour ne plus jamais dire \eng{I arrived to school} !

\textbf{Problématique :} comment reconnaître une préposition, comment construire un syntagme prépositionnel correct, et comment choisir la bonne préposition anglaise quand le français en utilise une autre ?

\courssub{Lecture et observation : un texte riche en prépositions}

Lis attentivement ce court texte. Un oncle camerounais décrit sa maison et son quartier. Pendant la lecture, \textbf{souligne tous les petits mots} qui semblent introduire un lieu, un moment ou une direction.

\begin{textebox}[My uncle's neighbourhood]
My uncle lives \textbf{in} Douala. His house is \textbf{on} a quiet street, \textbf{between} a church and a small shop. Every morning, \textbf{at} 6 o'clock, he walks \textbf{to} the market \textbf{with} his wife. They go \textbf{across} the bridge and \textbf{through} the crowd. \textbf{In} the evening, they sit \textbf{in front of} the house and talk \textbf{with} their neighbours.
\end{textebox}

\textbf{Glossaire :}

\begin{center}
\begin{tabularx}{\linewidth}{|l|l|X|}
\hline
\rowcolor{popblueL} \textbf{Mot anglais} & \textbf{Nature} & \textbf{Traduction} \\
\hline
uncle & nom & oncle \\
\hline
quiet & adjectif & calme, tranquille \\
\hline
church & nom & église \\
\hline
shop & nom & boutique, magasin \\
\hline
market & nom & marché \\
\hline
bridge & nom & pont \\
\hline
crowd & nom & foule \\
\hline
neighbours & nom & voisins \\
\hline
\end{tabularx}
\end{center}

\textbf{Questions de compréhension :}

\begin{enumerate}
\item Where does the uncle live? \emph{(Où habite l'oncle ?)}
\item What time does he walk to the market? \emph{(À quelle heure va-t-il au marché ?)}
\item What do they cross on their way? \emph{(Que traversent-ils en chemin ?)}
\item Where do they sit in the evening? \emph{(Où s'assoient-ils le soir ?)}
\end{enumerate}

\begin{experience}[Classement spontané]
En binôme, recopie les prépositions soulignées dans le texte (\eng{in, on, between, at, to, with, across, through, in front of}) et classe-les en quatre colonnes selon la question à laquelle elles répondent : \textbf{où ?} (place), \textbf{quand ?} (time), \textbf{comment / avec qui ?} (manner), \textbf{vers où ?} (direction). Compare ensuite ton classement avec celui de tes camarades : c'est exactement le plan de cette leçon !
\end{experience}

\courssub{Qu'est-ce qu'une préposition ?}

\begin{definition}[Preposition — la préposition]
Une \cle{préposition} (\eng{preposition}) est un \textbf{mot invariable} placé \textbf{devant} un nom, un pronom ou un groupe nominal pour introduire un complément. Elle exprime une relation : de lieu, de temps, de direction ou de manière.
\par
Exemples : \eng{in}, \eng{on}, \eng{at}, \eng{to}, \eng{with}, \eng{under}, \eng{between}, \eng{across}.
\end{definition}

\begin{propriete}[Trois propriétés essentielles]
\begin{enumerate}
\item \textbf{Invariable} : une préposition ne change jamais de forme (pas de pluriel, pas de conjugaison). On dit \eng{in the house} et \eng{in the houses} : \eng{in} reste identique.
\item \textbf{Toujours suivie d'un complément} : une préposition ne s'emploie jamais seule dans une phrase affirmative ; elle introduit un nom (\eng{to school}), un pronom (\eng{with me}) ou un groupe nominal (\eng{on the big table}).
\item \textbf{Pas de traduction mot à mot} : \eng{on} ne veut pas toujours dire « sur » (\eng{on Monday} = lundi, \eng{on foot} = à pied). C'est le piège numéro un des francophones.
\end{enumerate}
\end{propriete}

\begin{exemplebox}[Premiers exemples traduits]
\eng{The cat is under the table.} « Le chat est sous la table. »
\par
\eng{She arrived at night.} « Elle est arrivée la nuit. »
\par
\eng{The students are in the classroom.} « Les élèves sont dans la salle de classe. »
\par
\eng{He travels to Bamenda by bus.} « Il voyage à Bamenda en car. »
\end{exemplebox}

\courssub{Le syntagme prépositionnel (prepositional phrase)}

\begin{definition}[Prepositional phrase — le syntagme prépositionnel]
Un \cle{syntagme prépositionnel} (\eng{prepositional phrase}) est un groupe de mots formé d'une \textbf{préposition + son complément} (nom, pronom ou groupe nominal).
\par
Exemple : \eng{under the bed} = \eng{under} (préposition) + \eng{the bed} (complément) « sous le lit ».
\par
Autres exemples : \eng{in the garden} « dans le jardin », \eng{at 5 o'clock} « à 5 heures », \eng{with my sister} « avec ma sœur ».
\end{definition}

\begin{popfigure}
\begin{center}
\begin{tikzpicture}[
box/.style={draw=popblue, fill=popblueL, rounded corners=3pt, align=center, font=\small, minimum height=0.9cm},
lab/.style={font=\footnotesize\itshape, text=popmuted}
]
\node[box, fill=poporangeL, draw=poporange, font=\small\bfseries] (pp) at (0,0) {Prepositional phrase};
\node[box, align=center] (prep) at (-3.4,-1.8){Preposition\\ \eng{on}};
\node[box, fill=popgreenL, draw=popgreen, align=center] (comp) at (3.4,-1.8){Complement\\ \eng{the table}};
\node[font=\large\bfseries, text=popdark] at (0,-1.8) {=};
\node[font=\large\bfseries, text=popdark] at (1.7,-1.8) {+};
\draw[-{Stealth}, thick, poporange] (pp.south west) -- (prep.north);
\draw[-{Stealth}, thick, popgreen] (pp.south east) -- (comp.north);
\node[lab] at (-3.4,-2.7) {le petit mot invariable};
\node[lab] at (3.4,-2.7) {nom, pronom ou groupe nominal};
\node[font=\small\bfseries, text=popblue] at (0,-3.3) {\eng{on the table} = « sur la table »};
\end{tikzpicture}
\end{center}
\legende{Structure du syntagme prépositionnel : préposition + complément.}
\end{popfigure}

\begin{methode}[Repérer un syntagme prépositionnel en quatre étapes]
\begin{enumerate}
\item \textbf{Cherche le petit mot invariable} : \eng{in, on, at, to, from, with, under, between...}
\item \textbf{Vérifie qu'il est suivi d'un nom ou d'un pronom} : \eng{in} + \eng{the kitchen}, \eng{with} + \eng{her}. Si le mot est suivi d'un verbe conjugué, ce n'est pas une préposition.
\item \textbf{Pose la question} : \emph{où ?} (place), \emph{quand ?} (time), \emph{comment ? avec quoi ?} (manner), \emph{vers où ? d'où ?} (direction).
\item \textbf{Encadre tout le groupe} : la préposition \textbf{et} son complément forment un seul bloc. Dans \eng{The boy sits on the old wooden chair}, le syntagme prépositionnel est \eng{on the old wooden chair}, pas seulement \eng{on}.
\end{enumerate}
\end{methode}

\begin{exoresolu}[Repérage dans une phrase]
Énoncé : souligne tous les syntagmes prépositionnels dans la phrase suivante, puis indique la question à laquelle chacun répond.
\par
\eng{On Saturday morning, Amina goes to the stadium with her brother and they sit behind the goal.}
\tcblower
Solution détaillée :
\begin{itemize}
\item \eng{On Saturday morning} « samedi matin » : préposition \eng{on} + complément \eng{Saturday morning} ; répond à \emph{quand ?} (time).
\item \eng{to the stadium} « au stade » : préposition \eng{to} + complément \eng{the stadium} ; répond à \emph{vers où ?} (direction).
\item \eng{with her brother} « avec son frère » : préposition \eng{with} + complément \eng{her brother} ; répond à \emph{avec qui ?} (manner/accompagnement).
\item \eng{behind the goal} « derrière le but » : préposition \eng{behind} + complément \eng{the goal} ; répond à \emph{où ?} (place).
\end{itemize}
Vérification : chaque préposition est bien suivie d'un groupe nominal, et aucune n'est suivie d'un verbe conjugué.
\end{exoresolu}

\courssub{Les fonctions du syntagme prépositionnel}

Un syntagme prépositionnel peut jouer trois rôles principaux dans la phrase.

\begin{propriete}[Les trois fonctions]
\begin{enumerate}
\item \textbf{Complément du verbe} : il précise où, quand ou comment se déroule l'action. \eng{She lives in Douala.} « Elle habite à Douala. » (\eng{in Douala} complète le verbe \eng{lives}.)
\item \textbf{Complément du nom} : il précise un nom. \eng{The house on the hill is beautiful.} « La maison sur la colline est belle. » (\eng{on the hill} précise quel \eng{house}.)
\item \textbf{Complément de phrase} : il situe toute la phrase dans le temps ou l'espace ; souvent placé en tête et suivi d'une virgule. \eng{In the evening, they talk with their neighbours.} « Le soir, ils parlent avec leurs voisins. »
\end{enumerate}
\end{propriete}

\begin{attention}[Piège : préposition ou adverbe ?]
Certains mots (\eng{in}, \eng{up}, \eng{down}) peuvent être des adverbes quand ils sont \textbf{seuls}, sans complément. Compare : \eng{Come in!} « Entre ! » (\eng{in} est un adverbe, rien ne suit) et \eng{She is in the kitchen.} « Elle est dans la cuisine. » (\eng{in} est une préposition, suivie de \eng{the kitchen}). Le test est simple : \textbf{pas de complément, pas de préposition}.
\end{attention}

\courssub{Les quatre grandes catégories de la leçon}

Dans cette leçon, nous classerons les prépositions en quatre familles, selon la question à laquelle le syntagme répond.

\begin{center}
\begin{tabularx}{\linewidth}{|l|X|X|}
\hline
\rowcolor{popblueL} \textbf{Catégorie} & \textbf{Question} & \textbf{Exemples} \\
\hline
Place (lieu) & where? (où ?) & \eng{in the garden}, \eng{on the table}, \eng{between the houses} \\
\hline
Time (temps) & when? (quand ?) & \eng{at 6 o'clock}, \eng{on Monday}, \eng{in July} \\
\hline
Manner (manière) & how? with what? (comment ?) & \eng{by bus}, \eng{with a pen}, \eng{with pleasure} \\
\hline
Direction & where to? (vers où ?) & \eng{to the market}, \eng{across the bridge}, \eng{through the crowd} \\
\hline
\end{tabularx}
\end{center}

\begin{popfigure}
\begin{center}
\begin{tikzpicture}[
mindmap,
grow cyclic,
every node/.style={concept, align=center, font=\small},
concept color=popblue!60,
level 1/.append style={level distance=3.6cm, sibling angle=90},
level 2/.append style={level distance=2.4cm, sibling angle=45, font=\footnotesize}
]
\node[concept color=poporange!70, font=\small\bfseries] {PREPOSITIONS}
  child[concept color=popblue!50] { node[align=center]{PLACE\\ \emph{where?}}
    child { node[align=center]{\eng{in} \\ the house} }
    child { node[align=center]{\eng{on} \\ the table} }
    child { node[align=center]{\eng{under} \\ the bed} }
  }
  child[concept color=popgreen!50] { node[align=center]{TIME\\ \emph{when?}}
    child { node[align=center]{\eng{at} \\ 5 o'clock} }
    child { node[align=center]{\eng{on} \\ Monday} }
    child { node[align=center]{\eng{in} \\ 2024} }
  }
  child[concept color=poppurple!45] { node[align=center]{MANNER\\ \emph{how?}}
    child { node[align=center]{\eng{by} \\ bus} }
    child { node[align=center]{\eng{with} \\ a pen} }
    child { node[align=center]{\eng{with} \\ care} }
  }
  child[concept color=poppink!50] { node[align=center]{DIRECTION\\ \emph{where to?}}
    child { node[align=center]{\eng{to} \\ school} }
    child { node[align=center]{\eng{across} \\ the road} }
    child { node[align=center]{\eng{through} \\ the crowd} }
  };
\end{tikzpicture}
\end{center}
\legende{Carte mentale : les quatre familles de prépositions de la leçon.}
\end{popfigure}

\begin{savaistu}[Pourquoi « preposition » ?]
Le mot anglais \eng{preposition} vient du latin \emph{praeponere}, « placer devant » : la préposition est littéralement le mot « pré-posé », placé \textbf{avant} son complément. En anglais parlé familier, on entend parfois une préposition en fin de phrase (\eng{Who are you talking to?} « À qui parles-tu ? »), mais à l'écrit scolaire et au Bac, garde la structure classique : préposition + complément.
\end{savaistu}

\begin{exercice}[À toi de jouer]
Dans le texte « My uncle's neighbourhood », retrouve : 1) deux syntagmes prépositionnels de lieu ; 2) un syntagme de temps ; 3) un syntagme de direction ; 4) un syntagme de manière ou d'accompagnement. Recopie-les et traduis-les en français.
\end{exercice}

\begin{corrige}[Exercice « À toi de jouer »]
1) Lieu : \eng{on a quiet street} « dans une rue calme », \eng{in front of the house} « devant la maison » (aussi possible : \eng{in Douala}, \eng{between a church and a small shop}). 2) Temps : \eng{at 6 o'clock} « à 6 heures » (ou \eng{in the evening} « le soir »). 3) Direction : \eng{to the market} « au marché » (ou \eng{across the bridge}, \eng{through the crowd}). 4) Manière/accompagnement : \eng{with his wife} « avec sa femme » (ou \eng{with their neighbours}).
\end{corrige}

\begin{aretenir}[L'essentiel de la section]
\begin{itemize}
\item Une \cle{préposition} est un mot invariable placé devant un nom, un pronom ou un groupe nominal.
\item Un \cle{syntagme prépositionnel} = préposition + complément ; il répond aux questions \emph{où ? quand ? comment ? vers où ?}
\item Il peut être complément du verbe, complément du nom ou complément de phrase.
\item Les quatre familles de la leçon : \textbf{place}, \textbf{time}, \textbf{manner}, \textbf{direction}.
\item Les prépositions ne se traduisent pas mot à mot : apprends chaque emploi en contexte.
\end{itemize}
\end{aretenir}

\coursec{Lesson 8 — Prepositions and Prepositional Phrases : Place, Time, Manner and Direction}

\courssub{Observation et découverte : un week-end à Buea}

\begin{textebox}[Blog post — “A Weekend at the Foot of Mount Cameroon”]
Last Saturday, I travelled \textbf{to} Buea \textbf{by} bus. The journey took three hours. We arrived \textbf{at} the main park \textbf{at} noon. My cousin was waiting \textbf{for} me \textbf{at} the entrance, \textbf{next to} a bendskin stand. We walked \textbf{to} his house, which is \textbf{in} Molyko, \textbf{near} the university. It is a small house \textbf{with} a red roof, \textbf{behind} the big supermarket. \textbf{In} the evening, we went \textbf{up} the mountain \textbf{on} foot. The view \textbf{from} the top was amazing. We came \textbf{down} \textbf{by} car \textbf{with} a driver we met \textbf{at} the hotel. \textbf{On} Sunday, \textbf{at} 8 a.m., we took a taxi \textbf{to} the airport. I was \textbf{in} a hurry but the traffic moved \textbf{slowly}. I finally got \textbf{on} the plane \textbf{just in time}.
\end{textebox}

\begin{definition}[Vocabulaire utile — Mots et expressions du texte]
\begin{itemize}
\item \eng{travel} (v.) — voyager \quad \eng{journey} (n.) — trajet, voyage
\item \eng{arrive at / in} (v.) — arriver à (point / grande ville) \quad \eng{wait for} — attendre
\item \eng{bendskin} (n., Cam.) — moto-taxi \quad \eng{stand} (n.) — station, arrêt
\item \eng{Molyko} (n.) — quartier de Buea \quad \eng{roof} (n.) — toit
\item \eng{view} (n.) — vue, panorama \quad \eng{amazing} (adj.) — incroyable
\item \eng{in a hurry} (expr.) — pressé \quad \eng{just in time} (expr.) — juste à temps
\item \eng{on foot} (expr.) — à pied \quad \eng{by car / by bus / by plane} — en voiture / en bus / en avion
\end{itemize}
\end{definition}

\begin{experience}[Activité de découverte — Listening / Reading]
\textbf{Étape 1.} Écoute (ou lis) le texte ci-dessus. Souligne toutes les prépositions (\eng{to, by, at, for, next to, in, near, with, behind, on, up, from, down, on, at, in, on, just in time}). \\
\textbf{Étape 2.} Classe-les en quatre catégories dans ton cahier : \textbf{Lieu (Place)}, \textbf{Temps (Time)}, \textbf{Direction / Mouvement (Direction)}, \textbf{Manière (Manner)}. \\
\textbf{Étape 3.} Compare ton classement avec ton voisin. Quelles prépositions reviennent dans plusieurs catégories ? (Indice : \eng{in}, \eng{on}, \eng{at}).
\end{experience}

\courssub{Prépositions de lieu (Place)}

Dans la section précédente, nous avons observé le texte de découverte et relevé des expressions comme \eng{in Molyko}, \eng{on foot}, \eng{at the entrance}, \eng{next to a bendskin stand} ou \eng{behind the supermarket}. Toutes contiennent un petit mot placé devant un nom (ou un groupe nominal) : une \cle{préposition}. Ces petits mots sont redoutables pour les francophones, car le français utilise souvent « à », « dans » ou « sur » là où l'anglais fait un choix précis entre \eng{in}, \eng{on} et \eng{at}. Cette section te donne les clés pour ne plus jamais hésiter.

\courssub{La triade fondamentale : in / on / at}

\textbf{Observation}\par

Lis attentivement ces trois phrases tirées de la vie quotidienne au Cameroun :

\begin{exemplebox}[Trois phrases à observer]
\eng{The students are in the classroom.} « Les élèves sont dans la salle de classe. »

\eng{The books are on the desk.} « Les livres sont sur le bureau. »

\eng{Wait for me at the entrance.} « Attends-moi à l'entrée. »
\end{exemplebox}

Dans la première phrase, les élèves sont \emph{à l'intérieur} d'un espace fermé : \eng{in}. Dans la deuxième, les livres reposent \emph{sur une surface} : \eng{on}. Dans la troisième, l'entrée est envisagée comme \emph{un point précis} sur un parcours : \eng{at}. Toute la logique des prépositions de lieu anglaises tient dans cette image mentale : \textbf{un volume} (\eng{in}), \textbf{une surface} (\eng{on}), \textbf{un point} (\eng{at}).

\begin{propriete}[La règle de la triade in / on / at pour le lieu]
\begin{itemize}
\item \cle{in} = à l'intérieur d'un espace, clos ou vaste (volume) : \eng{in a room} (dans une pièce), \eng{in a box} (dans une boîte), \eng{in Limbe} (à Limbé), \eng{in Cameroon} (au Cameroun), \eng{in the world} (dans le monde), \eng{in the street} (GB) / \eng{on the street} (US) (dans la rue).
\item \cle{on} = sur une surface, une ligne ou une page/écran (surface) : \eng{on the table} (sur la table), \eng{on the wall} (au mur), \eng{on the floor} (par terre), \eng{on the road} (sur la route), \eng{on page 5} (à la page 5), \eng{on the map} (sur la carte), \eng{on the screen} (à l'écran).
\item \cle{at} = en un point précis, un repère, une adresse précise ou un événement (point) : \eng{at the door} (à la porte), \eng{at the bus stop} (à l'arrêt de bus), \eng{at the crossroads} (au carrefour), \eng{at the top} (en haut), \eng{at the back} (au fond), \eng{at 10 Main Street} (au 10 rue Main), \eng{at the party} (à la fête).
\end{itemize}
\textbf{Astuce de mémorisation :} plus l'espace est grand et englobant, plus on va vers \eng{in} ; plus il est réduit à un repère, plus on va vers \eng{at}. La surface est le cas intermédiaire : \eng{on}.
\end{propriete}

\begin{popfigure}
\begin{tikzpicture}[scale=1.0, every node/.style={font=\small}]
% IN : une boîte avec un point dedans
\draw[thick, popblue, fill=popblue!15] (0,0) rectangle (2.4,1.8);
\fill[popink] (1.2,0.9) circle (0.09);
\node[popblue, align=center] at (1.2,-0.55) {\textbf{in}\\ volume / espace 3D};
% ON : une ligne avec un point dessus
\draw[very thick, poporange] (4.0,0.4) -- (6.8,0.4);
\fill[popink] (5.4,0.4) circle (0.09);
\node[poporange, align=center] at (5.4,-0.55) {\textbf{on}\\ surface / ligne 2D};
% AT : un point avec une flèche qui arrive dessus
\draw[-{Stealth[length=3mm]}, very thick, popgreen] (8.2,1.6) -- (9.3,0.7);
\fill[popink] (9.5,0.55) circle (0.09);
\node[popgreen, align=center] at (9.3,-0.55) {\textbf{at}\\ point précis 0D};
\end{tikzpicture}
\legende{La logique spatiale de la triade : in = volume 3D, on = surface 2D, at = point 0D.}
\end{popfigure}

\textbf{Le choix dépend du point de vue du locuteur}\par

Un même lieu peut prendre \eng{in} ou \eng{at} selon la façon dont on l'envisage. Compare :

\eng{My father works at the hospital.} « Mon père travaille à l'hôpital. » (l'hôpital = simple repère, le lieu de son travail)

\eng{My mother is in hospital; she had an operation yesterday.} « Ma mère est à l'hôpital ; elle a été opérée hier. » (elle est \emph{à l'intérieur}, en tant que patiente — \eng{in hospital} sans article en GB)

De même : \eng{Let's meet at the stadium.} « Retrouvons-nous au stade » (point de rendez-vous), mais \eng{There were 20,000 fans in the stadium.} « Il y avait 20 000 supporters dans le stade » (à l'intérieur de l'enceinte).

\eng{She is at school.} (élève/enseignant, activité normale) vs \eng{She is in the school.} (visiteur, à l'intérieur du bâtiment).

\begin{exemplebox}[Exemples traduits — la triade en contexte camerounais]
\eng{My grandmother lives in Limbe.} « Ma grand-mère vit à Limbé. » (ville = espace vaste)

\eng{There is a family photo on the wall.} « Il y a une photo de famille au mur. » (surface verticale)

\eng{The bendskin driver is waiting at the junction.} « Le moto-taximan attend au carrefour. » (point de repère)

\eng{The cocoa beans are drying in the sun, on large mats.} « Les fèves de cacao sèchent au soleil, sur de grandes nattes. » (\eng{in the sun} = exposition ; \eng{on mats} = surface de contact)

\eng{She is sitting at her desk, in the last row.} « Elle est assise à sa table, au dernier rang. » (\eng{at desk} = position de travail ; \eng{in row} = à l'intérieur de l'alignement)
\end{exemplebox}

\begin{attention}[Ne traduis pas « à » systématiquement par « to » !]
C'est l'erreur numéro un des francophones. En français, « je suis \textbf{à} l'école » et « je vais \textbf{à} l'école » utilisent la même préposition. En anglais, non :
\begin{itemize}
\item Position (où ?) : \eng{I am \textbf{at} school.} « Je suis à l'école. »
\item Mouvement (vers où ?) : \eng{I go \textbf{to} school.} « Je vais à l'école. »
\end{itemize}
\eng{to} exprime le \textbf{mouvement vers}, jamais la position statique. On ne dit donc jamais \eng{I am to school}. Retiens : \textbf{position = at / in / on ; direction = to} (nous détaillerons la direction en section 4).
\end{attention}

\courssub{Les prépositions de position : situer les choses les unes par rapport aux autres}

\textbf{Observation}\par

Imagine la cour de récréation de ton lycée. Pour décrire où se trouvent les élèves, tu as besoin de mots plus précis que \eng{in}, \eng{on} ou \eng{at} :

\eng{The flag is above the notice board. The headmaster's office is opposite the library. Aminatu is standing between Moussa and Léa, next to the big mango tree.}

Ces prépositions de \cle{position} (souvent des \cle{locutions prépositionnelles} de deux ou trois mots) se regroupent naturellement par paires d'opposés.

\begin{center}
\begin{tabularx}{\linewidth}{|X|X|X|}
\hline
\rowcolor{popblueL} English & Français & Exemple \\
\hline
above & au-dessus de (niveau supérieur, sans contact) & \eng{The clock is above the blackboard.} (L'horloge est au-dessus du tableau.) \\
\hline
below & en dessous de (niveau inférieur, sans contact) & \eng{The valley is below the village.} (La vallée est en contrebas du village.) \\
\hline
over & au-dessus de (recouvrant, franchissant, verticalité) & \eng{She put a cloth over the table.} (Elle a mis un tissu \emph{sur} la table = le couvre.) \\
\hline
under & sous (recouvert, verticalité directe) & \eng{The cat is hiding under the bed.} (Le chat se cache sous le lit.) \\
\hline
behind & derrière & \eng{The kitchen is behind the house.} (La cuisine est derrière la maison.) \\
\hline
in front of & devant (côté avant) & \eng{A taxi is parked in front of the school.} (Un taxi est garé devant l'école.) \\
\hline
between & entre (deux éléments distincts) & \eng{Limbe lies between Douala and Buea.} (Limbé se trouve entre Douala et Buea.) \\
\hline
among & parmi (groupe, masse, trois+) & \eng{She felt happy among her cousins.} (Elle se sentait heureuse parmi ses cousins.) \\
\hline
next to / beside & à côté de (contact latéral proche) & \eng{I sit next to my best friend.} (Je suis assis à côté de mon meilleur ami.) \\
\hline
near & près de (proximité, pas forcément contact) & \eng{Our house is near the market.} (Notre maison est près du marché.) \\
\hline
opposite & en face de (de l'autre côté d'un espace : rue, place, couloir) & \eng{The pharmacy is opposite the church.} (La pharmacie est en face de l'église.) \\
\hline
\end{tabularx}
\end{center}

\begin{popfigure}
\begin{tikzpicture}[scale=0.9, every node/.style={font=\small}]
% Salle de classe
\draw[thick, popdark] (0,0) rectangle (10,6);
% Tableau
\draw[thick, popdark, fill=popgold!20] (3.5,5.4) rectangle (6.5,5.9);
\node at (5,5.65) {Blackboard};
% Bureau prof
\draw[thick, popblue, fill=popblue!15] (3.8,4.3) rectangle (6.2,5.0);
\node[align=center] at (5,4.65) {Teacher's\\ desk};
% Élèves rang 1
\draw[thick, popgreen, fill=popgreen!15] (1.2,2.6) rectangle (3.0,3.4);
\node[align=center] at (2.1,3.0) {Amina};
\draw[thick, popgreen, fill=popgreen!15] (4.0,2.6) rectangle (5.8,3.4);
\node[align=center] at (4.9,3.0) {Bruno};
\draw[thick, popgreen, fill=popgreen!15] (6.8,2.6) rectangle (8.6,3.4);
\node[align=center] at (7.7,3.0) {Carine};
% Élève rang 2 (devant)
\draw[thick, poporange, fill=poporange!15] (4.0,0.6) rectangle (5.8,1.4);
\node[align=center] at (4.9,1.0) {David};
% Placard
\draw[thick, poppurple, fill=poppurple!15] (0.3,4.3) rectangle (1.9,5.0);
\node[align=center] at (1.1,4.65) {Cupboard};
% Flèches et labels
\node[popink, align=center] at (9.0,4.65) {\textbf{opposite}\\ the cupboard};
\draw[-{Stealth[length=2.5mm]}, popink, thick] (8.7,4.65) -- (2.1,4.65);
\node[popink, align=center] at (4.9,2.15) {\textbf{behind} Bruno};
\draw[-{Stealth[length=2.5mm]}, popink, thick] (4.9,2.3) -- (4.9,2.55);
\node[popink, align=center] at (4.9,3.85) {\textbf{in front of} David};
\node[popink, align=center] at (7.7,3.85) {\textbf{next to} Bruno};
\node[popink, align=center] at (2.1,3.85) {\textbf{between} cupboard \& Bruno};
\node[popink, align=center] at (5,5.2) {\textbf{above} teacher's desk};
\draw[-{Stealth[length=2mm]}, popink, thick] (5,5.1) -- (5,5.05);
\end{tikzpicture}
\legende{Plan de salle : Bruno est \eng{between} Amina and Carine ; David est \eng{behind} Bruno ; le tableau est \eng{above} le bureau ; le placard est \eng{opposite} Carine.}
\end{popfigure}

\courssub{Between ou among ? Above ou over ? Les distinctions qui font la différence}

\begin{propriete}[Between = deux éléments distincts ; among = au milieu d'un groupe]
\begin{itemize}
\item \cle{between} s'emploie pour situer quelque chose \textbf{entre deux} éléments (ou entre des éléments bien identifiés, pris un à un) : \eng{The post office is between the bank and the supermarket.} « La poste est entre la banque et le supermarché. » \eng{An agreement between the three countries.} (Accord entre les trois pays — considérés individuellement).
\item \cle{among} (ou \eng{amongst}, plus littéraire) s'emploie pour situer quelque chose \textbf{au milieu d'un groupe} (trois éléments ou plus, vus comme une masse indistincte) : \eng{I found my little brother among the crowd.} « J'ai retrouvé mon petit frère dans la foule. » \eng{Divide the mangoes among the children.} « Partage les mangues entre les enfants. » (répartition dans un groupe).
\end{itemize}
\textbf{Règle pratique au lycée :} \textbf{deux = between ; plusieurs (groupe) = among}.
\end{propriete}

\begin{propriete}[Above / over et below / under : la nuance du recouvrement et de la verticalité]
\begin{itemize}
\item \cle{above} et \cle{below} indiquent simplement un niveau supérieur ou inférieur, \textbf{sans contact ni recouvrement nécessaire} : \eng{The temperature is above 30 degrees.} « La température dépasse 30 degrés. » \eng{The picture hangs above the sofa.} « Le tableau est accroché au-dessus du canapé. » (espace entre les deux).
\item \cle{over} et \cle{under} impliquent souvent une \textbf{verticalité directe}, un contact, un recouvrement ou un franchissement : \eng{She spread a tablecloth over the table.} « Elle a étalé une nappe \emph{sur} la table » (la nappe couvre la table). \eng{There is a bridge over the river Wouri.} « Il y a un pont \emph{au-dessus/qui franchit} le Wouri. » \eng{The baby is sleeping under a mosquito net.} « Le bébé dort sous une moustiquaire » (recouvert par elle). \eng{He is under 18.} « Il a moins de 18 ans » (seuil).
\end{itemize}
\textbf{Cas où les deux sont possibles :} \eng{He held an umbrella over / above her head.} Mais pour le recouvrement et le franchissement, \eng{over} est obligatoire : \eng{The plane flew over Mount Cameroon.} (L'avion a survolé le mont Cameroun — il l'a franchi verticalement).
\end{propriete}

\begin{exemplebox}[Exemples traduits — les distinctions en action]
\eng{There is a crack between the two walls.} « Il y a une fissure entre les deux murs. »

\eng{She is very popular among her classmates.} « Elle est très appréciée parmi ses camarades. »

\eng{The sun is directly over our heads at noon.} « Le soleil est juste au-dessus de nos têtes à midi. » (verticalité)

\eng{The village below the hill is my grandfather's birthplace.} « Le village en contrebas de la colline est le village natal de mon grand-père. »
\end{exemplebox}

\courssub{Emplois figés et exceptions à apprendre par cœur}

Certaines combinaisons ne suivent pas la logique générale : ce sont des \cle{expressions figées} (collocations) qu'il faut mémoriser comme du vocabulaire. L'article (\eng{a/an/the}) y est souvent absent.

\begin{center}
\begin{tabularx}{\linewidth}{|X|X|}
\hline
\rowcolor{popblueL} Expression figée (GB) & Traduction / Contexte \\
\hline
\eng{at home} & à la maison (position) \\
\hline
\eng{at school / at college / at university} & à l'école / au lycée / à la fac (en tant qu'élève/étudiant) \\
\hline
\eng{at work} & au travail \\
\hline
\eng{at the top / at the bottom} & en haut / en bas (point extrême) \\
\hline
\eng{at the beginning / at the end} & au début / à la fin (point dans le temps ou l'ordre) \\
\hline
\eng{in bed} & au lit (état) \\
\hline
\eng{in hospital} (GB) / \eng{in the hospital} (US) & à l'hôpital (comme patient) \\
\hline
\eng{in prison} & en prison (détenu) \\
\hline
\eng{in the middle of} & au milieu de \\
\hline
\eng{on the left / on the right} & à gauche / à droite \\
\hline
\eng{on the first floor} & au premier étage (UK : 1er au-dessus du rez-de-chaussée) \\
\hline
\eng{on the ground floor} & au rez-de-chaussée \\
\hline
\eng{on the way} & en route \\
\hline
\eng{at the corner} (point) vs \eng{on the corner} (bâtiment occupant l'angle) & au coin de la rue \\
\hline
\eng{in the corner} (à l'intérieur d'une pièce) & dans le coin \\
\hline
\end{tabularx}
\end{center}

\eng{Turn left at the crossroads; the school is on the right, in the middle of the street.} « Tourne à gauche au carrefour ; l'école est à droite, au milieu de la rue. »

\begin{attention}[Transports : in the car mais on the bus !]
Les francophones disent « dans le bus », donc écrivent spontanément \eng{in the bus} : c'est une erreur lorsqu'on parle d'un trajet ou de la position à l'intérieur en tant que passager. La règle anglaise est la suivante :
\begin{itemize}
\item Véhicules où l'on \textbf{monte, peut se tenir debout et circuler} (grands véhicules publics) : \cle{on} — \eng{on the bus}, \eng{on the train}, \eng{on the plane}, \eng{on a boat/ship}, \eng{on a motorbike}, \eng{on a bendskin}, \eng{on a bicycle}.
\item Véhicules petits, privés, où l'on \textbf{s'assoit bas et ne circule pas} : \cle{in} — \eng{in the car}, \eng{in a taxi}, \eng{in a canoe}, \eng{in a helicopter}.
\end{itemize}
\eng{I met an old friend on the bus to Bamenda.} « J'ai rencontré un vieil ami dans le bus pour Bamenda. »
\eng{We were five in the taxi.} « Nous étions cinq dans le taxi. »
\end{attention}

\begin{savaistu}[In the street ou on the street ? Différences UK / US]
Les Britanniques disent traditionnellement \eng{in the street} (« dans la rue », espace clos entre les bâtiments), tandis que les Américains préfèrent \eng{on the street} (« sur la rue », surface de roulement). Les deux formes sont correctes à l'écrit : \eng{The children are playing in the street} (GB) / \eng{on the street} (US). Au baccalauréat, les deux sont acceptés — mais reste \textbf{cohérent} dans ta copie : ne passe pas de l'une à l'autre dans la même rédaction. Autre différence célèbre : les Britanniques vivent \eng{in Oxford Street}, les Américains \eng{on Fifth Avenue}. Pour les adresses précises : \eng{at 10 Downing Street} (UK/US).
\end{savaistu}

\courssub{Entraînement guidé — Lieu}

\begin{exoresolu}[Complète avec la préposition de lieu correcte]
Énoncé — Complète chaque phrase avec \eng{in}, \eng{on}, \eng{at}, \eng{between}, \eng{among}, \eng{next to}, \eng{opposite}, \eng{above}, \eng{under}, \eng{over}.

1. My uncle's shop is \dots\ the market, \dots\ the bakery and the pharmacy.

2. There is a beautiful map of Cameroon \dots\ the wall of our classroom.

3. The meeting starts soon; everybody is waiting \dots\ the headmaster's office.

4. I found a 500-franc coin \dots\ the passengers \dots\ the bus to Yaoundé.

5. The cat is sleeping \dots\ the table, just \dots\ the chair.

6. The plane flew \dots\ Mount Cameroon yesterday.

7. She divided the sweets \dots\ the four children.

\tcblower

Solution détaillée :

1. \eng{My uncle's shop is \textbf{opposite} the market, \textbf{between} the bakery and the pharmacy.} — « en face du marché » = \eng{opposite} (de l'autre côté de la rue) ; « entre » deux éléments précis = \eng{between}.

2. \eng{There is a beautiful map of Cameroon \textbf{on} the wall of our classroom.} — une carte est fixée \emph{sur la surface} du mur : \eng{on}.

3. \eng{Everybody is waiting \textbf{at} the headmaster's office.} — le bureau du proviseur est envisagé comme un \emph{point de rendez-vous} précis : \eng{at}.

4. \eng{I found a 500-franc coin \textbf{among} the passengers \textbf{on} the bus to Yaoundé.} — « parmi » un groupe (plusieurs personnes) = \eng{among} ; un trajet en bus = \eng{on the bus} (jamais \eng{in the bus} pour un passager).

5. \eng{The cat is sleeping \textbf{under} the table, just \textbf{next to} the chair.} — « sous la table » (recouvert) = \eng{under} ; « juste à côté de la chaise » = \eng{next to}.

6. \eng{The plane flew \textbf{over} Mount Cameroon yesterday.} — franchissement vertical, survol = \eng{over} (pas \eng{above}).

7. \eng{She divided the sweets \textbf{among} the four children.} — répartition dans un groupe de plus de deux = \eng{among}.
\end{exoresolu}

\begin{exercice}[À toi de jouer — Place]
1. Traduis en anglais : « Ma tante habite à Limbé, près de la plage. Sa maison se trouve entre l'église et l'école, en face d'un grand manguier. »

2. Corrige les erreurs dans ces phrases d'élèves : \\
a) \eng{I am to school now.} \\
b) \eng{She travels in the bus every morning.} \\
c) \eng{The picture is in the wall.} \\
d) \eng{He sat between the crowd.} \\
e) \eng{The bird flew above the river.} (contexte : il le traverse / le survole de près)

3. Décris ta salle de classe en cinq phrases en utilisant : \eng{in front of}, \eng{behind}, \eng{between}, \eng{above}, \eng{next to}.
\end{exercice}

\begin{corrige}[Exercice « À toi de jouer » — Place]
1. \eng{My aunt lives in Limbe, near the beach. Her house is between the church and the school, opposite a big mango tree.}

2. a) \eng{I am \textbf{at} school now.} (position = \eng{at}, pas \eng{to}) \\
b) \eng{She travels \textbf{on} the bus every morning.} (trajet en bus = \eng{on}) \\
c) \eng{The picture is \textbf{on} the wall.} (sur la surface du mur = \eng{on}) \\
d) \eng{He sat \textbf{among} the crowd.} (un groupe de plusieurs personnes = \eng{among}) \\
e) \eng{The bird flew \textbf{over} the river.} (franchissement / verticalité proche = \eng{over})

3. Modèle de réponse : \eng{The blackboard is in front of the students. The teacher's desk is between the blackboard and the first row. A clock hangs above the blackboard. My friend sits next to me, and the cupboard is behind us.}
\end{corrige}

\begin{aretenir}[L'essentiel — Prépositions de lieu]
\begin{itemize}
\item Triade spatiale : \cle{in} = volume (dans), \cle{on} = surface (sur), \cle{at} = point (à).
\item \eng{to} = mouvement vers (direction), JAMAIS position statique.
\item \cle{between} = deux éléments distincts ; \cle{among} = groupe / masse.
\item \cle{over}/\cle{under} = verticalité, recouvrement, franchissement ; \cle{above}/\cle{below} = simple niveau.
\item \cle{opposite} = en face de (séparé par un espace) ; \cle{in front of} = devant (même côté).
\item Expressions figées sans article : \eng{at home}, \eng{at school}, \eng{at work}, \eng{in bed}, \eng{in hospital}, \eng{in prison}, \eng{in the middle of}, \eng{on the left/right}, \eng{on the floor}.
\item Transports : \eng{in the car / in a taxi} MAIS \eng{on the bus / on the train / on the plane / on a motorbike / on foot}.
\end{itemize}
\end{aretenir}

\courssub{Prépositions de temps (Time)}

La bonne nouvelle : la triade \eng{in} / \eng{on} / \cle{at} fonctionne \textbf{exactement avec la même logique spatiale} pour le temps ! Visualise le temps comme une ligne.

\begin{propriete}[La triade in / on / at pour le temps]
\begin{itemize}
\item \cle{in} = \textbf{grande durée} (volume temporel) : parties de la journée (\eng{in the morning}, \eng{in the afternoon}, \eng{in the evening} — mais \eng{at night}), mois (\eng{in January}), saisons (\eng{in summer}), années (\eng{in 2024}), siècles (\eng{in the 21st century}), durées (\eng{in two hours}, \eng{in a few minutes}).
\item \cle{on} = \textbf{unité de temps délimitée} (surface : un jour, une date) : jours de la semaine (\eng{on Monday}, \eng{on Fridays}), dates précises (\eng{on 20 May}, \eng{on Christmas Day}), parties de journée précisées (\eng{on Monday morning}, \eng{on the evening of the 14th}), jours fériés (\eng{on Easter Monday}).
\item \cle{at} = \textbf{point précis sur l'horloge} (point temporel) : heures exactes (\eng{at 7:30}, \eng{at noon}, \eng{at midnight}, \eng{at dawn}, \eng{at dusk}), moments précis (\eng{at the moment}, \eng{at present}, \eng{at the same time}, \eng{at night} = la nuit comme période de sommeil), âges (\eng{at the age of 18}).
\end{itemize}
\end{propriete}

\begin{center}
\begin{tabularx}{\linewidth}{|X|X|X|}
\hline
\rowcolor{popblueL} \textbf{in} (Période large) & \textbf{on} (Jour / Date) & \textbf{at} (Heure / Point précis) \\
\hline
in the morning / afternoon / evening & on Monday / on Tuesdays & at 8 o'clock / at 10:15 \\
in January / in June & on 14 February / on 14th Feb & at noon / at midnight \\
in summer / in the dry season & on Christmas Day / on New Year's Eve & at night \\
in 2023 / in the 90s & on my birthday / on holiday & at the moment / at present \\
in two hours / in a minute & on the weekend (US) / at the weekend (GB) & at the beginning / at the end \\
\hline
\end{tabularx}
\end{center}

\begin{attention}[Pièges fréquents francophones — Temps]
\begin{itemize}
\item \textbf{« Le matin »} : \eng{in the morning} (général) mais \eng{on Monday morning} (précis). \textbf{Pas} \eng{at the morning}.
\item \textbf{« La nuit »} : \eng{at night} (période de sommeil, habitude) vs \eng{in the night} (au milieu d'une nuit précise : \eng{I woke up in the night}).
\item \textbf{« Le week-end »} : \eng{at the weekend} (GB, standard au Cameroun) / \eng{on the weekend} (US). Choisis \eng{at the weekend}.
\item \textbf{« À l'heure » / « À temps »} : \eng{on time} = ponctuel (selon l'horaire) ; \eng{in time} = assez tôt pour faire qqch (avant une échéance). \eng{The train arrived on time.} / \eng{We got there in time for the match.}
\item \textbf{Durée future} : \eng{in} + durée = « dans » (d'ici là). \eng{I'll be ready in 5 minutes.} « Je serai prêt dans 5 min. » Ne pas confondre avec \eng{after} (après un événement passé).
\end{itemize}
\end{attention}

\begin{propriete}[Autres prépositions de temps importantes]
\begin{itemize}
\item \cle{for} + durée : \eng{for two hours} (pendant deux heures — action accomplie ou habitude).
\item \cle{since} + point de départ : \eng{since Monday} / \eng{since 2010} (depuis — action qui continue, souvent \emph{present perfect}).
\item \cle{during} + nom (événement/période) : \eng{during the lesson}, \eng{during the holidays} (pendant — quand ?).
\item \cle{by} + échéance : \eng{by Friday} (au plus tard vendredi — \textbf{deadline}).
\item \cle{until / till} + fin : \eng{until 6 p.m.} (jusqu'à 18h — continuité).
\item \cle{from ... to / until} : \eng{from 8 to 4} (de 8 à 16h).
\end{itemize}
\end{propriete}

\begin{exemplebox}[Exemples traduits — Temps en contexte]
\eng{The final exam starts \textbf{at} 8 a.m. \textbf{on} Monday, \textbf{in} June.} « L'examen final commence à 8h le lundi, en juin. »

\eng{I have lived in Yaoundé \textbf{since} 2015 \textbf{for} nine years.} « J'habite Yaoundé depuis 2015 depuis neuf ans. »

\eng{Please submit your assignment \textbf{by} Friday.} « Veuillez rendre votre devoir pour vendredi au plus tard. »

\eng{We stayed awake \textbf{during} the night \textbf{until} dawn.} « Nous sommes restés éveillés pendant la nuit jusqu'à l'aube. »

\eng{The meeting lasted \textbf{for} three hours.} « La réunion a duré trois heures. »
\end{exemplebox}

\courssub{Direction, mouvement et manière}

\courssub{Mouvement et direction : to, into, onto, towards, through, across, along...}

Les prépositions de lieu deviennent des prépositions de \textbf{mouvement} quand elles suivent un verbe de déplacement (\eng{go, come, run, walk, drive, fly, fall, put, place...}). La logique spatiale (volume/surface/point) reste la même.

\begin{center}
\begin{tabularx}{\linewidth}{|X|X|X|}
\hline
\rowcolor{popblueL} Préposition & Logique spatiale & Exemple traduit \\
\hline
\eng{to} & Vers un point / destination & \eng{Go to school.} (direction générale) \\
\hline
\eng{into} & Vers l'intérieur (volume) & \eng{He walked into the room.} (il entre \emph{dans} la pièce) \\
\hline
\eng{onto} & Vers une surface & \eng{The cat jumped onto the table.} (il saute \emph{sur} la table) \\
\hline
\eng{towards} & Direction (pas nécessairement arrivée) & \eng{We are walking towards the market.} (on marche \emph{vers} le marché) \\
\hline
\eng{through} & Traversée d'un volume 3D (forêt, foule, tunnel) & \eng{The road goes through the forest.} (la route traverse la forêt) \\
\hline
\eng{across} & Traversée d'une surface 2D (rue, place, pont, rivière) & \eng{He swam across the river.} (il a traversé la rivière à la nage) \\
\hline
\eng{along} & Le long d'une ligne (route, rivière, couloir) & \eng{We walked along the beach.} (on a marché le long de la plage) \\
\hline
\eng{past} & Devant, en passant devant (point) & \eng{Walk past the church.} (passe devant l'église) \\
\hline
\eng{up / down} & Montée / Descente (pente, rue, escalier) & \eng{Go up the hill.} / \eng{Run down the stairs.} \\
\hline
\eng{out of} & Sortie d'un volume & \eng{Get out of the car.} (sors de la voiture) \\
\hline
\eng{off} & Descente d'une surface / véhicule « on » & \eng{Get off the bus.} (descends du bus) \\
\hline
\end{tabularx}
\end{center}

\begin{attention}[Verbes de mouvement + prépositions : pièges classiques]
\begin{itemize}
\item \textbf{Arrive} : \eng{arrive \textbf{at} a place} (bâtiment, point précis : \eng{at the station, at home}) vs \eng{arrive \textbf{in} a city/country} (grand espace : \eng{in Douala, in Cameroon}). \textbf{Jamais} \eng{arrive to}.
\item \textbf{Reach} : verbe transitif direct — \eng{reach the village} (pas de préposition !). \eng{We reached Yaoundé at noon.}
\item \textbf{Enter} : verbe transitif direct — \eng{enter the room} (pas de \eng{into} après \eng{enter}).
\item \textbf{Home} : \textbf{Jamais de préposition} avec \eng{home} (adverbe de lieu) : \eng{go home}, \eng{arrive home}, \eng{stay home}, \eng{come home}. \eng{Go to home} = FAUX.
\item \textbf{Get} : \eng{get in/into the car/taxi} MAIS \eng{get on/onto the bus/train/plane} ; \eng{get out of the car} MAIS \eng{get off the bus}.
\end{itemize}
\end{attention}

\courssub{La manière (Manner) : by, on, in, with, like}

La manière répond à la question « Comment ? » (\eng{How?}).

\begin{propriete}[Prépositions de manière]
\begin{itemize}
\item \cle{by} + \textbf{nom sans article} (moyen de transport, communication, méthode) : \eng{by car}, \eng{by bus}, \eng{by train}, \eng{by plane}, \eng{by phone}, \eng{by email}, \eng{by post}, \eng{by hand}, \eng{by mistake}, \eng{by chance}. \textbf{Exception :} \eng{on foot} (à pied — jamais \eng{by foot}).
\item \cle{on} + \textbf{pied / véhicule « on »} : \eng{on foot}, \eng{on horseback}.
\item \cle{in} + \textbf{véhicule « in » / état physique / langue} : \eng{in the car}, \eng{in a taxi}, \eng{in English}, \eng{in pencil}, \eng{in ink}, \eng{in a hurry}, \eng{in silence}.
\item \cle{with} + \textbf{instrument / accompagnement / manière (émotion)} : \eng{cut with a knife}, \eng{go with my brother}, \eng{with pleasure}, \eng{with difficulty}, \eng{with care}.
\item \cle{like} + \textbf{nom / pronom / gérondif} (comparaison, similitude) : \eng{He runs like the wind.}, \eng{She sings like a bird.}, \eng{Do it like this.}
\item \cle{Locutions prépositionnelles de manière} : \eng{in a ... way} (\eng{in a friendly way}), \eng{on purpose} (exprès), \eng{by accident} (par accident).
\end{itemize}
\end{propriete}

\begin{exemplebox}[Manière en contexte]
\eng{She came \textbf{by} bus but went back \textbf{on} foot because the traffic was bad.}
\eng{Write \textbf{in} pen, not \textbf{in} pencil.}
\eng{He solved the problem \textbf{with} ease / \textbf{in} an easy way.}
\eng{Don't behave \textbf{like} a child!}
\end{exemplebox}

\courssub{Comparaison français/anglais et pièges des francophones — Synthèse}

\begin{center}
\begin{tabularx}{\linewidth}{|X|X|X|}
\hline
\rowcolor{popblueL} Français (souvent unique) & Anglais (choix contextuel) & Exemple clé \\
\hline
\textbf{À} (lieu) & \eng{at} (point), \eng{in} (ville/pays), \eng{on} (rue US), \eng{to} (mouvement) & \eng{at the door}, \eng{in Paris}, \eng{to school} \\
\hline
\textbf{Dans} & \eng{in} (volume), \eng{into} (mouvement vers int.), \eng{during} (temps) & \eng{in the box}, \eng{go into the room}, \eng{during the night} \\
\hline
\textbf{Sur} & \eng{on} (surface/contact), \eng{onto} (mouvement vers surf.), \eng{about} (sujet) & \eng{on the table}, \eng{jump onto the bed}, \eng{a book about lions} \\
\hline
\textbf{De / Du / Des} (origine) & \eng{from} & \eng{I come from Bafoussam.} \\
\hline
\textbf{De / Du / Des} (appartenance, matière) & \eng{of} & \eng{The capital of Cameroon}, \eng{A cup of tea} \\
\hline
\textbf{Pour} (but, destinataire) & \eng{for} & \eng{This gift is for you.} \\
\hline
\textbf{Pour} (but, intention) & \eng{to} + infinitif & \eng{I study to succeed.} \\
\hline
\textbf{Par} (moyen, agent passif) & \eng{by} & \eng{by car}, \eng{written by Shakespeare} \\
\hline
\textbf{En} (durée, état, moyen) & \eng{in} (durée future, état), \eng{by} (moyen), \eng{on} (pied) & \eng{in 2024}, \eng{in trouble}, \eng{by train}, \eng{on foot} \\
\hline
\end{tabularx}
\end{center}

\begin{attention}[Top 5 des erreurs fatales au Bac]
\begin{enumerate}
\item \textbf{Confondre position et mouvement :} \eng{I am \textbf{to} Yaoundé} (FAUX) $\rightarrow$ \eng{I am \textbf{in} Yaoundé} / \eng{I go \textbf{to} Yaoundé}.
\item \textbf{Omettre la préposition après certains verbes :} \eng{listen music} (FAUX) $\rightarrow$ \eng{listen \textbf{to} music} ; \eng{wait the bus} (FAUX) $\rightarrow$ \eng{wait \textbf{for} the bus} ; \eng{look the picture} (FAUX) $\rightarrow$ \eng{look \textbf{at} the picture} ; \eng{enter the room} (CORRECT, transitif) mais \eng{go into the room}.
\item \textbf{Mauvaise préposition de transport :} \eng{in the bus}, \eng{in the train}, \eng{by foot} (FAUX) $\rightarrow$ \eng{on the bus}, \eng{on the train}, \eng{on foot} / \eng{by bus}.
\item \textbf{Traduction littérale « à » = \eng{at} partout :} \eng{at the morning}, \eng{at Monday}, \eng{at the wall} (FAUX) $\rightarrow$ \eng{in the morning}, \eng{on Monday}, \eng{on the wall}.
\item \textbf{Since / For :} \eng{I am here since two hours} (FAUX, présent + since) $\rightarrow$ \eng{I \textbf{have been} here \textbf{for} two hours} / \eng{\textbf{since} 10 a.m.}.
\end{enumerate}
\end{attention}

\courssub{Entraînement guidé et production — Lesson 8 Complete}

\begin{methode}[Comment réussir un exercice de prépositions (Gap-fill / Transformation)]
\begin{enumerate}
\item \textbf{Identifie le nom/groupe nominal} qui suit le trou (lieu ? temps ? moyen ?).
\item \textbf{Détermine la relation} : Position statique ? Mouvement vers ? Durée ? Moment précis ? Moyen ?
\item \textbf{Applique la logique spatiale/temporelle} : Volume (\eng{in}) / Surface (\eng{on}) / Point (\eng{at}).
\item \textbf{Vérifie les expressions figées} (\eng{at home}, \eng{on foot}, \eng{by car}, \eng{in the morning}, \eng{at night}).
\item \textbf{Relis la phrase complète} pour vérifier l'ordre des mots (Sujet + Verbe + Compléments : Place avant Temps ? Non, généralement \textbf{Place avant Temps} en fin de phrase : \eng{I met him \textbf{at the market} \textbf{yesterday}.})
\end{enumerate}
\end{methode}

\begin{exoresolu}[Exercices mixtes — Place, Temps, Mouvement, Manière]
Énoncé — Choisis la bonne préposition.

1. We arrived \dots\ Douala \dots\ midnight \dots\ a very old bus.
2. The keys are \dots\ the drawer \dots\ the kitchen.
3. She has been living in Bamenda \dots\ 2010.
4. Don't walk \dots\ the grass; use the path.
5. He threw the ball \dots\ the window and it broke.
6. I prefer travelling \dots\ train rather than \dots\ plane.
7. The meeting is \dots\ Monday \dots\ 10 a.m. sharp.
8. Divide the oranges \dots\ the three boys.

\tcblower

Solution détaillée :

1. \eng{We arrived \textbf{in} Douala \textbf{at} midnight \textbf{by} bus.} (Ville = \eng{in} ; Heure précise = \eng{at} ; Moyen transport = \eng{by}).
2. \eng{The keys are \textbf{in} the drawer \textbf{in} the kitchen.} (Tiroir = volume ; Pièce = volume). Note : \eng{on the drawer} si posées \emph{sur} le tiroir fermé.
3. \eng{She has been living in Bamenda \textbf{since} 2010.} (Point de départ au passé $\rightarrow$ \eng{since} + Present Perfect).
4. \eng{Don't walk \textbf{on} the grass; use the path.} (Surface = \eng{on}). \eng{Across the grass} possible si traversée.
5. \eng{He threw the ball \textbf{through} the window and it broke.} (Traversée volume 3D vitre/fenêtre ouverte = \eng{through}). \eng{Out of the window} aussi possible (sortie).
6. \eng{I prefer travelling \textbf{by} train rather than \textbf{by} plane.} (Moyen = \eng{by} + nom sans article).
7. \eng{The meeting is \textbf{on} Monday \textbf{at} 10 a.m. sharp.} (Jour = \eng{on} ; Heure = \eng{at}).
8. \eng{Divide the oranges \textbf{among} the three boys.} (Groupe de 3 = \eng{among}).
\end{exoresolu}

\begin{exercice}[Entraînement complet — À rendre]
\textbf{A. Complète avec la préposition appropriée (in, on, at, to, into, onto, by, on, for, since, during, until, among, between, opposite, behind).}

1. My father works \dots\ the Ministry \dots\ Yaoundé. He goes \dots\ work \dots\ car.
2. The match starts \dots\ 4 p.m. \dots\ Sunday. Be \dots\ the stadium \dots\ 3:30.
3. There is a big tree \dots\ our house and the river.
4. She was hiding \dots\ the crowd so the police couldn't see her.
5. I haven't seen him \dots\ Christmas. We will meet \dots\ the holidays.
6. The teacher walked \dots\ the classroom and sat \dots\ his desk.
7. Put the books \dots\ the shelf, please. Not \dots\ the floor.
8. We travelled \dots\ foot \dots\ the forest \dots\ two hours.

\textbf{B. Traduction (Français $\rightarrow$ Anglais).}
1. « Depuis combien de temps attends-tu ? J'attends depuis une heure. »
2. « Les enfants courent vers l'école. Ils sont pressés. »
3. « Ma maison est juste en face du marché, à côté de la pharmacie. »
4. « Il est monté dans le bus, s'est assis près de la fenêtre et a dormi pendant le trajet. »
5. « Ne marche pas sur l'herbe ! Fais le tour. »

\textbf{C. Production écrite (Writing).}
Rédige un court paragraphe (80-100 mots) pour décrire ton trajet de la maison à l'école. Utilise au moins : 3 prépositions de lieu, 3 de temps, 2 de direction/mouvement, 1 de manière. Souligne-les.
\end{exercice}

\begin{corrige}[Exercice « Entraînement complet »]
\textbf{A.}
1. \eng{My father works \textbf{at} the Ministry \textbf{in} Yaoundé. He goes \textbf{to} work \textbf{by} car.}
2. \eng{The match starts \textbf{at} 4 p.m. \textbf{on} Sunday. Be \textbf{at} the stadium \textbf{by} 3:30.} (\eng{by} = au plus tard).
3. \eng{There is a big tree \textbf{between} our house and the river.} (Deux éléments distincts).
4. \eng{She was hiding \textbf{among} the crowd so the police couldn't see her.} (Groupe/masse).
5. \eng{I haven't seen him \textbf{since} Christmas. We will meet \textbf{during} the holidays.}
6. \eng{The teacher walked \textbf{into} the classroom and sat \textbf{at} his desk.} (Mouvement vers int. $\rightarrow$ \eng{into} ; Position travail $\rightarrow$ \eng{at}).
7. \eng{Put the books \textbf{on} the shelf, please. Not \textbf{on} the floor.} (Surface étagère/sol).
8. \eng{We travelled \textbf{on} foot \textbf{through} the forest \textbf{for} two hours.} (\eng{on foot} figé ; \eng{through} volume forêt ; \eng{for} durée).

\textbf{B.}
1. \eng{How long have you been waiting? I have been waiting \textbf{for} an hour.}
2. \eng{The children are running \textbf{to} school. They are \textbf{in} a hurry.}
3. \eng{My house is just \textbf{opposite} the market, \textbf{next to} the pharmacy.}
4. \eng{He got \textbf{on} the bus, sat \textbf{by/near} the window, and slept \textbf{during} the journey.}
5. \eng{Don't walk \textbf{on} the grass! Go \textbf{around} / \eng{Walk \textbf{around} it.}}

\textbf{C. Modèle de production :}
\eng{Every morning, I wake up \textbf{at} 5:30 \textbf{in} the morning. I leave \textbf{at} 6:00 and walk \textbf{to} the main road \textbf{on} foot. It takes \textbf{about} 20 minutes. I wait \textbf{at} the bus stop \textbf{for} the bendskin. I sit \textbf{behind} the driver. We go \textbf{along} the expressway \textbf{through} the traffic. I arrive \textbf{at} school \textbf{at} 7:15, just \textbf{in} time for assembly.}
\end{corrige}

\begin{aretenir}[Synthèse finale — Lesson 8 : Prépositions et Locutions prépositionnelles]
\begin{itemize}
\item \textbf{LIEU (Position) :} \eng{in} (volume), \eng{on} (surface), \eng{at} (point). \eng{between} (2), \eng{among} (groupe). \eng{over/under} (verticalité/contact), \eng{above/below} (niveau). \eng{opposite} (en face), \eng{in front of} (devant).
\item \textbf{TEMPS :} \eng{in} (période large), \eng{on} (jour/date), \eng{at} (heure point). \eng{for} (durée), \eng{since} (début), \eng{by} (échéance), \eng{until} (fin), \eng{during} (quand).
\item \textbf{MOUVEMENT/DIRECTION :} \eng{to} (destination), \eng{into} (vers int.), \eng{onto} (vers surf.), \eng{towards} (direction), \eng{through} (traversée 3D), \eng{across} (traversée 2D), \eng{along} (le long), \eng{past} (devant), \eng{up/down}, \eng{out of/off} (sortie).
\item \textbf{MANIÈRE :} \eng{by} + transport/moyen (sans article), \eng{on foot}, \eng{in} + véhicule privé/langue/état, \eng{with} + instrument/émotion, \eng{like} (comparaison).
\item \textbf{RÈGLES D'OR :} \eng{arrive at/in} (jamais \eng{to}), \eng{reach/enter} (sans prép.), \eng{go home} (sans prép.), \eng{get on/off} (bus) vs \eng{get in/out of} (car), \eng{at the weekend} (GB), \eng{on time} vs \eng{in time}.
\end{itemize}
\end{aretenir}

\begin{savaistu}[Culture \& Langue : Les prépositions dans l'anglais camerounais (CamE)]
L'anglais camerounais (Cameroon English) suit les normes britanniques pour l'orthographe et la grammaire standard (ex. \eng{at the weekend}, \eng{in hospital}, \eng{programme}, \eng{colour}). Cependant, l'influence du français et des langues locales crée des particularités :
\begin{itemize}
\item \eng{At} s'étend souvent là où l'anglais standard utiliserait \eng{in} ou \eng{on} : \eng{at the market} (standard), mais on entend \eng{at the house} (pour \eng{in the house}), \eng{at the road} (pour \eng{on the road}).
\item \eng{Since} utilisé pour \eng{for} : \eng{I am here since two hours} (erreur courante, à corriger : \eng{for two hours}).
\item \eng{By} utilisé pour \eng{at} (heure) : \eng{Come by 8 o'clock} (standard : \eng{at 8} ou \eng{by 8} si échéance).
\item \textbf{Conseil :} Pour les examens officiels (Bac, Probatoire, GCE), utilise l'\textbf{anglais standard britannique} présenté dans cette leçon.
\end{itemize}
\end{savaistu}

\coursec{Prépositions de temps (time)}

Dans la section précédente, nous avons vu que les prépositions de lieu (\eng{in}, \eng{on}, \eng{at}, \eng{next to}, \eng{between}...) permettent de situer les personnes et les objets \emph{dans l'espace}. Bonne nouvelle : la logique est presque la même pour situer les événements \emph{dans le temps}. En anglais comme en français, on ne dit pas n'importe quoi : on ne peut pas traduire mécaniquement « le lundi » par \eng{in Monday} ni « en 2024 » par \eng{on 2024}. Cette section te donne les clés pour ne plus jamais hésiter devant une date, une heure ou une durée.

\courssub{Observation : un emploi du temps bien rempli}

\begin{textebox}[A busy week for Amina]
Amina is a Form Five student at Government High School in Buea. Her week is always very busy. She wakes up \eng{at 5.30 a.m.} every morning and arrives at school \eng{at 7 o'clock}. \eng{On Monday}, she has mathematics \eng{in the morning} and sports \eng{in the afternoon}. \eng{On 15th May}, she will take her English test. She has lived in Buea \eng{since 2020}, and she has studied English \eng{for six years}. \eng{During the holidays}, she visits her grandmother in Bamenda. Classes run \eng{from 8 a.m. to 3 p.m.}, and students must finish their homework \eng{by Friday}. \eng{At night}, Amina reads her notes before sleeping. She was born \eng{in 2008}, \eng{in December}, and she always celebrates her birthday \eng{at Christmas time} with the whole family.
\end{textebox}

\textbf{Glossaire}

\begin{itemize}
\item \eng{to wake up} (verbe) : se réveiller
\item \eng{to take a test} (verbe) : passer un test
\item \eng{homework} (nom, indénombrable) : les devoirs
\item \eng{to run} (verbe, ici) : se dérouler
\end{itemize}

\textbf{Questions d'observation}

\begin{enumerate}
\item Quelle préposition est utilisée devant les heures (\eng{5.30 a.m.}, \eng{7 o'clock}) ?
\item Quelle préposition est utilisée devant les jours et les dates (\eng{Monday}, \eng{15th May}) ?
\item Quelle préposition est utilisée devant les mois, les années et les saisons (\eng{December}, \eng{2008}) ?
\item Relevez les deux expressions avec \eng{since} et \eng{for}. Quelle différence voyez-vous ?
\end{enumerate}

\textbf{Réponses attendues} : 1. \eng{at} ; 2. \eng{on} ; 3. \eng{in} ; 4. \eng{since 2020} désigne un point de départ, \eng{for six years} désigne une durée.

\courssub{La triade in / on / at pour le temps}

\begin{propriete}[La triade temporelle : at, on, in]
Pour exprimer le moment d'un événement, l'anglais utilise trois prépositions selon le degré de précision :
\begin{itemize}
\item \cle{at} : un \textbf{point précis} dans le temps — l'heure exacte et les moments ponctuels : \eng{at 7 a.m.}, \eng{at noon} (à midi), \eng{at midnight}, \eng{at night}, \eng{at Christmas}, \eng{at the weekend} (GB).
\item \cle{on} : un \textbf{jour ou une date} : \eng{on Monday}, \eng{on 15th May}, \eng{on my birthday}, \eng{on Christmas Day}, \eng{on Independence Day}.
\item \cle{in} : une \textbf{période longue} — mois, années, saisons, siècles, parties de la journée : \eng{in July}, \eng{in 2024}, \eng{in the rainy season}, \eng{in the 21st century}, \eng{in the morning}.
\end{itemize}
Image mentale : plus la période est grande, plus on « entre dedans » (\eng{in}) ; plus elle est précise, plus c'est un point (\eng{at}).
\end{propriete}

\begin{popfigure}
\begin{tikzpicture}
\fill[popgreenL] (0,0) -- (6,0) -- (4.6,1.5) -- (1.4,1.5) -- cycle;
\fill[popgoldL] (1.4,1.7) -- (4.6,1.7) -- (3.8,2.9) -- (2.2,2.9) -- cycle;
\fill[poporange] (2.2,3.1) -- (3.8,3.1) -- (3,4.3) -- cycle;
\node[align=center] at (3,0.75) {\textbf{IN} — mois, années, saisons\\ \eng{in July, in 2024, in the morning}};
\node[align=center] at (3,2.3) {\textbf{ON} — jours, dates\\ \eng{on Monday, on 15th May}};
\node[align=center, text=white] at (3,3.55) {\textbf{AT} — heure précise};
\node[align=center] at (7.1,3.55) {\eng{at 7 a.m.}\\ \eng{at noon, at night}};
\end{tikzpicture}
\legende{La pyramide de la triade temporelle : \eng{at} au sommet (le plus précis), \eng{in} à la base (les périodes les plus longues).}
\end{popfigure}

\begin{exemplebox}[La triade en action]
\begin{itemize}
\item \eng{The exam is on 12th June.} « L'examen a lieu le 12 juin. »
\item \eng{School starts at 7.30 a.m. in Cameroon.} « Les cours commencent à 7 h 30 au Cameroun. »
\item \eng{The rainy season begins in March.} « La saison des pluies commence en mars. »
\item \eng{My grandfather was born in 1955.} « Mon grand-père est né en 1955. »
\item \eng{We celebrate Christmas at midnight.} « Nous fêtons Noël à minuit. »
\item \eng{I always do my homework in the evening.} « Je fais toujours mes devoirs le soir. »
\end{itemize}
\end{exemplebox}

\begin{attention}[Les exceptions de la triade]
Deux cas piègent régulièrement les élèves :
\begin{enumerate}
\item On dit \eng{in the morning}, \eng{in the afternoon}, \eng{in the evening}, \textbf{mais} \eng{at night} (et non \eng{in the night} pour « la nuit » en général).
\item Quand un jour est combiné avec une partie de la journée, c'est le \textbf{jour} qui commande : \eng{on Monday morning}, \eng{on Friday evening}, \eng{on the afternoon of 5th June} — jamais \eng{in Monday morning}.
\end{enumerate}
\end{attention}

\begin{center}
\begin{tabularx}{\linewidth}{|X|X|X|}\hline
\rowcolor{popblueL} \textbf{AT (point précis)} & \textbf{ON (jour / date)} & \textbf{IN (période longue)} \\ \hline
\eng{at 6 o'clock} (à 6 heures) & \eng{on Tuesday} (mardi) & \eng{in April} (en avril) \\ \hline
\eng{at noon / at midnight} & \eng{on 1st January} (le 1er janvier) & \eng{in 2024} (en 2024) \\ \hline
\eng{at night} (la nuit) & \eng{on my birthday} (à mon anniversaire) & \eng{in the dry season} (en saison sèche) \\ \hline
\eng{at Christmas} (à Noël) & \eng{on Christmas Day} (le jour de Noël) & \eng{in the morning} (le matin) \\ \hline
\eng{at the weekend} (GB) & \eng{on Monday morning} (lundi matin) & \eng{in the 20th century} (au XXe siècle) \\ \hline
\end{tabularx}
\end{center}

\courssub{Exprimer la durée : for, since, during, from...to}

\begin{propriete}[For, since, during]
\begin{itemize}
\item \cle{for} + une \textbf{durée} (un segment de temps) : \eng{for two weeks}, \eng{for six years}, \eng{for a long time}.
\item \cle{since} + un \textbf{point de départ} (une date, un événement) : \eng{since 2020}, \eng{since Monday}, \eng{since I was a child}. Avec \eng{since}, on utilise presque toujours le \emph{present perfect}.
\item \cle{during} + un \textbf{nom} (un événement, une période nommée) : \eng{during the holidays}, \eng{during the lesson}, \eng{during the match}. Attention : \eng{during} est suivi d'un \emph{nom}, jamais d'une durée chiffrée.
\item \cle{from...to/until} : encadrent une période : \eng{from 8 a.m. to 4 p.m.}, \eng{from Monday to Friday}.
\end{itemize}
\end{propriete}

\begin{popfigure}
\begin{tikzpicture}
\draw[-{Stealth}, thick, popdark] (0,0) -- (11,0);
\foreach \x/\lab in {1/{2020}, 5/{2022}, 9/{2024}}
  \draw[popdark] (\x,-0.12) -- (\x,0.12) node[above=2pt] {\lab};
\fill[poporange] (1,0) circle (0.14);
\node[above=0.55cm, align=center, poporange] at (1,0) {\textbf{since 2020}\\ point de départ};
\draw[very thick, popgreen] (1.3,0.3) -- (8.7,0.3);
\draw[very thick, popgreen] (1.3,0.15) -- (1.3,0.45);
\draw[very thick, popgreen] (8.7,0.15) -- (8.7,0.45);
\node[above=0.75cm, align=center, popgreen] at (5,0) {\textbf{for four years}\\ durée};
\fill[poppurple] (9,0) circle (0.14);
\node[above=0.55cm, align=center, poppurple] at (9.6,0) {\textbf{by / until Friday}\\ limite finale};
\end{tikzpicture}
\legende{Frise du temps : \eng{since} marque le point de départ, \eng{for} mesure le segment de durée, \eng{by} et \eng{until} fixent la limite finale.}
\end{popfigure}

\begin{exemplebox}[Durée et limites]
\begin{itemize}
\item \eng{I have lived here since 2019.} « J'habite ici depuis 2019. »
\item \eng{We studied for three hours.} « Nous avons étudié pendant trois heures. »
\item \eng{During the holidays, I help my father at the market.} « Pendant les vacances, j'aide mon père au marché. »
\item \eng{The library is open from 9 a.m. to 5 p.m.} « La bibliothèque est ouverte de 9 h à 17 h. »
\item \eng{Finish your composition by Friday.} « Terminez votre rédaction au plus tard vendredi. »
\item \eng{Wait until Friday, please.} « Attendez jusqu'à vendredi, s'il vous plaît. »
\end{itemize}
\end{exemplebox}

\begin{propriete}[By ou until ?]
\begin{itemize}
\item \cle{by} = « au plus tard » : l'action doit être \textbf{terminée} à ce moment-là. \eng{Send the letter by Monday.} « Envoie la lettre au plus tard lundi. »
\item \cle{until} (ou \eng{till}, familier) = « jusqu'à » : l'action ou l'état \textbf{continue} jusqu'à ce moment. \eng{The shop stays open until 8 p.m.} « La boutique reste ouverte jusqu'à 20 h. »
\end{itemize}
Compare : \eng{I must finish by 5 p.m.} (à 17 h, c'est fini) / \eng{I worked until 5 p.m.} (j'ai travaillé sans arrêt jusqu'à 17 h).
\end{propriete}

\courssub{Ago, before, after : situer avant et après}

\begin{propriete}[Ago, before, after]
\begin{itemize}
\item \cle{ago} = « il y a » : se place \textbf{après} la durée, contrairement au français ! \eng{two days ago} « il y a deux jours » ; \eng{a long time ago} « il y a longtemps ». S'emploie avec le \emph{prétérit}.
\item \cle{before} = « avant » : \eng{before the exam} « avant l'examen » ; \eng{before 6 o'clock} « avant 6 heures ».
\item \cle{after} = « après » : \eng{after lunch} « après le déjeuner » ; \eng{after school} « après l'école ».
\end{itemize}
\end{propriete}

\begin{exemplebox}[Avant et après]
\begin{itemize}
\item \eng{My uncle left Douala three years ago.} « Mon oncle a quitté Douala il y a trois ans. »
\item \eng{Wash your hands before the meal.} « Lave-toi les mains avant le repas. »
\item \eng{We play football after school.} « Nous jouons au football après l'école. »
\end{itemize}
\end{exemplebox}

\courssub{Les cas sans préposition}

\begin{propriete}[Zéro préposition]
En anglais, on ne met \textbf{jamais} de préposition devant les expressions de temps introduites par \cle{last}, \cle{next}, \cle{this}, \cle{that}, \cle{every}, \cle{today}, \cle{yesterday}, \cle{tomorrow} :
\begin{itemize}
\item \eng{I saw her last Monday.} « Je l'ai vue lundi dernier. » (jamais \eng{on last Monday})
\item \eng{We are travelling next year.} « Nous voyagerons l'année prochaine. » (jamais \eng{in next year})
\item \eng{I have a test this morning.} « J'ai un test ce matin. »
\item \eng{She calls me every Sunday.} « Elle m'appelle chaque dimanche. »
\item \eng{See you tomorrow!} « À demain ! »
\end{itemize}
\end{propriete}

\begin{attention}[Erreurs fréquentes des francophones]
\begin{enumerate}
\item « Depuis » se traduit par \eng{since} (point de départ) ou \eng{for} (durée), \textbf{jamais} par \eng{during} : \eng{I have known her for ten years.} « Je la connais depuis dix ans. » — et non \eng{during ten years}.
\item Ne traduis pas « il y a » mot à mot : « il y a deux jours » = \eng{two days ago} (l'ordre est inversé !).
\item Pas de préposition devant \eng{last}, \eng{next}, \eng{this}, \eng{every} : \eng{last week}, pas \eng{on last week}.
\item « Le matin » = \eng{in the morning} avec l'article \eng{the} obligatoire, mais « la nuit » = \eng{at night} \emph{sans} article.
\item « Pendant » + durée se traduit par \eng{for} : \eng{I slept for eight hours.} « J'ai dormi pendant huit heures. » — et non \eng{during eight hours}.
\end{enumerate}
\end{attention}

\begin{savaistu}[British or American?]
Pour parler du week-end, les Britanniques disent \eng{at the weekend} tandis que les Américains disent \eng{on the weekend}. Les deux formes sont correctes : c'est une simple question de variété d'anglais. Mentionner cette différence dans une copie ou à l'oral peut impressionner l'examinateur au Bac ! Autre curiosité : les Américains écrivent les dates « mois + jour » (\eng{May 15th}) alors que les Britanniques préfèrent « jour + mois » (\eng{15th May}) — d'où l'importance de bien choisir \eng{on} devant toutes les dates, quelle que soit la variété.
\end{savaistu}

\begin{exoresolu}[Complète avec la bonne préposition — ou rien]
Énoncé — Complète avec \eng{at}, \eng{on}, \eng{in}, \eng{for}, \eng{since}, \eng{during}, \eng{by}, \eng{until}, \eng{ago} ou \eng{Ø} (aucune préposition) :
\begin{enumerate}
\item The plane landed \dots\ 6.45 p.m.
\item My aunt has lived in Yaoundé \dots\ 2015.
\item We played football \dots\ two hours.
\item \dots\ last Friday, we visited the museum.
\item The meeting is \dots\ 3rd March.
\item I was born \dots\ 2009.
\item You must return the books \dots\ Thursday.
\item She arrived in Douala two weeks \dots.
\end{enumerate}
\tcblower
Solution détaillée :
\begin{enumerate}
\item \eng{at} — heure précise : \eng{at 6.45 p.m.}
\item \eng{since} — point de départ avec present perfect : \eng{since 2015}.
\item \eng{for} — durée : \eng{for two hours} (jamais \eng{during} + durée chiffrée).
\item \eng{Ø} — aucune préposition devant \eng{last} : \eng{Last Friday, we visited the museum.}
\item \eng{on} — date : \eng{on 3rd March}.
\item \eng{in} — année : \eng{in 2009}.
\item \eng{by} — limite « au plus tard » : \eng{by Thursday}.
\item \eng{ago} — « il y a deux semaines » : \eng{two weeks ago} (attention à l'ordre inversé).
\end{enumerate}
\end{exoresolu}

\begin{exercice}[À toi de jouer]
\begin{enumerate}
\item Traduis en anglais : « Mon frère étudie l'anglais depuis cinq ans. » / « Le match commence samedi à 16 heures. » / « Nous sommes arrivés il y a une heure. » / « Je me lève à 5 h 30 le matin. »
\item Corrige les erreurs : \eng{I have lived here during ten years.} / \eng{She was born on 2005.} / \eng{We met on last Monday.} / \eng{He sleeps in the night.}
\item Rédige cinq phrases sur ta semaine en utilisant au moins cinq prépositions de temps différentes (\eng{at}, \eng{on}, \eng{in}, \eng{for}, \eng{until}, \eng{ago}...).
\end{enumerate}
\end{exercice}

\begin{corrige}[Exercice : À toi de jouer]
\begin{enumerate}
\item \eng{My brother has studied English for five years.} / \eng{The match starts on Saturday at 4 p.m.} / \eng{We arrived an hour ago.} / \eng{I get up at 5.30 in the morning.}
\item \eng{I have lived here for ten years.} / \eng{She was born in 2005.} / \eng{We met last Monday.} / \eng{He sleeps at night.}
\item Modèle : \eng{On Monday, I have maths at 8 a.m. I have studied English for six years. I usually read in the evening. Last week, I visited my cousins. I do my homework until 9 p.m.}
\end{enumerate}
\end{corrige}

\begin{aretenir}[L'essentiel des prépositions de temps]
\begin{itemize}
\item \eng{at} = point précis (heure, \eng{noon}, \eng{night}) ; \eng{on} = jour ou date ; \eng{in} = mois, année, saison, partie de la journée.
\item \eng{since} + point de départ ; \eng{for} + durée ; \eng{during} + nom d'événement ; \eng{from...to} pour encadrer.
\item \eng{by} = au plus tard ; \eng{until} = jusqu'à.
\item \eng{ago} se place après la durée : \eng{three days ago}.
\item Jamais de préposition devant \eng{last}, \eng{next}, \eng{this}, \eng{every}, \eng{today}, \eng{yesterday}, \eng{tomorrow}.
\end{itemize}
\end{aretenir}

\coursec{Direction, mouvement et manière}

Dans la section précédente, nous avons étudié les prépositions de \emph{temps} (\eng{at}, \eng{on}, \eng{in}, \eng{since}, \eng{for}\ldots), qui répondent à la question « quand ? ». Nous passons maintenant à un autre grand emploi des prépositions : exprimer le \cle{mouvement} (d'où l'on vient, où l'on va, par où l'on passe) et la \cle{manière} (comment l'action est faite, avec quel moyen). Ces prépositions sont essentielles pour raconter un voyage, décrire un trajet à Douala ou expliquer comment on se rend au lycée.

\courssub{Observation : un trajet à Douala}

\begin{textebox}[Getting to my uncle's house — a short narrative]
Last Saturday, I left home early and walked \textbf{to} the bus station. I got \textbf{into} the bus and sat near the window. The bus drove \textbf{along} the main road, went \textbf{through} a small tunnel, and went \textbf{over} a bridge. It passed \textbf{through} several villages and drove \textbf{past} a big market without stopping. After two hours, we arrived \textbf{in} Douala. I got \textbf{out of} the bus and walked \textbf{towards} my uncle's house. I walked \textbf{across} a busy street, climbed \textbf{up} a small hill, and finally reached his gate. My uncle was waiting for me \textbf{with} a big smile. We travelled \textbf{from} the bus station \textbf{to} his house \textbf{on foot}, because he lives nearby. In the evening, we spoke \textbf{in} English and laughed a lot. My uncle works \textbf{as} a mechanic, and he sings \textbf{like} a professional singer!
\end{textebox}

\textbf{Glossaire :} \emph{gate} (n., portail) ; \emph{nearby} (adv., tout près) ; \emph{to reach} (v., atteindre) ; \emph{without stopping} (sans s'arrêter).

\textbf{Questions d'observation :}
\begin{enumerate}
\item Souligne toutes les prépositions suivies d'un lieu. Lesquelles expriment un \emph{mouvement} ? Lesquelles expriment une \emph{position} ?
\item Pourquoi dit-on \eng{arrived in Douala} et non \eng{arrived to Douala} ?
\item Comment le narrateur se déplace-t-il à la fin : \eng{by foot} ou \eng{on foot} ?
\end{enumerate}

\courssub{Direction et mouvement : to, into, out of, towards, from...to}

\begin{propriete}[La règle du mouvement]
\begin{itemize}
\item \cle{to} marque la \textbf{destination} : \eng{He went to school.} (Il est allé à l'école.)
\item \cle{into} marque l'\textbf{entrée dans} un espace (mouvement vers l'intérieur) : \eng{The children ran into the classroom.} (Les enfants sont entrés en courant dans la salle.)
\item \cle{out of} marque la \textbf{sortie} d'un espace : \eng{She ran out of the house.} (Elle est sortie de la maison en courant.)
\item \cle{towards} marque la \textbf{direction}, sans certitude d'arrivée : \eng{He walked towards the market.} (Il a marché en direction du marché — on ne sait pas s'il y est arrivé.)
\item \cle{from...to} indique l'\textbf{origine et la destination} : \eng{We travelled from Yaoundé to Buea.} (Nous avons voyagé de Yaoundé à Buéa.)
\end{itemize}
\textbf{Exception importante :} après \eng{arrive}, on utilise \eng{at} ou \eng{in}, jamais \eng{to} : \eng{We arrived in Douala.} De même, \eng{home} ne prend pas de préposition : \eng{I went home} (pas \eng{to home}).
\end{propriete}

\begin{exemplebox}[Exemples traduits]
\begin{itemize}
\item \eng{The children ran into the classroom.} = Les enfants sont entrés en courant dans la salle.
\item \eng{We travelled from Yaoundé to Buea.} = Nous avons voyagé de Yaoundé à Buéa.
\item \eng{The cat jumped out of the box.} = Le chat a sauté hors de la boîte.
\item \eng{She walked towards the river.} = Elle a marché vers la rivière.
\item \eng{They went home after the match.} = Ils sont rentrés à la maison après le match.
\end{itemize}
\end{exemplebox}

\begin{attention}[Arrive + at/in, jamais to]
Les francophones traduisent souvent « arriver à » par \eng{arrive to}. C'est une erreur grave. On dit :
\begin{itemize}
\item \eng{We arrived \textbf{in} Douala at six o'clock.} (ville, pays → \eng{in})
\item \eng{We arrived \textbf{at} the station late.} (lieu précis → \eng{at})
\end{itemize}
De même, \eng{to go to school} (mouvement) se distingue de \eng{to be at school} (position) : \eng{He went \textbf{to} school} (il s'est déplacé) mais \eng{He is \textbf{at} school} (il y est).
\end{attention}

\courssub{Le trajet : across, through, along, past, over, up, down}

Ces prépositions décrivent \emph{par où} passe le mouvement. La distinction la plus délicate pour un francophone est celle entre \cle{across} et \cle{through} :

\begin{propriete}[Across ou through ?]
\begin{itemize}
\item \cle{across} = d'un côté à l'autre d'une \textbf{surface plane} (une route, un pont, un terrain, une rivière) : \eng{She walked across the road.} (Elle a traversé la route.)
\item \cle{through} = à travers un \textbf{volume, un espace en trois dimensions} (une forêt, un tunnel, une foule) : \eng{The bus went through the forest.} (Le bus a traversé la forêt.)
\item \cle{along} = le long de : \eng{They walked along the river.} (Ils ont marché le long de la rivière.)
\item \cle{past} = devant, sans s'arrêter : \eng{We drove past the stadium.} (Nous sommes passés devant le stade.)
\item \cle{over} = par-dessus : \eng{The bird flew over the roof.} (L'oiseau a volé par-dessus le toit.)
\item \cle{up / down} = en montant / en descendant : \eng{He climbed up the hill.} / \eng{She ran down the stairs.}
\end{itemize}
\end{propriete}

\begin{popfigure}
\begin{center}
\begin{tikzpicture}[scale=0.9, >=Stealth]
% to : A -> B
\fill[popblue] (0,0) circle (0.12) node[below=2pt, popdark]{A};
\fill[poporange] (2.6,0) circle (0.12) node[below=2pt, popdark]{B};
\draw[->, very thick, popblue] (0.25,0) -- (2.35,0) node[midway, above, popdark]{to};
% into / out of : cercle
\draw[thick, poppurple] (5.6,0) circle (0.85);
\draw[->, very thick, popgreen] (3.9,0) -- (4.95,0) node[midway, above, popdark]{into};
\draw[->, very thick, poporange] (6.25,0) -- (7.4,0) node[midway, above, popdark]{out of};
% across : rectangle traversé
\draw[thick, popblue, fill=popblueL] (8.6,-0.5) rectangle (10.4,0.5);
\draw[->, very thick, popdark] (8.3,0) -- (10.7,0) node[midway, above=2pt, popdark]{across};
% through : tunnel
\draw[thick, popgold, fill=popgoldL] (12,-0.5) -- (13.6,-0.5) arc(-90:90:0.5) -- (12,0.5) arc(90:270:0.5);
\draw[->, very thick, popdark] (11.6,0) -- (14,0) node[midway, above=2pt, popdark]{through};
\end{tikzpicture}
\end{center}
\legende{Schéma des prépositions de mouvement : \eng{to} (destination), \eng{into}/\eng{out of} (entrée/sortie), \eng{across} (surface plane), \eng{through} (volume).}
\end{popfigure}

\begin{center}
\begin{tabularx}{\linewidth}{|X|X|X|}
\hline
\rowcolor{popblueL} Préposition & Type d'espace & Exemple traduit \\
\hline
across & surface plane (2 dimensions) & \eng{He swam across the river.} = Il a traversé la rivière à la nage. \\
\hline
through & volume, espace fermé (3 dimensions) & \eng{The train went through the tunnel.} = Le train a traversé le tunnel. \\
\hline
along & une ligne, une longueur & \eng{We walked along the beach.} = Nous avons marché le long de la plage. \\
\hline
past & devant un point, sans arrêt & \eng{She went past my house.} = Elle est passée devant chez moi. \\
\hline
over & par-dessus un obstacle & \eng{He jumped over the fence.} = Il a sauté par-dessus la clôture. \\
\hline
up / down & montée / descente & \eng{They walked up the hill.} = Ils ont monté la colline. \\
\hline
\end{tabularx}
\end{center}

\courssub{Manière et moyen : by, on foot, with, in}

\begin{propriete}[By + moyen de transport, sans article]
Pour exprimer le \textbf{moyen de transport}, on utilise \cle{by} + nom \textbf{sans article} : \eng{by bus}, \eng{by car}, \eng{by plane}, \eng{by train}, \eng{by bike}, \eng{by bendskin}.
\begin{itemize}
\item \eng{She goes to school by bus.} = Elle va à l'école en bus.
\end{itemize}
\textbf{Exception célèbre :} \cle{on foot} (à pied) — jamais \eng{by foot} !
\end{propriete}

\begin{propriete}[With + instrument ; in + langue ou façon]
\begin{itemize}
\item \cle{with} + \textbf{instrument} (avec article) : \eng{Cut the bread with a knife.} (Coupe le pain avec un couteau.) ; \eng{He wrote the letter with a pen.}
\item \cle{in} + \textbf{langue} : \eng{She answered in English.} (Elle a répondu en anglais.)
\item \cle{in} + \textbf{façon, manière} : \eng{He spoke in a low voice.} (Il a parlé à voix basse.) ; \eng{She answered with a smile.} (Elle a répondu avec un sourire.)
\item \cle{without} = sans : \eng{He left without a word.} (Il est parti sans un mot.)
\end{itemize}
\end{propriete}

\begin{exemplebox}[Exemples traduits]
\begin{itemize}
\item \eng{She goes to school on foot.} = Elle va à l'école à pied.
\item \eng{Cut the bread with a knife.} = Coupe le pain avec un couteau.
\item \eng{He spoke in a low voice.} = Il a parlé à voix basse.
\item \eng{They travelled to Bamenda by car.} = Ils ont voyagé à Bamenda en voiture.
\item \eng{He left without a word.} = Il est parti sans un mot.
\end{itemize}
\end{exemplebox}

\begin{attention}[By foot et with the car : deux erreurs classiques]
\begin{itemize}
\item « à pied » se dit \eng{\textbf{on} foot}, jamais \eng{by foot} : \eng{I go to school on foot.}
\item Pour le moyen de transport, on n'utilise pas \eng{with} : « en voiture » = \eng{\textbf{by} car}, pas \eng{with the car}. En revanche, \eng{with} convient pour l'instrument : \eng{with a knife}, \eng{with a pen}.
\item Astuce : transport → \eng{by} (sans article) ; instrument → \eng{with} (avec article) ; langue → \eng{in}.
\end{itemize}
\end{attention}

\courssub{Like ou as ?}

\begin{propriete}[Like = comparaison ; as = fonction]
\begin{itemize}
\item \cle{like} = « comme », pour une \textbf{comparaison} (il n'est pas vraiment un ange, il chante seulement comme un ange) : \eng{He sings like an angel.} (Il chante comme un ange.) ; \eng{She runs like a cheetah.} (Elle court comme un guépard.)
\item \cle{as} = « en tant que », pour une \textbf{fonction, un rôle réel} : \eng{He works as a teacher.} (Il travaille comme professeur — c'est bien son métier.) ; \eng{I use this room as an office.} (J'utilise cette pièce comme bureau.)
\end{itemize}
Test : si la personne \emph{est réellement} ce que dit le nom qui suit, utilise \eng{as} ; si c'est une simple image, utilise \eng{like}.
\end{propriete}

\begin{center}
\begin{tabularx}{\linewidth}{|X|X|X|}
\hline
\rowcolor{popblueL} Français & English & Remarque \\
\hline
à pied & on foot & jamais \eng{by foot} \\
\hline
en voiture / en bus & by car / by bus & sans article \\
\hline
avec un couteau & with a knife & instrument → \eng{with} \\
\hline
en anglais & in English & langue → \eng{in} \\
\hline
à voix basse & in a low voice & manière → \eng{in} \\
\hline
sans un mot & without a word & \eng{without} + nom \\
\hline
comme un ange (image) & like an angel & comparaison \\
\hline
comme professeur (métier) & as a teacher & fonction réelle \\
\hline
\end{tabularx}
\end{center}

\begin{exoresolu}[Complète avec la préposition correcte]
Énoncé : complète avec \eng{to}, \eng{into}, \eng{out of}, \eng{through}, \eng{across}, \eng{by}, \eng{on}, \eng{with}, \eng{in}, \eng{like} ou \eng{as}.
\begin{enumerate}
\item The boys ran \_\_\_ the classroom when the bell rang.
\item My father goes to work \_\_\_ car, but I go \_\_\_ foot.
\item The train passed \_\_\_ a long tunnel.
\item She works \_\_\_ a nurse at the hospital, and she sings \_\_\_ a bird.
\item Please write your name \_\_\_ a pen, and answer \_\_\_ English.
\end{enumerate}
\tcblower
Solution détaillée :
\begin{enumerate}
\item \eng{into} — entrée dans un espace (mouvement vers l'intérieur de la salle).
\item \eng{by} car (moyen de transport, sans article) ; \eng{on} foot (exception obligatoire : jamais \eng{by foot}).
\item \eng{through} — le tunnel est un volume en trois dimensions (une surface plane exigerait \eng{across}).
\item \eng{as} a nurse (fonction réelle : c'est son métier) ; \eng{like} a bird (simple comparaison : elle n'est pas un oiseau).
\item \eng{with} a pen (instrument, avec article) ; \eng{in} English (langue).
\end{enumerate}
\end{exoresolu}

\begin{exercice}[À ton tour]
\begin{enumerate}
\item Traduis : « Nous sommes arrivés à Douala à huit heures. » (attention à \eng{arrive} !)
\item Traduis : « Elle va au marché à pied et revient en taxi. »
\item Choisis : \eng{The swimmer crossed the pool (across / through).} Justifie ta réponse.
\item Choisis : \eng{He works (like / as) a driver in Buea.} Justifie.
\item Rédige trois phrases sur ton trajet quotidien vers le lycée en utilisant \eng{from...to}, \eng{along} et \eng{past}.
\end{enumerate}
\end{exercice}

\begin{corrige}[Exercice « À ton tour »]
\begin{enumerate}
\item \eng{We arrived in Douala at eight o'clock.} (ville → \eng{in} ; heure → \eng{at} ; jamais \eng{arrived to}.)
\item \eng{She goes to the market on foot and comes back by taxi.}
\item \eng{across} — la piscine est traversée comme une surface, d'un bord à l'autre.
\item \eng{as} — « driver » est sa fonction réelle, son métier.
\item Modèle : \eng{I walk from my house to the lycée. I walk along the main road and I go past the stadium every morning.}
\end{enumerate}
\end{corrige}

\begin{aretenir}[L'essentiel de la section]
\begin{itemize}
\item Mouvement : \eng{to} (destination), \eng{into} (entrée), \eng{out of} (sortie), \eng{towards} (direction), \eng{from...to} (origine-destination).
\item Trajet : \eng{across} = surface plane ; \eng{through} = volume ; \eng{along} = le long de ; \eng{past} = devant ; \eng{over} = par-dessus ; \eng{up/down} = montée/descente.
\item Moyen et manière : \eng{by} + transport sans article, mais \eng{on foot} ; \eng{with} + instrument ; \eng{in} + langue ou voix ; \eng{without} = sans.
\item \eng{like} = comparaison (comme) ; \eng{as} = fonction réelle (en tant que).
\item Pièges : \eng{arrive at/in} (jamais \eng{to}) ; \eng{on foot} (jamais \eng{by foot}) ; \eng{by car} (jamais \eng{with the car}).
\end{itemize}
\end{aretenir}

\coursec{Observation et découverte}

\begin{textebox}[Dialogue : Une réunion de famille à Yaoundé]
\textbf{Context :} The Ngo family gathers for Sunday lunch at their home in Mvog-Ada, Yaoundé. Auntie Marie has just arrived from Douala.

\eng{Auntie Marie:} “Good afternoon, everyone! I \textbf{arrived} \textbf{at} the house ten minutes \textbf{ago}. The traffic \textbf{on} the highway \textbf{was} terrible.”

\eng{Mother:} “Welcome, my sister! You \textbf{are} just \textbf{in} time for lunch. The children \textbf{are} \textbf{in} the garden \textbf{playing} football.”

\eng{Father:} “Marie, you \textbf{came} \textbf{by} bus, \textbf{didn't} you? \textbf{How} long \textbf{did} it \textbf{take} you?”

\eng{Auntie Marie:} “Yes, I \textbf{took} the bus \textbf{from} Douala. It \textbf{took} me four hours. I \textbf{got} \textbf{off} \textbf{at} the Total station \textbf{in} Mvan, then \textbf{walked} \textbf{to} the house.”

\eng{Cousin Paul:} “Auntie, \textbf{look at} your bag! It \textbf{is} \textbf{on} the table \textbf{next to} the fridge.”

\eng{Auntie Marie:} “Oh, thank you Paul. I \textbf{depend} \textbf{on} you young ones to \textbf{look after} me!”
\end{textebox}

\begin{experience}[Repérage guidé : Chasse aux prépositions]
\begin{enumerate}
\item En binôme, relisez le dialogue et soulignez \textbf{toutes} les prépositions (mots courts comme \eng{at, in, on, by, to, from, off, next to}).
\item Classez-les en trois colonnes sur votre cahier : \textbf{Lieu / Position}, \textbf{Temps / Durée}, \textbf{Mouvement / Direction / Manière}.
\item Comparez avec la classe. Le professeur complète au tableau.
\end{enumerate}
\end{experience}

\begin{aretenir}[Ce que nous allons maîtriser]
À la fin de cette leçon, vous saurez :
\begin{itemize}
\item Choisir la bonne préposition de \textbf{lieu} (\eng{at}, \eng{in}, \eng{on}) et les locutions de lieu (\eng{next to}, \eng{in front of}, \eng{behind}…).
\item Utiliser correctement les prépositions de \textbf{temps} (\eng{at}, \eng{on}, \eng{in}, \eng{for}, \eng{since}, \eng{during}…).
\item Distinguer \textbf{mouvement} (\eng{to}, \eng{into}, \eng{towards}) et \textbf{position}, et exprimer la \textbf{manière} (\eng{by}, \eng{with}, \eng{on foot}).
\item Éviter les pièges classiques des francophones : traduction mot à mot, faux amis (\eng{attend} vs \eng{wait for}), verbes à construction différente (\eng{enter}, \eng{marry}, \eng{depend on}…).
\end{itemize}
\end{aretenir}

\coursec{Prépositions de lieu (Place)}

\courssub{Les trois piliers : \eng{at}, \eng{in}, \eng{on}}

\begin{propriete}[Règle de base : \eng{at} / \eng{in} / \eng{on} pour la position]
\begin{itemize}
\item \textbf{\eng{at}} : point précis, adresse, événement, lieu de rendez-vous. \eng{at the bus stop}, \eng{at 15 Avenue Kennedy}, \eng{at the party}, \eng{at school}, \eng{at home}.
\item \textbf{\eng{in}} : espace clos, volume, ville, pays, quartier, rue (souvent), pièce. \eng{in the classroom}, \eng{in Douala}, \eng{in Cameroon}, \eng{in the kitchen}, \eng{in Mvog-Ada}.
\item \textbf{\eng{on}} : surface, ligne, étage, rive, île, écran, page. \eng{on the table}, \eng{on the wall}, \eng{on the first floor}, \eng{on the river}, \eng{on the screen}, \eng{on page 12}.
\end{itemize}
\end{propriete}

\begin{exemplebox}[Exemples camerounais]
\begin{itemize}
\item \eng{My father works \textbf{at} the Ministry \textbf{in} Yaoundé.} (point précis / ville)
\item \eng{We live \textbf{in} Buea, \textbf{on} the slopes of Mount Cameroon.} (ville / surface-pente)
\item \eng{The bendskin stopped \textbf{at} the roundabout \textbf{in} front of the market.} (point / locution)
\item \eng{There is a stain \textbf{on} your shirt.} (surface)
\item \eng{She is \textbf{in} the kitchen cooking \textbf{ndolé}.} (pièce close)
\end{itemize}
\end{exemplebox}

\courssub{Locutions prépositionnelles de lieu (compounds)}

\begin{propriete}[Principales locutions de lieu à connaître]
\begin{center}
\begin{tabularx}{\linewidth}{|X|X|X|}
\hline
\rowcolor{popblueL} \textbf{Locution} & \textbf{Sens} & \textbf{Exemple} \\
\hline
\eng{next to} / \eng{beside} & à côté de & \eng{The pharmacy is \textbf{next to} the bakery.} \\
\hline
\eng{in front of} & devant (extérieur) & \eng{Wait \textbf{in front of} the school gate.} \\
\hline
\eng{behind} & derrière & \eng{The garden is \textbf{behind} the house.} \\
\hline
\eng{opposite} & en face de & \eng{The bank is \textbf{opposite} the post office.} \\
\hline
\eng{between} & entre (deux) & \eng{The market is \textbf{between} the church and the mosque.} \\
\hline
\eng{among} & au milieu de (plusieurs) & \eng{He stood \textbf{among} the crowd.} \\
\hline
\eng{across from} & en face de (rue) & \eng{The hotel is \textbf{across from} the station.} \\
\hline
\eng{near} / \eng{close to} & près de & \eng{We live \textbf{near} the university.} \\
\hline
\eng{far from} & loin de & \eng{My village is \textbf{far from} the city.} \\
\hline
\eng{at the corner of} & au coin de & \eng{Meet me \textbf{at the corner of} Kennedy Avenue.} \\
\hline
\end{tabularx}
\end{center}
\end{propriete}

\begin{attention}[Pièges \eng{in} vs \eng{on} vs \eng{at} pour les francophones]
\begin{itemize}
\item \textbf{Rue} : \eng{in} \emph{ou} \eng{on} selon la variété. UK : \eng{in Oxford Street} ; US : \eng{on Oxford Street}. Au Bac camerounais (anglais britannique) : privilégiez \eng{in} pour la rue.
\item \textbf{Transport} : \eng{in} a car / a taxi / a canoe ; \eng{on} a bus / a train / a plane / a bike / a bendskin / a motorbike / a ship. Règle mnémotechnique : si on peut \textbf{se tenir debout} dedans $\rightarrow$ \eng{on} ; si on est \textbf{assis dans un habitacle} $\rightarrow$ \eng{in}.
\item \textbf{École / travail} : \eng{at school}, \eng{at work}, \eng{at university} (position générale) ; \eng{in the classroom}, \eng{in the office} (dans le bâtiment/pièce).
\end{itemize}
\end{attention}

\begin{exoresolu}[Choix de la préposition de lieu]
Énoncé : Complétez avec \eng{at}, \eng{in}, \eng{on} ou une locution.
a) The children are playing \_\_\_\_ the garden. \\
b) My uncle works \_\_\_\_ the hospital \_\_\_\_ Douala. \\
c) There is a beautiful painting \_\_\_\_ the wall. \\
d) The taxi stopped \_\_\_\_ the traffic lights. \\
e) Buea is \_\_\_\_ the foot of Mount Cameroon. \\
f) The book is \_\_\_\_ the shelf \_\_\_\_ the window.
\tcblower
Solution : a) \eng{in} (espace clos) ; b) \eng{at} (lieu de travail précis), \eng{in} (ville) ; c) \eng{on} (surface verticale) ; d) \eng{at} (point précis) ; e) \eng{at} (point précis : \eng{at the foot of}) ; f) \eng{on} (surface), \eng{next to} / \eng{beside} (locution).
\end{exoresolu}

\coursec{Prépositions de temps (Time)}

\courssub{Les trois piliers temporels : \eng{at}, \eng{on}, \eng{in}}

\begin{propriete}[Règle \eng{at} / \eng{on} / \eng{in} pour le temps]
\begin{itemize}
\item \textbf{\eng{at}} : heure précise, moment ponctuel, fête (sans \eng{day}). \eng{at 7:30}, \eng{at noon}, \eng{at midnight}, \eng{at the moment}, \eng{at Christmas}, \eng{at Easter}, \eng{at night}, \eng{at the weekend} (UK).
\item \textbf{\eng{on}} : jour, date, jour férié avec \eng{day}, partie de journée précise. \eng{on Monday}, \eng{on 20 May}, \eng{on Christmas Day}, \eng{on Monday morning}, \eng{on the morning of the exam}.
\item \textbf{\eng{in}} : période longue (mois, année, saison, siècle), partie de journée (générale), durée future. \eng{in June}, \eng{in 2024}, \eng{in the dry season}, \eng{in the morning}, \eng{in the afternoon}, \eng{in the evening}, \eng{in two weeks} (dans deux semaines).
\end{itemize}
\end{propriete}

\begin{exemplebox}[Exemples ancrés dans le calendrier camerounais]
\begin{itemize}
\item \eng{The match starts \textbf{at} 3:30 p.m. \textbf{on} Sunday \textbf{in} January.}
\item \eng{We celebrate Youth Day \textbf{on} 11 February.}
\item \eng{The rainy season begins \textbf{in} March.}
\item \eng{I revise my lessons \textbf{in} the evening \textbf{at} 8 p.m.}
\item \eng{The results will be out \textbf{in} three days.}
\end{itemize}
\end{exemplebox}

\courssub{Durée, antériorité, intervalle : \eng{for}, \eng{since}, \eng{during}, \eng{from... to}, \eng{until/till}}

\begin{propriete}[Autres prépositions de temps essentielles]
\begin{itemize}
\item \textbf{\eng{for}} + durée : \eng{for two hours}, \eng{for a long time}, \eng{for ages}. (Répond à \eng{How long?})
\item \textbf{\eng{since}} + point de départ (date, heure, événement) : \eng{since Monday}, \eng{since 2010}, \eng{since I arrived}. Souvent avec \textbf{present perfect} : \eng{I have lived here \textbf{since} 2015.}
\item \textbf{\eng{during}} + nom (événement/période) : \eng{during the lesson}, \eng{during the holidays}, \eng{during the match}. \textbf{Jamais} \eng{during two hours} (utilisez \eng{for}).
\item \textbf{\eng{from ... to / until / till}} : intervalle. \eng{from Monday to Friday}, \eng{from 8 a.m. until 4 p.m.}
\item \textbf{\eng{by}} : échéance (au plus tard). \eng{Submit your assignment \textbf{by} Friday.}
\item \textbf{\eng{before}} / \textbf{\eng{after}} + nom / gérondif : \eng{before lunch}, \eng{after leaving school}.
\end{itemize}
\end{propriete}

\begin{attention}[Erreurs fréquentes : \eng{for} vs \eng{since} vs \eng{during}]
\begin{itemize}
\item \eng{I have waited \textbf{for} two hours.} (durée) $\neq$ \eng{I have waited \textbf{since} 10 a.m.} (départ).
\item \eng{We talked \textbf{during} the break.} (événement) $\neq$ \eng{We talked \textbf{for} twenty minutes.} (durée).
\item \textbf{Jamais} \eng{since two years} $\rightarrow$ \eng{for two years}.
\item \textbf{Jamais} \eng{during two hours} $\rightarrow$ \eng{for two hours}.
\end{itemize}
\end{attention}

\begin{exoresolu}[Prépositions de temps : correction]
Énoncé : Corrigez si nécessaire.
a) I have known him \textbf{during} ten years. \\
b) The shop opens \textbf{at} Monday \textbf{at} 8 a.m. \\
c) We will finish the work \textbf{until} Friday. \\
d) She has been here \textbf{for} January.
\tcblower
Solution : a) \eng{for} ten years (durée) ; b) \eng{on} Monday (jour), \eng{at} 8 a.m. (heure) ; c) \eng{by} Friday (échéance) ou \eng{until} Friday (continuité jusqu'à) ; d) \eng{since} January (point de départ).
\end{exoresolu}

\coursec{Direction, mouvement et manière}

\courssub{Mouvement vers un lieu : \eng{to}, \eng{into}, \eng{towards}, \eng{through}, \eng{across}}

\begin{propriete}[Choix selon la nature du mouvement]
\begin{itemize}
\item \textbf{\eng{to}} : direction générale, destination. \eng{go to school}, \eng{walk to the market}, \eng{travel to Douala}. \textbf{Standard} pour verbes de mouvement (\eng{go, come, walk, run, drive, fly, travel, send}).
\item \textbf{\eng{into}} : mouvement \textbf{vers l'intérieur} d'un volume (insiste sur l'entrée). \eng{walk into the classroom}, \eng{get into the car}, \eng{fall into the river}. \textit{Contraste :} \eng{He is in the room} (position) vs \eng{He came into the room} (mouvement d'entrée).
\item \textbf{\eng{towards}} : direction \textbf{vers} (mais pas nécessairement arrivée). \eng{He walked towards the door.}
\item \textbf{\eng{through}} : traverser un volume (forêt, foule, tunnel). \eng{drive through the tunnel}, \eng{walk through the forest}.
\item \textbf{\eng{across}} : traverser une surface (rue, place, pont, frontière). \eng{walk across the road}, \eng{swim across the river}.
\end{itemize}
\end{propriete}

\begin{exemplebox}[Mouvement dans un trajet quotidien]
\eng{Every morning, Pierre leaves his house \textbf{in} Mvan \textbf{at} 6:30. He walks \textbf{to} the bus stop. He gets \textbf{into} a bendskin. They drive \textbf{through} the traffic jam \textbf{on} the bridge. He gets \textbf{off} \textbf{at} the school gate and walks \textbf{into} the building.}
\end{exemplebox}

\courssub{Manière, moyen, instrument : \eng{by}, \eng{with}, \eng{on}, \eng{in}}

\begin{propriete}[Prépositions de manière et de moyen]
\begin{itemize}
\item \textbf{\eng{by}} + moyen de transport (sans article) : \eng{by bus}, \eng{by car}, \eng{by train}, \eng{by plane}, \eng{by bike}, \eng{by bendskin}, \eng{by foot} (mais \eng{on foot} plus courant).
\item \textbf{\eng{on}} + \eng{foot} : \eng{on foot} = à pied.
\item \textbf{\eng{in}} / \textbf{\eng{on}} + article + véhicule : \eng{in a car}, \eng{in a taxi}, \eng{on a bus}, \eng{on a motorbike}, \eng{on a bendskin}. (Rappel : \eng{in} habitacle, \eng{on} surface/straddle).
\item \textbf{\eng{with}} + instrument / outil / partie du corps : \eng{write with a pen}, \eng{cut with a knife}, \eng{see with my eyes}.
\item \textbf{\eng{by}} + gérondif (action) : \eng{by working hard}, \eng{by taking the bus}.
\item \textbf{\eng{like}} / \textbf{\eng{as}} : \eng{He runs \textbf{like} the wind}, \eng{He works \textbf{as} a teacher}.
\end{itemize}
\end{propriete}

\begin{attention}[Piège \eng{by} vs \eng{with} pour l'agent (passif) et l'instrument]
\begin{itemize}
\item \textbf{Agent (passif)} : \eng{by} + personne. \eng{The book was written \textbf{by} Mongo Beti.}
\item \textbf{Instrument} : \eng{with} + objet. \eng{The door was opened \textbf{with} a key.}
\item \textbf{Moyen de transport} : \eng{by} (règle générale sans article) mais \eng{in/on} + \emph{article} si précision.
\end{itemize}
\end{attention}

\begin{exoresolu}[Direction et manière]
Énoncé : Choisissez la bonne préposition.
a) They swam \_\_\_\_ the river to the other side. (\eng{across} / \eng{through}) \\
b) We went \_\_\_\_ the village \_\_\_\_ foot. (\eng{to} / \eng{towards}) (\eng{by} / \eng{on}) \\
c) The thief entered \_\_\_\_ the house \_\_\_\_ the window. (\eng{into} / \eng{in}) (\eng{through} / \eng{across}) \\
d) She sent the letter \_\_\_\_ post. (\eng{by} / \eng{with})
\tcblower
Solution : a) \eng{across} (surface d'eau) ; b) \eng{to} (destination atteinte), \eng{on foot} (locution figée) ; c) \eng{into} (mouvement vers intérieur), \eng{through} (traverser volume/ouverture) ; d) \eng{by post} (moyen).
\end{exoresolu}

\coursec{Comparaison français/anglais et pièges des francophones}

Dans les sections précédentes, nous avons vu que le choix entre \eng{to}, \eng{into}, \eng{at} et \eng{in} dépend d'une question simple : y a-t-il \cle{mouvement} ou \cle{position} ? Mais il reste un piège beaucoup plus redoutable pour les francophones : croire que les prépositions anglaises se traduisent mot à mot depuis le français. C'est faux dans la grande majorité des cas. Cette section est sans doute la plus rentable de toute la leçon pour votre moyenne au Bac.

\courssub{Le principe général : ne jamais traduire mot à mot}

\begin{propriete}[La règle d'or des prépositions]
Une préposition ne se traduit presque \cle{jamais} mot à mot entre le français et l'anglais. Le français dit « dépendre \emph{de} », l'anglais dit \eng{depend \textbf{on}} ; le français dit « entrer \emph{dans} », l'anglais dit \eng{enter} tout court. Il faut donc apprendre chaque verbe \textbf{avec} sa préposition (ou son absence de préposition) comme un seul bloc indivisible : \eng{listen to}, \eng{wait for}, \eng{depend on}, \eng{marry somebody}.
\end{propriete}

Observez d'abord ces phrases tirées de la vie scolaire camerounaise. Repérez la préposition anglaise et comparez avec la traduction française :

\begin{exemplebox}[Observation : la préposition anglaise surprend toujours]
\begin{itemize}
\item \eng{The students are listening to the headmaster.} = Les élèves écoutent le proviseur. (français : zéro préposition ; anglais : \eng{to})
\item \eng{We are waiting for the bendskin.} = Nous attendons la moto-taxi. (français : zéro préposition ; anglais : \eng{for})
\item \eng{She entered the classroom silently.} = Elle est entrée \emph{dans} la salle en silence. (français : « dans » ; anglais : zéro préposition)
\item \eng{My cousin married a nurse from Buea.} = Mon cousin a épousé une infirmière de Buea. (français : « épouser quelqu'un », mais beaucoup de francophones traduisent mentalement « se marier \emph{avec} » ; anglais : zéro préposition)
\item \eng{Success depends on hard work.} = La réussite dépend \emph{du} travail. (français : « de » ; anglais : \eng{on})
\end{itemize}
Conclusion : aucune correspondance fiable. Il faut mémoriser les blocs.
\end{exemplebox}

\begin{methode}[Comment mémoriser les constructions verbe + préposition]
\begin{enumerate}
\item N'apprenez jamais un verbe seul. Apprenez le \textbf{bloc} complet : \eng{listen to}, \eng{wait for}, \eng{depend on}, \eng{apologize for}.
\item Fabriquez une \textbf{phrase personnelle} avec chaque bloc, liée à votre vie : \eng{I listen to makossa on the radio.} / \eng{I am waiting for my results.}
\item Relisez vos phrases à voix haute chaque semaine : l'oreille finit par « entendre » la bonne préposition.
\item Dans vos rédactions, relisez une dernière fois en ne traquant \textbf{que} les prépositions.
\end{enumerate}
\end{methode}

\courssub{Anglais avec préposition, français sans : \eng{listen to}, \eng{look at}, \eng{wait for}, \eng{ask for}, \eng{pay for}}

Ces verbes sont transitifs directs en français (« écouter quelque chose », « attendre quelqu'un »), mais exigent une préposition en anglais.

\begin{propriete}[Verbes anglais à préposition obligatoire]
\begin{itemize}
\item \eng{to listen \textbf{to}} = écouter quelque chose/quelqu'un
\item \eng{to look \textbf{at}} = regarder
\item \eng{to wait \textbf{for}} = attendre
\item \eng{to ask \textbf{for}} = demander (quelque chose)
\item \eng{to pay \textbf{for}} = payer (quelque chose)
\end{itemize}
La préposition reste obligatoire même si le complément est un pronom : \eng{Listen to me!} = Écoute-moi !
\end{propriete}

\begin{exemplebox}[Exemples traduits]
\begin{itemize}
\item \eng{I listen to the radio every morning.} = J'écoute la radio tous les matins.
\item \eng{Look at the blackboard, please.} = Regardez le tableau, s'il vous plaît.
\item \eng{They are waiting for the bus in front of the school.} = Ils attendent le bus devant l'école.
\item \eng{She asked for a glass of water.} = Elle a demandé un verre d'eau.
\item \eng{My father paid for my school fees yesterday.} = Mon père a payé mes frais de scolarité hier.
\end{itemize}
\end{exemplebox}

\begin{attention}[Erreur fréquente des francophones]
\eng{I listen the radio.} est \textbf{FAUX} : il faut \eng{I listen \textbf{to} the radio.} De même, \eng{I am waiting the bus.} est \textbf{FAUX} : il faut \eng{I am waiting \textbf{for} the bus.} Et on dit \eng{He paid for the books.}, jamais \eng{He paid the books} quand on parle de l'objet acheté (on peut dire \eng{He paid the driver} = il a payé le chauffeur, la personne).
\end{attention}

\courssub{Anglais sans préposition, français avec : \eng{enter}, \eng{marry}, \eng{phone}}

C'est le piège symétrique : le français met une préposition (« entrer \emph{dans} », « se marier \emph{avec} », « téléphoner \emph{à} »), l'anglais n'en met \textbf{aucune}.

\begin{propriete}[Verbes anglais transitifs directs]
\begin{itemize}
\item \eng{to enter a room} = entrer \emph{dans} une pièce (pas de \eng{into} après \eng{enter} : le verbe contient déjà l'idée de mouvement vers l'intérieur)
\item \eng{to marry somebody} = épouser quelqu'un (pas de \eng{with})
\item \eng{to phone somebody} = téléphoner \emph{à} quelqu'un (pas de \eng{to})
\end{itemize}
Attention : le nom \eng{marriage} et l'adjectif \eng{married} prennent eux la préposition \eng{to} : \eng{She is married \textbf{to} a teacher.} = Elle est mariée à un enseignant.
\end{propriete}

\begin{attention}[Erreur fréquente des francophones]
\eng{He entered into the room.} est lourd et fautif : dites \eng{He entered the room.} (\eng{enter into} existe mais seulement dans un sens abstrait : \eng{enter into an agreement} = conclure un accord). \eng{She married with him.} est \textbf{FAUX} : dites \eng{She married him.} ou \eng{She got married to him.} Enfin, \eng{I phoned to my mother.} est faux : \eng{I phoned my mother.}
\end{attention}

\courssub{Français avec préposition, anglais avec une préposition différente : \eng{depend on}, \eng{consist in}, \eng{succeed in}, \eng{insist on}, \eng{apologize for}}

Ici, les deux langues mettent une préposition, mais \textbf{pas la même}. C'est le cas le plus délicat, car la présence d'une préposition en français donne une fausse confiance.

\begin{propriete}[Constructions à mémoriser comme des blocs]
\begin{itemize}
\item \eng{to depend \textbf{on}} = dépendre \emph{de}
\item \eng{to consist \textbf{in}} = consister \emph{en} / à
\item \eng{to succeed \textbf{in} (+ -ing)} = réussir \emph{à} (faire)
\item \eng{to insist \textbf{on} (+ -ing)} = insister \emph{pour} (faire)
\item \eng{to apologize \textbf{for} (+ -ing)} = s'excuser \emph{de} (faire)
\end{itemize}
Après \eng{succeed in}, \eng{insist on} et \eng{apologize for}, le verbe suivant est au gérondif (\emph{-ing}), jamais à l'infinitif.
\end{propriete}

\begin{exemplebox}[Exemples traduits]
\begin{itemize}
\item \eng{It depends on the weather.} = Cela dépend du temps (qu'il fait).
\item \eng{He succeeded in passing the exam.} = Il a réussi à passer l'examen.
\item \eng{Happiness consists in helping others.} = Le bonheur consiste à aider les autres.
\item \eng{She insisted on paying for the meal.} = Elle a insisté pour payer le repas.
\item \eng{He apologized for being late.} = Il s'est excusé d'être en retard.
\end{itemize}
\end{exemplebox}

\courssub{Faux amis de construction : \eng{attend}, \eng{assist}, et les autres}

\begin{attention}[« Attend » n'est pas « attendre »]
\eng{to attend} signifie \textbf{fréquenter} ou \textbf{assister à} (être présent) : \eng{She attends Government High School Bamenda.} = Elle fréquente le lycée de Bamenda. \eng{We attended a meeting.} = Nous avons assisté à une réunion. « Attendre » se dit \eng{to wait for}. Ne dites jamais \eng{I am attending the bus} pour « j'attends le bus » !
\end{attention}

\begin{exemplebox}[Autres faux amis de construction à connaître]
\begin{itemize}
\item \eng{to assist} = aider (soutenu) ; « assister à » un événement = \eng{to attend}.
\item \eng{to join somebody} = rejoindre quelqu'un (pas de \eng{with}) : \eng{May I join you?} = Puis-je me joindre à vous ?
\item \eng{to answer a question} = répondre \emph{à} une question (pas de \eng{to}).
\item \eng{to resemble somebody} = ressembler \emph{à} quelqu'un (pas de \eng{to}) : \eng{She resembles her mother.}
\item \eng{to reach a place} = arriver \emph{à} un lieu (pas de \eng{to} ni \eng{at}) : \eng{We reached Yaoundé at noon.}
\item \eng{to approach a place} = s'approcher de (pas de \eng{to}) : \eng{The plane approached the airport.}
\end{itemize}
\end{exemplebox}

\begin{popfigure}
\begin{center}
\begin{tikzpicture}[font=\small]
\node[draw=popblue, fill=popblueL, rounded corners, align=center, minimum width=5.2cm] (f1) at (0,3) {\textbf{Le français dit...}\\ « écouter la radio »};
\node[draw=popgreen, fill=popgreenL, rounded corners, align=center, minimum width=5.2cm] (e1) at (7.5,3) {\textbf{L'anglais dit...}\\ \eng{listen \textbf{to} the radio}};
\draw[-{Stealth}, very thick, poporange] (f1) -- (e1);
\node[draw=popblue, fill=popblueL, rounded corners, align=center, minimum width=5.2cm] (f2) at (0,1.5) {« entrer \emph{dans} la salle »};
\node[draw=popgreen, fill=popgreenL, rounded corners, align=center, minimum width=5.2cm] (e2) at (7.5,1.5) {\eng{enter the room} (rien !)};
\draw[-{Stealth}, very thick, poporange] (f2) -- (e2);
\node[draw=popblue, fill=popblueL, rounded corners, align=center, minimum width=5.2cm] (f3) at (0,0) {« se marier \emph{avec} lui »};
\node[draw=popgreen, fill=popgreenL, rounded corners, align=center, minimum width=5.2cm] (e3) at (7.5,0) {\eng{marry him} (rien !)};
\draw[-{Stealth}, very thick, poporange] (f3) -- (e3);
\node[draw=popblue, fill=popblueL, rounded corners, align=center, minimum width=5.2cm] (f4) at (0,-1.5) {« dépendre \emph{du} temps »};
\node[draw=popgreen, fill=popgreenL, rounded corners, align=center, minimum width=5.2cm] (e4) at (7.5,-1.5) {\eng{depend \textbf{on} the weather}};
\draw[-{Stealth}, very thick, poporange] (f4) -- (e4);
\end{tikzpicture}
\end{center}
\legende{Correspondances trompeuses : la préposition française ne prédit jamais la préposition anglaise.}
\end{popfigure}

\begin{popfigure}
\begin{center}
\begin{tikzpicture}[font=\small, node distance=0.9cm]
\node[draw=popdark, fill=popgoldL, rounded corners, align=center] (q) {Le verbe exprime...};
\node[draw=poporange, fill=poporangeL, rounded corners, align=center, below left=0.9cm and 0.2cm of q] (mov) {\textbf{Mouvement ?}\\ \eng{to}, \eng{into}, \eng{towards}\\ \eng{go to school}, \eng{enter the room}};
\node[draw=popblue, fill=popblueL, rounded corners, align=center, below right=0.9cm and 0.2cm of q] (pos) {\textbf{Position ?}\\ \eng{at}, \eng{in}, \eng{on}\\ \eng{live in Douala}, \eng{stay at home}};
\draw[-{Stealth}, thick, popdark] (q) -- (mov);
\draw[-{Stealth}, thick, popdark] (q) -- (pos);
\end{tikzpicture}
\end{center}
\legende{Organigramme de décision : mouvement (\eng{to}/\eng{into}) ou position (\eng{at}/\eng{in}) ?}
\end{popfigure}

\courssub{Les 10 erreurs les plus fréquentes au Bac}

\begin{center}
\begin{tabularx}{\linewidth}{|X|X|X|}
\hline
\rowcolor{popblueL} \textbf{Erreur fréquente (FAUX)} & \textbf{Correction} & \textbf{Français} \\
\hline
\eng{I listen the radio.} & \eng{I listen \textbf{to} the radio.} & J'écoute la radio. \\
\hline
\eng{She is waiting the bus.} & \eng{She is waiting \textbf{for} the bus.} & Elle attend le bus. \\
\hline
\eng{He entered into the room.} & \eng{He entered the room.} & Il est entré dans la pièce. \\
\hline
\eng{She married with him.} & \eng{She married him.} & Elle l'a épousé. \\
\hline
\eng{I phoned to my uncle.} & \eng{I phoned my uncle.} & J'ai téléphoné à mon oncle. \\
\hline
\eng{It depends of the price.} & \eng{It depends \textbf{on} the price.} & Cela dépend du prix. \\
\hline
\eng{He succeeded to pass.} & \eng{He succeeded \textbf{in passing}.} & Il a réussi à passer. \\
\hline
\eng{She insisted to come.} & \eng{She insisted \textbf{on coming}.} & Elle a insisté pour venir. \\
\hline
\eng{I am attending the bus.} & \eng{I am waiting \textbf{for} the bus.} & J'attends le bus. \\
\hline
\eng{He apologized to be late.} & \eng{He apologized \textbf{for being} late.} & Il s'est excusé d'être en retard. \\
\hline
\end{tabularx}
\end{center}

\begin{exoresolu}[Correction de phrases]
Énoncé : corrigez les erreurs de prépositions dans les phrases suivantes. a) \eng{My sister is waiting the taxi.} b) \eng{They entered into the church.} c) \eng{It depends of you.} d) \eng{He succeeded to win the match.}
\tcblower
Solution détaillée : a) \eng{wait} exige \eng{for} : \eng{My sister is waiting \textbf{for} the taxi.} b) \eng{enter} est transitif direct : \eng{They entered the church.} c) « dépendre de » se dit \eng{depend on} : \eng{It depends \textbf{on} you.} d) \eng{succeed} se construit avec \eng{in} + gérondif : \eng{He succeeded \textbf{in winning} the match.}
\end{exoresolu}

\coursec{Entraînement guidé et production}

\courssub{Exercices d'application}

\begin{exercice}[Complétion : prépositions de lieu et de temps]
Complétez avec la préposition ou locution appropriée (\eng{at, in, on, next to, behind, opposite, between, for, since, during, from, to, until, by}).
\begin{enumerate}
\item The final exam starts \_\_\_\_ 8 a.m. \_\_\_\_ Monday \_\_\_\_ June.
\item My house is \_\_\_\_ the church and the mosque, \_\_\_\_ the market.
\item I have been learning English \_\_\_\_ 2018. I study \_\_\_\_ two hours every evening.
\item \_\_\_\_ the holidays, we travelled \_\_\_\_ Kribi \_\_\_\_ car.
\item Please submit your project \_\_\_\_ Friday \_\_\_\_ the latest.
\end{enumerate}
\end{exercice}

\begin{exercice}[Transformation : mouvement vs position]
Réécrivez les phrases en changeant la position en mouvement (ou l'inverse) et adaptez la préposition.
\begin{enumerate}
\item \eng{The children are in the garden.} (mouvement : \eng{run})
\item \eng{He walked into the office.} (position)
\item \eng{The bus stops at the station.} (mouvement vers : \eng{drive})
\item \eng{She is on the bike.} (mouvement : \eng{get})
\end{enumerate}
\end{exercice}

\begin{exercice}[Verbes à construction piège : choix multiple]
Choisissez la bonne forme.
\begin{enumerate}
\item I need to \_\_\_\_ the director. (speak / speak to / speak with)
\item The cake consists \_\_\_\_ flour, eggs and sugar. (of / in / from)
\item He apologized \_\_\_\_ the mistake. (for / of / to)
\item We must \_\_\_\_ the problem. (discuss / discuss about / discuss on)
\item She resembles \_\_\_\_ her grandmother. (to / -- / with)
\end{enumerate}
\end{exercice}

\courssub{Production écrite guidée}

\begin{experience}[Rédaction : « Mon trajet pour l'école »]
Rédigez un paragraphe de 80-100 mots en anglais pour décrire votre trajet quotidien de la maison à l'école. Utilisez \textbf{au moins} :
\begin{itemize}
\item 3 prépositions de lieu (\eng{at, in, on, next to, behind}…)
\item 3 prépositions de temps (\eng{at, on, in, for, since, during}…)
\item 2 prépositions de mouvement (\eng{to, into, towards, through, across}…)
\item 1 préposition de manière (\eng{by, on foot, with}…)
\item 2 blocs verbe+préposition vus en section 5 (\eng{wait for, listen to, depend on, succeed in}…)
\end{itemize}
\textbf{Modèle de début :} \eng{Every morning, I wake up at... I leave my house in... at... I walk to...}
\end{experience}

\begin{experience}[Expression orale : « Guide touristique de mon quartier »]
En binôme : l'un est un touriste qui arrive à la gare routière de votre ville, l'autre est un guide. Le touriste demande comment aller à trois endroits (marché, lycée, hôpital, cybercafé...). Le guide explique le chemin en utilisant des prépositions de lieu, de direction et des locutions (\eng{go straight on, turn left at, cross, next to, opposite, between}). Échangez les rôles. Le professeur écoute et note les prépositions correctes/incorrectes.
\end{experience}

\begin{corrige}[Exercices d'application]
\textbf{Exercice 1 (Complétion) :}
1. \eng{at} 8 a.m. \eng{on} Monday \eng{in} June.
2. \eng{between} the church and the mosque, \eng{opposite} / \eng{near} the market.
3. \eng{since} 2018. \eng{for} two hours.
4. \eng{During} the holidays, we travelled \eng{to} Kribi \eng{by} car.
5. \eng{by} Friday (échéance).

\textbf{Exercice 2 (Transformation) :}
1. \eng{The children run into the garden.}
2. \eng{He is in the office.}
3. \eng{The bus drives to the station.}
4. \eng{She gets on the bike.}

\textbf{Exercice 3 (Choix multiple) :}
1. \eng{speak to} (parler à quelqu'un) — \eng{speak with} possible mais \eng{to} plus courant pour initier.
2. \eng{consists in} (UK) / \eng{consists of} (US) — au Bac camerounais (UK) : \eng{in}.
3. \eng{apologized for}.
4. \eng{discuss} (transitif direct, pas de préposition).
5. \eng{resembles} (transitif direct, pas de préposition).
\end{corrige}

\begin{savaistu}[Le marqueur n°1 des correcteurs]
À l'oral comme à l'écrit du Bac, les correcteurs repèrent d'abord les erreurs de prépositions : c'est le marqueur numéro 1 du niveau d'un candidat francophone. Un texte avec \eng{depend on}, \eng{wait for} et \eng{succeed in} correctement employés donne immédiatement une impression de maîtrise. En anglais parlé, les locuteurs natifs eux-mêmes hésitent rarement sur ces blocs parce qu'ils les ont appris entiers, jamais séparément : faites exactement pareil.
\end{savaistu}

\begin{savaistu}[Contexte camerounais : le bendskin et les prépositions]
Au Cameroun, le \eng{bendskin} (moto-taxi) est roi. On dit : \eng{I go to school \textbf{by} bendskin} (moyen général) mais \eng{I get \textbf{on} a bendskin} (action de monter sur la moto). On ne dit \textbf{jamais} \eng{in a bendskin} (sauf si vous êtes dans un side-car fermé !). De même : \eng{wait for a bendskin}, \eng{get off a bendskin}, \eng{pay the bendskin driver} (pas \eng{pay for the driver}).
\end{savaistu}

\coursec{Entraînement guidé et production}

Après avoir comparé le français et l'anglais et repéré les pièges les plus fréquents des francophones (\eng{married with}, \eng{listen the radio}, \eng{in the morning} confondu avec \eng{on Monday morning}...), il est temps de passer à la pratique. C'est en s'entraînant qu'une règle devient un réflexe. Cette section vous propose un parcours progressif en cinq étapes : identifier, compléter, corriger, transformer, puis produire. C'est exactement le cheminement que l'on attend de vous au baccalauréat.

\begin{popfigure}
\begin{center}
\begin{tikzpicture}[node distance=0.35cm]
\node[draw=popblue, fill=popblueL, rounded corners, minimum width=2.4cm, minimum height=0.9cm, align=center, font=\small\bfseries] (a) {1. Identifier};
\node[draw=popgreen, fill=popgreenL, rounded corners, minimum width=2.4cm, minimum height=0.9cm, align=center, font=\small\bfseries, right=of a] (b) {2. Compléter};
\node[draw=poporange, fill=poporangeL, rounded corners, minimum width=2.4cm, minimum height=0.9cm, align=center, font=\small\bfseries, right=of b] (c) {3. Corriger};
\node[draw=poppurple, fill=poppurpleL, rounded corners, minimum width=2.4cm, minimum height=0.9cm, align=center, font=\small\bfseries, right=of c] (d) {4. Transformer};
\node[draw=poppink, fill=poppinkL, rounded corners, minimum width=2.4cm, minimum height=0.9cm, align=center, font=\small\bfseries, right=of d] (e) {5. Produire};
\draw[-{Stealth}, thick, popdark] (a) -- (b);
\draw[-{Stealth}, thick, popdark] (b) -- (c);
\draw[-{Stealth}, thick, popdark] (c) -- (d);
\draw[-{Stealth}, thick, popdark] (d) -- (e);
\end{tikzpicture}
\end{center}
\legende{La frise de progression : de l'identification à la production libre.}
\end{popfigure}

\courssub{La méthode pour choisir la bonne préposition}

Avant de vous lancer dans les exercices, retenez cette démarche en trois questions. Elle fonctionne pour tous les types d'exercices : phrases à trous, QCM, correction d'erreurs.

\begin{methode}[Trois questions pour choisir la bonne préposition]
\begin{enumerate}
\item \textbf{Identifier la fonction.} Le complément exprime-t-il un \cle{lieu} (où ?), un \cle{temps} (quand ?), un \cle{mouvement} (vers où ? d'où ?) ou une \cle{manière} (comment ? avec quoi ?).
\item \textbf{Appliquer la triade \eng{in / on / at}.} Pour le lieu : \eng{in} (espace fermé, ville, pays), \eng{on} (surface, étage, rue), \eng{at} (point précis). Pour le temps : \eng{in} (mois, années, saisons, parties de la journée), \eng{on} (jours et dates), \eng{at} (heures et moments ponctuels).
\item \textbf{Vérifier la construction figée.} S'agit-il d'un verbe à préposition obligatoire (\eng{listen to}, \eng{wait for}, \eng{get married to}) ou d'une expression toute faite (\eng{on foot}, \eng{by bus}, \eng{at home}, \eng{in time}) ? Dans ce cas, il faut la connaître par cœur : aucune logique ne peut la deviner.
\end{enumerate}
\end{methode}

\begin{propriete}[Rappel de synthèse : les grandes familles de prépositions]
\begin{itemize}
\item \textbf{Position} : \eng{at} (point précis : \eng{at the bus stop}), \eng{in} (espace fermé : \eng{in the kitchen}, \eng{in Douala}), \eng{on} (surface : \eng{on the table}, \eng{on the second floor}).
\item \textbf{Mouvement} : \eng{to} (vers : \eng{go to school}), \eng{into} (entrée : \eng{come into the room}), \eng{out of} (sortie : \eng{get out of the taxi}).
\item \textbf{Durée} : \eng{for} + durée (\eng{for two years}), \eng{since} + point de départ (\eng{since 2020}).
\item \textbf{Moyen de transport} : \eng{by} + moyen sans article (\eng{by bus}, \eng{by taxi}) ; exception : \eng{on foot}.
\item \textbf{Instrument / accompagnement} : \eng{with} (\eng{cut it with a knife}, \eng{live with my grandmother}).
\end{itemize}
\end{propriete}

\begin{exemplebox}[Un modèle à imiter]
\eng{My father works in a bank, on the second floor. He leaves home at 7 a.m. and goes to work by taxi.}

« Mon père travaille dans une banque, au deuxième étage. Il quitte la maison à 7 h et va au travail en taxi. »

Observez : cinq prépositions, cinq fonctions différentes — lieu fermé (\eng{in a bank}), surface/étage (\eng{on the second floor}), heure précise (\eng{at 7 a.m.}), mouvement vers (\eng{to work}), moyen de transport (\eng{by taxi}). C'est cette précision que l'on attend dans vos productions.
\end{exemplebox}

\courssub{Exercice 1 : phrases à trous (fonctions mélangées)}

\begin{exercice}[Complete with the correct preposition]
Complétez chaque phrase avec la préposition qui convient : \eng{at -- in -- on -- to -- into -- out of -- for -- since -- by -- with}.

\begin{enumerate}
\item My sister has lived \ldots\ Bamenda \ldots\ 2019.
\item The children are playing \ldots\ the garden, \ldots\ front of the house.
\item We go \ldots\ church \ldots\ Sundays \ldots\ 8 o'clock.
\item My uncle came \ldots\ of his car and walked \ldots\ the shop.
\item She has been waiting \ldots\ the bus \ldots\ thirty minutes.
\item My grandmother cuts vegetables \ldots\ a very sharp knife.
\item They travelled \ldots\ Yaoundé \ldots\ train last December.
\item There is a beautiful picture \ldots\ the wall \ldots\ the sitting room.
\end{enumerate}
\end{exercice}

\begin{corrige}[Exercice 1]
\begin{enumerate}
\item \eng{in Bamenda, since 2019} — ville : \eng{in} ; point de départ d'une durée : \eng{since}.
\item \eng{in the garden, in front of the house} — espace délimité : \eng{in} ; locution figée : \eng{in front of}.
\item \eng{to church on Sundays at 8 o'clock} — mouvement : \eng{to} ; jour : \eng{on} ; heure : \eng{at}.
\item \eng{out of his car, into the shop} — sortie : \eng{out of} ; entrée : \eng{into}.
\item \eng{waiting for the bus for thirty minutes} — verbe à préposition : \eng{wait for} ; durée : \eng{for}.
\item \eng{with a very sharp knife} — instrument : \eng{with}.
\item \eng{to Yaoundé by train} — destination : \eng{to} ; moyen : \eng{by}.
\item \eng{on the wall, in the sitting room} — surface : \eng{on} ; pièce : \eng{in}.
\end{enumerate}
\end{corrige}

\courssub{Exercice 2 : QCM de type baccalauréat}

\begin{exercice}[Choose the correct option]
Entourez la bonne réponse (A, B, C ou D).

\begin{enumerate}
\item My parents got married \ldots\ 1995.\\
A. on \hspace{0.8cm} B. in \hspace{0.8cm} C. at \hspace{0.8cm} D. since
\item The meeting starts \ldots\ 10 a.m. \ldots\ Monday.\\
A. at / on \hspace{0.8cm} B. in / at \hspace{0.8cm} C. on / in \hspace{0.8cm} D. at / in
\item She goes to school \ldots\ foot every morning.\\
A. by \hspace{0.8cm} B. with \hspace{0.8cm} C. on \hspace{0.8cm} D. in
\item The cat jumped \ldots\ the box and hid inside.\\
A. onto \hspace{0.8cm} B. into \hspace{0.8cm} C. out of \hspace{0.8cm} D. at
\item We have known our neighbours \ldots\ ten years.\\
A. since \hspace{0.8cm} B. during \hspace{0.8cm} C. for \hspace{0.8cm} D. at
\item My brother is interested \ldots\ football.\\
A. for \hspace{0.8cm} B. on \hspace{0.8cm} C. at \hspace{0.8cm} D. in
\end{enumerate}
\end{exercice}

\begin{corrige}[Exercice 2]
\begin{enumerate}
\item \textbf{B} — \eng{in 1995} : année $\rightarrow$ \eng{in}.
\item \textbf{A} — \eng{at 10 a.m. on Monday} : heure $\rightarrow$ \eng{at}, jour $\rightarrow$ \eng{on}.
\item \textbf{C} — \eng{on foot} : expression figée, exception à la règle de \eng{by}.
\item \textbf{B} — \eng{jumped into the box} : mouvement d'entrée $\rightarrow$ \eng{into} (et non \eng{in}, qui indique la position).
\item \textbf{C} — \eng{for ten years} : durée $\rightarrow$ \eng{for} (\eng{since} exige un point de départ : \eng{since 2015}).
\item \textbf{D} — \eng{interested in} : adjectif à préposition obligatoire.
\end{enumerate}
\end{corrige}

\courssub{Exercice 3 : corriger les erreurs typiques des francophones}

\begin{exoresolu}[Correct the mistakes]
Énoncé. Chaque phrase contient une erreur de préposition. Soulignez-la et corrigez-la.

\begin{enumerate}
\item My cousin is married with a nurse from Buea.
\item We listen the radio every evening.
\item I go at school by bendskin.
\item She arrived to Douala last week.
\item He has lived here since five years.
\end{enumerate}

\tcblower
Solution détaillée.

\begin{enumerate}
\item \eng{married with} $\rightarrow$ \eng{married \textbf{to}}. Faux ami : le français dit « marié avec », l'anglais dit \eng{married to}. \eng{My cousin is married to a nurse from Buea.}
\item \eng{listen the radio} $\rightarrow$ \eng{listen \textbf{to} the radio}. Le verbe \eng{listen} exige toujours \eng{to} devant son complément.
\item \eng{go at school} $\rightarrow$ \eng{go \textbf{to} school}. Le mouvement vers un lieu exige \eng{to}, jamais \eng{at} (qui marque la position).
\item \eng{arrived to Douala} $\rightarrow$ \eng{arrived \textbf{in} Douala}. Le verbe \eng{arrive} se construit avec \eng{in} (ville, pays) ou \eng{at} (point précis : \eng{arrive at the station}), jamais avec \eng{to}.
\item \eng{since five years} $\rightarrow$ \eng{\textbf{for} five years}. \eng{Since} introduit un point de départ (\eng{since 2020}) ; une durée s'exprime avec \eng{for}.
\end{enumerate}
\end{exoresolu}

\begin{exercice}[À vous de jouer]
Corrigez de la même façon :
\begin{enumerate}
\item My mother depends of my father's salary.
\item They discussed about the family budget.
\item I am good in Mathematics.
\item We entered into the room quietly.
\item The market is full of people in Saturdays.
\end{enumerate}
\end{exercice}

\begin{corrige}[Exercice 3 (suite)]
\begin{enumerate}
\item \eng{depends of} $\rightarrow$ \eng{depends \textbf{on}}.
\item \eng{discussed about} $\rightarrow$ \eng{discussed} (sans préposition : \eng{discuss} est transitif direct).
\item \eng{good in} $\rightarrow$ \eng{good \textbf{at}} (\eng{to be good at something}).
\item \eng{entered into} $\rightarrow$ \eng{entered} (\eng{enter} est transitif direct ; on garde \eng{into} seulement avec un verbe de mouvement comme \eng{come}, \eng{go}, \eng{walk}).
\item \eng{in Saturdays} $\rightarrow$ \eng{\textbf{on} Saturdays} (jour $\rightarrow$ \eng{on}).
\end{enumerate}
\end{corrige}

\courssub{Exercice 4 : transformations}

\begin{exercice}[Transform the sentences]
\begin{enumerate}
\item Remplacez l'adverbe par un syntagme prépositionnel équivalent : \eng{She sang beautifully.} (utilisez \eng{with} + nom)
\item Remplacez : \eng{He drives his car to work.} (utilisez \eng{by})
\item Transformez en question : \eng{The keys are on the table.}
\item Transformez en question : \eng{They arrived at six o'clock.}
\item Remplacez l'adverbe : \eng{He spoke angrily to his son.} (utilisez \eng{in} + nom)
\end{enumerate}
\end{exercice}

\begin{corrige}[Exercice 4]
\begin{enumerate}
\item \eng{She sang with great beauty.} (ou mieux : \eng{She sang with a beautiful voice.})
\item \eng{He goes to work by car.}
\item \eng{Where are the keys?} ou \eng{Are the keys on the table?}
\item \eng{What time did they arrive?} ou \eng{Did they arrive at six o'clock?}
\item \eng{He spoke to his son in an angry voice.} (ou \eng{in anger})
\end{enumerate}
Remarquez la correspondance : adverbe de manière en \eng{-ly} $\leftrightarrow$ \eng{with} + nom ou \eng{in a\ldots way / voice / manner}. C'est une transformation classique au baccalauréat.
\end{corrige}

\courssub{Exercice 5 : traduction (français $\rightarrow$ anglais)}

\begin{exercice}[Translate into English]
Traduisez les phrases suivantes sur le thème de la vie familiale et sociale.

\begin{enumerate}
\item Ma grand-mère habite chez nous depuis dix ans.
\item Le mariage aura lieu à l'église, samedi, à dix heures.
\item Mon frère va au marché à pied tous les matins.
\item Nous sommes arrivés à Douala en train le mois dernier.
\item Elle découpe le pain avec un couteau et le pose sur la table.
\end{enumerate}
\end{exercice}

\begin{corrige}[Exercice 5]
\begin{enumerate}
\item \eng{My grandmother has lived with us for ten years.} (« chez nous » = \eng{with us} ; durée = \eng{for}.)
\item \eng{The wedding will take place at the church, on Saturday, at ten o'clock.}
\item \eng{My brother goes to the market on foot every morning.} (jamais \eng{by foot} !)
\item \eng{We arrived in Douala by train last month.} (\eng{arrive in} + ville.)
\item \eng{She cuts the bread with a knife and puts it on the table.}
\end{enumerate}
\end{corrige}

\courssub{Production finale : décrire sa maison, son quartier et son emploi du temps}

\begin{experience}[My home, my neighbourhood, my day]
Rédigez un paragraphe de 6 à 8 lignes en anglais. Vous devez :
\begin{itemize}
\item décrire votre maison (pièces, position des objets) ;
\item situer votre quartier (près de quoi ? en face de quoi ?) ;
\item raconter votre journée type (heures, moyens de transport, activités) ;
\item utiliser \textbf{au moins 8 prépositions différentes} et les souligner.
\end{itemize}
\end{experience}

\begin{exemplebox}[Modèle de production]
\eng{I live in a small house in Mendong, a quiet neighbourhood of Yaoundé. Our house is opposite a primary school and next to a provision store. In my bedroom, there is a bookshelf on the wall and my schoolbag always lies under my bed. I wake up at 5.30 a.m. and I leave home at 6.45. I go to school by taxi with my two sisters. Lessons begin at 7.30 and finish at 3 p.m. In the evening, I do my homework at the dining table, and on Saturdays, I help my mother at the market. I have lived in this neighbourhood since my childhood, and I feel at home here.}

« J'habite dans une petite maison à Mendong, un quartier calme de Yaoundé. Notre maison est en face d'une école primaire et à côté d'une boutique. Dans ma chambre, il y a une étagère au mur et mon cartable est toujours sous mon lit. Je me réveille à 5 h 30 et je quitte la maison à 6 h 45. Je vais à l'école en taxi avec mes deux sœurs. Les cours commencent à 7 h 30 et finissent à 15 h. Le soir, je fais mes devoirs à la table de la salle à manger, et le samedi, j'aide ma mère au marché. J'habite ce quartier depuis mon enfance et je m'y sens chez moi. »

Comptez les prépositions : \eng{in} (x5), \eng{of}, \eng{opposite}, \eng{next to}, \eng{on} (x2), \eng{under}, \eng{at} (x6), \eng{by}, \eng{with}, \eng{since}, \eng{to}. Soit onze prépositions différentes : objectif largement atteint.
\end{exemplebox}

\begin{attention}[La dernière passe de correction : la stratégie la plus rentable au Bac]
Dans une production écrite, ne relisez pas tout le texte en cherchant toutes les erreurs à la fois : vous n'y arriverez pas sous pression. Consacrez une \textbf{dernière passe de relecture uniquement aux prépositions}. Pour chaque préposition, posez-vous les trois questions de la méthode : fonction ? triade \eng{in/on/at} ? construction figée ? Cette stratégie ciblée élimine la majorité des fautes qui coûtent des points en \emph{essay writing} et en \emph{guided composition}. Les correcteurs du baccalauréat sont particulièrement sensibles aux erreurs du type \eng{married with}, \eng{in foot}, \eng{arrive to} : ce sont des erreurs « signature » des candidats francophones.
\end{attention}

\begin{aretenir}[Ce qu'il faut retenir de cette séance d'entraînement]
\begin{itemize}
\item Avant de choisir une préposition : identifier la fonction, appliquer \eng{in/on/at}, vérifier les constructions figées.
\item Position : \eng{at / in / on} ; mouvement : \eng{to / into / out of} ; durée : \eng{for / since} ; moyen : \eng{by} ; instrument : \eng{with}.
\item Les erreurs les plus coûteuses : \eng{married with}, \eng{listen} sans \eng{to}, \eng{arrive to}, \eng{by foot}, \eng{since} + durée.
\item En production écrite : une dernière relecture consacrée exclusivement aux prépositions.
\end{itemize}
\end{aretenir}

\newpage

\coursec{Activité d'intégration}

\coursec{Lesson 8 — Grammar: Prepositions and Prepositional Phrases (Place, Time, Manner and Direction)}

\courssub{Objectifs de la leçon}

À l'issue de cette leçon, l'élève sera capable de :

\begin{itemize}
\item identifier et classifier les \cle{prépositions} et \cle{syntagmes prépositionnels} (préposition + groupe nominal) selon leur fonction : lieu (\eng{place}), temps (\eng{time}), direction/mouvement (\eng{direction/movement}), manière (\eng{manner}) ;
\item construire des phrases correctes en respectant la syntaxe anglaise (position du syntagme prépositionnel, choix de la préposition selon le contexte) ;
\item comprendre des instructions, des indications de chemin ou des informations horaires dans un anglais standard (britannique) ;
\item produire à l'oral et à l'écrit une description de lieu, un itinéraire ou un récit en mobilisant au moins dix syntagmes prépositionnels variés ;
\item éviter les erreurs systématiques des francophones : confusion \eng{in/on/at}, \eng{to} superflu après \eng{arrive/reach/enter}, \eng{of} intrusif après \eng{near/opposite}, omission de la préposition verbale (\eng{listen to}, \eng{wait for}).
\end{itemize}

\courssub{Observation et découverte — Situation de communication}

\begin{textebox}[A letter from an English visitor]
“Dear students,

My name is Sarah Collins and I am a travel blogger from Manchester, England. Next month, I will spend three days in Cameroon to write an article about everyday life in a Cameroonian neighbourhood. I will arrive in Yaoundé on Saturday 15th June at 10 o'clock in the morning. I do not know the town at all, so I am looking for young guides!

I would like to visit a real neighbourhood, not the tourist places. I want to see the market, the school, the church or the mosque, the health centre and the football field. I would like to walk through the streets, sit in a small restaurant and talk to people. I am free from Saturday morning until Monday evening.

Could you be my guide for one afternoon? Please write back and tell me about your neighbourhood: where the important places are, how we can get there, and when they open.

Best regards,

Sarah Collins”
\end{textebox}

\begin{center}
\begin{tabularx}{\linewidth}{|X|X|}
\hline
\rowcolor{popblueL} \textbf{Mot anglais} & \textbf{Traduction française} \\
\hline
a travel blogger & un blogueur de voyage / une blogueuse de voyage \\
\hline
everyday life & la vie quotidienne \\
\hline
a neighbourhood & un quartier \\
\hline
a health centre & un centre de santé \\
\hline
to look for & chercher \\
\hline
to write back & répondre (par écrit) \\
\hline
a stall & un étal (marché) \\
\hline
a junction & un carrefour \\
\hline
\end{tabularx}
\end{center}

\courssub{Prépositions de lieu (Place) — Règle et emplois}

\begin{propriete}[Prepositions of Place : \eng{in}, \eng{on}, \eng{at} et prépositions de position]
\textbf{Règle générale (du grand au petit) :}
\begin{itemize}
\item \eng{IN} : espaces fermés, volumes, grandes zones (villes, pays, quartiers, pièces, parcs, rues \emph{au sens de « dans la rue »}). \eng{In Yaoundé}, \eng{in the market}, \eng{in the street}, \eng{in the classroom}, \eng{in Cameroon}.
\item \eng{ON} : surfaces, lignes, étages, rues \emph{au sens d'adresse}, moyens de transport publics. \eng{On the table}, \eng{on the wall}, \eng{on the first floor}, \eng{on Main Street}, \eng{on the bus}, \eng{on the train}, \eng{on foot}.
\item \eng{AT} : points précis, adresses exactes, événements, arrêts, coins. \eng{At the bus stop}, \eng{at 10 Main Street}, \eng{at the corner}, \eng{at the entrance}, \eng{at the party}, \eng{at home}, \eng{at school} (en tant qu'institution/activité).
\end{itemize}

\textbf{Prépositions de position relative (souvent composées) :}
\begin{center}
\begin{tabular}{|l|l|l|}
\hline
\rowcolor{popblueL} \textbf{Préposition} & \textbf{Sens} & \textbf{Exemple traduit} \\
\hline
\eng{opposite} & en face de & \eng{The bank is opposite the market.} « La banque est en face du marché. » \\
\hline
\eng{next to} / \eng{beside} & à côté de & \eng{The school is next to the church.} « L'école est à côté de l'église. » \\
\hline
\eng{between} & entre (deux) & \eng{The pharmacy is between the bakery and the tailor's shop.} « La pharmacie est entre la boulangerie et la couture. » \\
\hline
\eng{among} & au milieu de (plus de deux) & \eng{The well is among the houses.} « Le puits est au milieu des maisons. » \\
\hline
\eng{behind} & derrière & \eng{The football field is behind the health centre.} « Le terrain de foot est derrière le centre de santé. » \\
\hline
\eng{in front of} & devant (extérieur) & \eng{There is a taxi rank in front of the station.} « Il y a un parc de taxis devant la gare. » \\
\hline
\eng{near} / \eng{close to} & près de & \eng{We live near the university.} « Nous habitons près de l'université. » \\
\hline
\eng{at the corner of} & au coin de & \eng{Meet me at the corner of the street.} « Retrouve-moi au coin de la rue. » \\
\hline
\end{tabular}
\end{center}

\emph{Note :} \eng{opposite}, \eng{behind}, \eng{near} ne prennent \textbf{jamais} \eng{of}. \eng{Next to}, \eng{close to}, \eng{far from} intègrent la particule.
\end{propriete}

\courssub{Prépositions de temps (Time) — Règle et emplois}

\begin{propriete}[Prepositions of Time : \eng{at}, \eng{on}, \eng{in} et durées]
\textbf{Règle générale (du court au long) :}
\begin{itemize}
\item \eng{AT} : moment précis (heure, fête, \eng{night}, \eng{noon}, \eng{midnight}). \eng{At 10 o'clock}, \eng{at Christmas}, \eng{at night}, \eng{at the moment}.
\item \eng{ON} : jours, dates, jours + partie de journée, jours fériés précis. \eng{On Saturday}, \eng{on 15th June}, \eng{on Saturday morning}, \eng{on Christmas Day}, \eng{on my birthday}.
\item \eng{IN} : périodes longues (mois, années, saisons, siècles), parties de journée (sauf \eng{night}), délais futurs. \eng{In June}, \eng{in 2025}, \eng{in the morning}, \eng{in the afternoon}, \eng{in the evening}, \eng{in five minutes}, \eng{in two hours}.
\end{itemize}

\textbf{Autres prépositions temporelles fréquentes :}
\begin{center}
\begin{tabular}{|l|l|l|}
\hline
\rowcolor{popblueL} \textbf{Préposition} & \textbf{Emploi} & \textbf{Exemple traduit} \\
\hline
\eng{for} & durée (passé, présent, futur) & \eng{We walked for two hours.} « Nous avons marché pendant deux heures. » \\
\hline
\eng{since} & point de départ (present perfect) & \eng{She has lived here since 2020.} « Elle vit ici depuis 2020. » \\
\hline
\eng{from ... to / until / till} & intervalle (début $\rightarrow$ fin) & \eng{The centre opens from 8 a.m. to 4 p.m.} « Le centre ouvre de 8 h à 16 h. » \\
\hline
\eng{by} & échéance (au plus tard) & \eng{Be back by 6 p.m.} « Sois de retour pour 18 h au plus tard. » \\
\hline
\eng{during} & quand ? (dans un événement) & \eng{During the visit, it rained.} « Pendant la visite, il a plu. » \\
\hline
\end{tabular}
\end{center}

\emph{Attention :} Pas de préposition devant \eng{last}, \eng{next}, \eng{this}, \eng{every}, \eng{each} : \eng{last Saturday} (pas \emph{on last Saturday}), \eng{every morning} (pas \emph{in every morning}).
\end{propriete}

\courssub{Direction, mouvement et manière}

\begin{propriete}[Prepositions of Direction / Movement]
Contrairement au lieu (statique), ces prépositions indiquent un \textbf{trajet} ou un \textbf{but}. Elles s'utilisent avec des verbes de mouvement (\eng{go, walk, run, drive, fly, arrive, come, enter, cross, turn}).
\begin{center}
\begin{tabular}{|l|l|l|}
\hline
\rowcolor{popblueL} \textbf{Préposition} & \textbf{Sens / Image} & \textbf{Exemple traduit} \\
\hline
\eng{to} & but, destination (général) & \eng{We go to school.} « Nous allons à l'école. » \\
\hline
\eng{into} & entrée dans un volume & \eng{She walked into the shop.} « Elle est entrée dans la boutique. » \\
\hline
\eng{onto} & mouvement vers une surface & \eng{He jumped onto the stage.} « Il a sauté sur la scène. » \\
\hline
\eng{towards} & direction (orientation) & \eng{Walk towards the stadium.} « Marchez en direction du stade. » \\
\hline
\eng{along} & le long de (ligne) & \eng{We drove along the coast.} « Nous avons roulé le long de la côte. » \\
\hline
\eng{across} & traverser (surface/ligne) & \eng{Cross the road.} / \eng{Walk across the bridge.} « Traversez la route / le pont. » \\
\hline
\eng{through} & traverser (volume/espace 3D) & \eng{We walked through the forest.} « Nous avons traversé la forêt. » \\
\hline
\eng{past} & passer devant & \eng{Walk past the mosque.} « Passez devant la mosquée. » \\
\hline
\eng{round / around} & contourner / tout autour & \eng{Go round the building.} « Contournez le bâtiment. » \\
\hline
\eng{up / down} & monter / descendre (pente, rue) & \eng{Walk up the hill.} / \eng{Go down the street.} \\
\hline
\end{tabular}
\end{center}

\textbf{Verbes de mouvement sans préposition :} \eng{arrive} + \eng{at} (lieu précis) / \eng{in} (ville, pays) — \textbf{jamais} \emph{arrive to}. \eng{Reach}, \eng{enter} : transitifs directs (\eng{reach the village}, \eng{enter the house} — pas de \eng{to/into} après).
\end{propriete}

\begin{propriete}[Prepositions of Manner]
\begin{center}
\begin{tabular}{|l|l|l|}
\hline
\rowcolor{popblueL} \textbf{Préposition} & \textbf{Emploi} & \textbf{Exemple traduit} \\
\hline
\eng{by} + \eng{-ing} / nom (transport, communication) & moyen, instrument & \eng{by car, by bus, by bike, by bendskin, by phone, by email, by post} \\
\hline
\eng{on} + \eng{foot} (figé) & à pied & \eng{We go on foot.} « Nous y allons à pied. » \\
\hline
\eng{in} + \textit{mon véhicule personnel} & \eng{in my father's car}, \eng{in a taxi} \\
\hline
\eng{with} + nom (instrument, accompagnement, manière) & \eng{with a pen}, \eng{with pleasure}, \eng{with great care}, \eng{with my friends} \\
\hline
\eng{like} + nom / \eng{-ing} & comparaison (manière) & \eng{He runs like a cheetah.} « Il court comme un guépard. » \eng{She sings like an angel.} \\
\hline
\eng{without} + \eng{-ing} / nom & absence & \eng{He left without saying goodbye.} « Il est parti sans dire au revoir. » \\
\hline
\end{tabular}
\end{center}
\end{propriete}

\courssub{Comparaison français/anglais et pièges des francophones}

\begin{attention}[Erreurs fréquentes — Check-list de révision]
\begin{itemize}
\item \textbf{\eng{arrive}} : \eng{arrive \textbf{at} the station} (lieu précis), \eng{arrive \textbf{in} Yaoundé} (ville/pays). \textcolor{red}{Jamais \eng{arrive to}.}
\item \textbf{Transport} : \eng{by car/bus/bike/bendskin} MAIS \eng{in \textbf{my} car}, \eng{in \textbf{a} taxi}, \eng{on \textbf{the} bus}, \eng{on \textbf{the} train}, \eng{on foot} (sans article).
\item \textbf{Proximité} : \eng{near the school}, \eng{close to the school}. \textcolor{red}{Jamais \eng{near of}.}
\item \textbf{Position} : \eng{next \textbf{to}} (deux mots). \textcolor{red}{Jamais \eng{next of}.} \eng{opposite} (sans \eng{to}).
\item \textbf{Temps} : \eng{in the morning / afternoon / evening} MAIS \eng{on Monday morning}, \eng{on Tuesday afternoon}. \eng{at night} (pas \emph{in the night} sauf « au milieu de la nuit »).
\item \textbf{Verbes à préposition obligatoire} : \eng{listen \textbf{to}}, \eng{wait \textbf{for}}, \eng{look \textbf{at}} / \eng{look \textbf{for}}, \eng{talk \textbf{to}} / \eng{talk \textbf{about}}, \eng{ask \textbf{for}}, \eng{pay \textbf{for}}. En français, le verbe est souvent transitif direct ; en anglais, la préposition est requise.
\item \textbf{Adjectes + préposition} : \eng{good \textbf{at}}, \eng{interested \textbf{in}}, \eng{afraid \textbf{of}}, \eng{keen \textbf{on}}, \eng{responsible \textbf{for}}, \eng{married \textbf{to}}, \eng{similar \textbf{to}}, \eng{different \textbf{from}}.
\end{itemize}
\end{attention}

\begin{center}
\begin{tikzpicture}[mindmap, grow cyclic, every node/.style=concept, concept color=popblue,
    level 1/.append style={level distance=3.5cm, sibling angle=90},
    level 2/.append style={level distance=2.5cm, sibling angle=45}]
\node[concept, text=white, font=\bfseries, align=center]{Prepositional\\ Phrases}
    child[concept color=popgreen] { node[concept, align=center]{Place\\ (\eng{in/on/at}, \eng{next to}...)} }
    child[concept color=poporange] { node[concept, align=center]{Time\\ (\eng{at/on/in}, \eng{for/since}...)} }
    child[concept color=poppurple] { node[concept, align=center]{Direction\\ (\eng{to/into/towards}...)} }
    child[concept color=poppink] { node[concept, align=center]{Manner\\ (\eng{by/with/like}...)} };
\end{tikzpicture}
\legende{Cartographie mentale des quatre familles de prépositions étudiées.}
\end{center}

\courssub{Entraînement guidé et production — Situation d'intégration}

\courssub{Consignes de l'activité « Visite guidée pour Sarah »}

\textbf{Étape 1 — Préparation individuelle (10 min).} Dessinez rapidement un plan mental de votre quartier (Yaoundé, Douala, Bamenda, Buea, ou votre village). Notez au moins huit lieux en anglais : \eng{the market}, \eng{the school}, \eng{the church}, \eng{the mosque}, \eng{the health centre}, \eng{the football field}, \eng{the bus stop}, \eng{the bakery}, \eng{the tailor's shop}. Pour chaque lieu, écrivez un syntagme prépositionnel de lieu : \eng{opposite the market}, \eng{behind the school}, \eng{next to the mosque}.

\textbf{Étape 2 — Préparation du guide (5 min).} L'élève-guide prépare au moins \textbf{dix syntagmes prépositionnels variés} : quatre de lieu, trois de direction, deux de temps, un de manière. Exemples : \eng{We walk along this road.} « Nous marchons le long de cette route. » / \eng{The church is between the school and the health centre.} « L'église est entre l'école et le centre de santé. » / \eng{We will arrive at the football field in five minutes.} « Nous arriverons au terrain de foot dans cinq minutes. » / \eng{We go there on foot.} « Nous y allons à pied. »

\textbf{Étape 3 — Préparation du touriste (5 min).} L'élève-Sarah prépare \textbf{trois questions} : \eng{Where is the nearest pharmacy?} « Où est la pharmacie la plus proche ? » / \eng{How do we get to the market?} « Comment allons-nous au marché ? » / \eng{When does the health centre open?} « Quand le centre de santé ouvre-t-il ? »

\textbf{Étape 4 — Jeu de rôle (10 min par binôme).} Le guide accueille Sarah (\eng{Welcome to my neighbourhood!}) et conduit la visite. Sarah pose ses questions ; le guide répond avec des prépositions correctes. Le professeur circule et note les erreurs sans interrompre.

\textbf{Étape 5 — Inversion des rôles.} Les élèves échangent : le touriste devient guide de son propre quartier.

\textbf{Étape 6 — Production écrite individuelle (20 min).} Chaque élève rédige le compte rendu de la visite pour le blog de Sarah, en 8 à 10 lignes, à la troisième personne ou à la première, au prétérit ou au présent, en réutilisant \textbf{au moins dix syntagmes prépositionnels différents} (soulignez-les dans votre copie).

\textbf{Étape 7 — Correction et affichage.} Relecture croisée entre binômes (vérification des prépositions), correction collective au tableau des erreurs les plus fréquentes, affichage des meilleurs textes.

\courssub{Ressources à mobiliser}

\begin{aretenir}[Boîte à outils de la visite guidée]
\begin{itemize}
\item \textbf{Lieu} : \eng{opposite the market}, \eng{next to the school}, \eng{between the church and the mosque}, \eng{behind the health centre}, \eng{near the bus stop}, \eng{in front of the bakery}, \eng{at the corner of the street}, \eng{among the houses}.
\item \textbf{Direction et mouvement} : \eng{walk along the road}, \eng{go through the market}, \eng{turn left at the junction}, \eng{walk towards the stadium}, \eng{go across the bridge}, \eng{arrive at the football field}, \eng{get to the centre}, \eng{go past the mosque}, \eng{walk up the hill}.
\item \textbf{Temps} : \eng{at 10 o'clock}, \eng{in the morning}, \eng{on Saturday}, \eng{in five minutes}, \eng{from 8 a.m. until 4 p.m.}, \eng{for two hours}, \eng{since 2020}, \eng{by 6 p.m.}.
\item \textbf{Manière} : \eng{on foot}, \eng{by bike}, \eng{by bendskin}, \eng{with pleasure}, \eng{with great energy}, \eng{like a real guide}, \eng{without a map}.
\item \textbf{Questions utiles} : \eng{Where is...?}, \eng{How do we get to...?}, \eng{When does it open?}, \eng{How far is it from here?}, \eng{What time does it close?}
\end{itemize}
\end{aretenir}

\courssub{Proposition de production et corrigé commenté}

\begin{exoresolu}[A guided tour of Mokolo neighbourhood — Modèle de compte rendu]
Rédigez le compte rendu écrit de la visite (8 à 10 lignes) pour le blog de Sarah, avec au moins dix syntagmes prépositionnels soulignés.
\tcblower
\textbf{Production modèle :}

\begin{textebox}[Modèle de rédaction]
“On Saturday morning, I met Sarah \uline{at the bus station} \uline{at 10 o'clock}. First, we walked \uline{along the main road} \uline{towards Mokolo market}, which is \uline{opposite the post office}. We went \uline{through the market} and Sarah bought some fruit \uline{at a small stall} \uline{between two shops}. Then we visited the health centre, which is \uline{next to the primary school}, \uline{behind a big mango tree}. The centre opens \uline{at 7.30} \uline{in the morning} and closes \uline{at 4 p.m.} After that, we crossed the road and arrived \uline{at the football field} \uline{in five minutes}. Some boys were playing \uline{on the field} \uline{with great energy}. \uline{In the afternoon}, we had lunch \uline{in a small restaurant} \uline{near the church}, and we came back \uline{on foot} because Sarah wanted to take photos. She said my neighbourhood was \uline{like a village} \uline{inside the town}!”
\end{textebox}

\textbf{Commentaire des choix de langue (analyse grammaticale) :}

\begin{itemize}
\item \uline{\eng{On Saturday morning}} : préposition \eng{on} obligatoire car le jour est précisé (règle : \eng{on + day}), même si l'on dit \eng{in the morning}.
\item \uline{\eng{at the bus station at 10 o'clock}} : \eng{at} pour un lieu ponctuel (\eng{station}) et pour l'heure exacte (\eng{at 10 o'clock}).
\item \uline{\eng{along the main road towards Mokolo market}} : deux prépositions de direction bien distinguées : \eng{along} = le long de (suivi d'une ligne) ; \eng{towards} = en direction de (orientation vers un but).
\item \uline{\eng{opposite the post office}} : préposition de lieu sans \eng{to} (contrairement à \eng{next \textbf{to}}).
\item \uline{\eng{through the market}} : mouvement à l'intérieur d'un espace traversé (volume 3D), d'où \eng{through} et non \eng{across}.
\item \uline{\eng{at a small stall between two shops}} : \eng{at} pour un point précis (l'étal) ; \eng{between} pour la relation entre deux entités.
\item \uline{\eng{next to the primary school}}, \uline{\eng{behind a big mango tree}}, \uline{\eng{near the church}} : trois prépositions de position relative variées et correctes (pas de \eng{of} après \eng{near/behind}).
\item \uline{\eng{arrived at the football field in five minutes}} : \eng{arrive \textbf{at}} (lieu précis) — \textcolor{red}{jamais \eng{arrive to}} ; \eng{in five minutes} = délai futur (dans cinq minutes).
\item \uline{\eng{on the field with great energy}} : \eng{on} pour une surface de jeu (\eng{on the pitch/field/court}) ; \eng{with} + nom = manière.
\item \uline{\eng{In the afternoon}} : \eng{in} pour partie de journée (sauf \eng{at night}).
\item \uline{\eng{on foot}} : expression figée sans article (\textcolor{red}{pas \eng{by foot}, pas \eng{on the foot}}).
\item \uline{\eng{like a village inside the town}} : \eng{like} + nom = comparaison (manière) ; \eng{inside} = préposition de lieu (intérieur).
\end{itemize}

Le texte contient \textbf{18 syntagmes prépositionnels} soulignés couvrant les quatre catégories de la leçon : lieu (6), direction (3), temps (5), manière (4).
\end{exoresolu}

\courssub{Exercices d'application}

\begin{exercice}[Complétion — Choisissez la préposition correcte]
Complétez les phrases avec \eng{in}, \eng{on}, \eng{at}, \eng{to}, \eng{into}, \eng{onto}, \eng{through}, \eng{across}, \eng{along}, \eng{towards}, \eng{for}, \eng{since}, \eng{by}, \eng{with}, \eng{like}, \eng{opposite}, \eng{next to}, \eng{behind}, \eng{near}.
\begin{enumerate}
\item Sarah arrived \_\_\_\_\_ Yaoundé \_\_\_\_\_ Saturday \_\_\_\_\_ 10 a.m.
\item The pharmacy is \_\_\_\_\_ the bakery and the tailor's shop.
\item We walked \_\_\_\_\_ the bridge and then \_\_\_\_\_ the forest \_\_\_\_\_ the village.
\item She has lived in this neighbourhood \_\_\_\_\_ 2015.
\item They go to school \_\_\_\_\_ foot, but today they go \_\_\_\_\_ bendskin.
\item The meeting lasts \_\_\_\_\_ two hours, \_\_\_\_\_ 2 p.m. \_\_\_\_\_ 4 p.m.
\item He drives \_\_\_\_\_ his father's car. (Attention : véhicule personnel)
\item The children are playing \_\_\_\_\_ the field \_\_\_\_\_ great joy.
\item The mosque is \_\_\_\_\_ the health centre.
\item She speaks English \_\_\_\_\_ a native speaker.
\end{enumerate}
\end{exercice}

\begin{corrige}[Exercice 1]
\begin{enumerate}
\item \eng{in} (ville), \eng{on} (jour), \eng{at} (heure précise).
\item \eng{between} (deux repères).
\item \eng{across} (traverser surface/ligne), \eng{through} (traverser volume), \eng{into} (entrer dans) ou \eng{towards} (direction). Accepté : \eng{across the bridge, through the forest, towards/into the village}.
\item \eng{since} (point de départ + present perfect implicite).
\item \eng{on} (figé \eng{on foot}), \eng{by} (moyen de transport \eng{by bendskin}).
\item \eng{for} (durée), \eng{from}, \eng{to/until} (intervalle).
\item \eng{in} (véhicule personnel : \eng{in my father's car}).
\item \eng{on} (surface de jeu), \eng{with} (manière).
\item \eng{opposite} ou \eng{near} ou \eng{behind} (selon contexte imaginé). \eng{Opposite} sans \eng{to}.
\item \eng{like} (comparaison).
\end{enumerate}
\end{corrige}

\begin{exercice}[Transformation — Réécrivez en corrigeant l'erreur]
Chaque phrase contient une erreur de préposition (francophone). Réécrivez la phrase correcte.
\begin{enumerate}
\item We arrive to the station at noon.
\item The school is near of the market.
\item He goes to work by foot every day.
\item She is waiting the bus since 20 minutes.
\item They are interesting in Cameroonian culture.
\item We entered into the room quietly.
\item The match starts in Sunday afternoon.
\item He came in my car. (Contexte : c'est sa propre voiture)
\item The book is between the three students.
\item She looks the children playing.
\end{enumerate}
\end{exercice}

\begin{corrige}[Exercice 2]
\begin{enumerate}
\item We arrive \textbf{at} the station at noon. (\eng{arrive at/in}, jamais \eng{to}).
\item The school is \textbf{near} the market. (Pas \eng{of} après \eng{near}).
\item He goes to work \textbf{on} foot every day. (\eng{on foot}, pas \eng{by foot}).
\item She has been waiting \textbf{for} the bus \textbf{for} 20 minutes. (\eng{wait for} + \eng{for} durée / \eng{since} 20 minutes ago).
\item They are interested \textbf{in} Cameroonian culture. (\eng{interested in}).
\item They \textbf{entered} the room quietly. (\eng{enter} transitif direct, pas \eng{into}).
\item The match starts \textbf{on} Sunday afternoon. (\eng{on} + jour précis).
\item He came \textbf{in his} car. (Véhicule personnel $\rightarrow$ \eng{in} + possessif).
\item The book is \textbf{among} the three students. (\eng{among} pour $>2$).
\item She looks \textbf{at} the children playing. (\eng{look at} = regarder ; \eng{watch} possible sans préposition).
\end{enumerate}
\end{corrige}

\courssub{Bilan de la leçon}

Cette leçon a permis d'ancrer la maîtrise des \cle{syntagmes prépositionnels} à travers les quatre dimensions fondamentales : lieu, temps, direction, manière. La progression \textbf{observation $\rightarrow$ règle explicite $\rightarrow$ comparaison contrastive $\rightarrow$ entraînement guidé $\rightarrow$ production intégrée} respecte l'approche par les compétences du programme MINESEC.

\begin{itemize}
\item \textbf{Grammaire} : L'élève distingue désormais \eng{in/on/at} selon l'échelle spatiale/temporelle, évite \eng{arrive to} et \eng{near of}, et place correctement le syntagme prépositionnel en fin de phrase (structure SVO + Adverbial).
\item \textbf{Vocabulaire} : Le lexique du quartier (\eng{health centre, tailor's shop, stall, junction, bendskin}) et les verbes de mouvement (\eng{cross, walk along, turn, arrive}) sont réactivés en contexte.
\item \textbf{Compétences} : \textit{Listening/Reading} (lettre de Sarah), \textit{Speaking} (jeu de rôle guide/touriste), \textit{Writing} (compte rendu blog 8--10 lignes), \textit{Grammar} (exercices ciblés).
\item \textbf{Méthodologie} : L'oral précède l'écrit pour sécuriser la production ; la relecture par les pairs développe l'autonomie et l'esprit critique.
\end{itemize}

\emph{Poursuite possible :} Envoi des meilleurs comptes rendus à un établissement partenaire anglophone (projet eTwinning / correspondance) ou enregistrement audio des visites guidées pour le portfolio numérique de l'élève.

\newpage

\coursec{Méthodes et exercices résolus}

\coursec{Lesson 8 — Grammar: Prepositions and prepositional phrases (place, time, manner and direction)}

\begin{popfigure}
\begin{tikzpicture}[
  mindmap,
  concept color=popblue,
  level 1 concept/.append style={level distance=3.5cm, sibling angle=90},
  level 2 concept/.append style={level distance=2.5cm},
  every node/.style={font=\small, align=center},
  grow cyclic
]
\node[concept, align=center]{Prepositional\\ Phrases}
  child[concept color=popgreen] { node[concept, align=center]{Place\\ (Where?)}
    child { node[concept] {in (area)} }
    child { node[concept] {on (surface)} }
    child { node[concept] {at (point)} }
  }
  child[concept color=poporange] { node[concept, align=center]{Time\\ (When?)}
    child { node[concept] {in (month/year)} }
    child { node[concept] {on (day/date)} }
    child { node[concept] {at (hour)} }
  }
  child[concept color=poppurple] { node[concept, align=center]{Direction\\ (Where to?)}
    child { node[concept] {to / into / onto} }
    child { node[concept] {from / out of / off} }
    child { node[concept] {through / across / along} }
  }
  child[concept color=poppink] { node[concept, align=center]{Manner\\ (How?)}
    child { node[concept] {by + transport} }
    child { node[concept] {with + tool/person} }
    child { node[concept] {in + style/language} }
    child { node[concept] {like + comparison} }
  };
\end{tikzpicture}
\legende{Les quatre familles de phrases prépositionnelles : lieu, temps, direction, manière.}
\end{popfigure}

\courssub{Identifier une préposition et son complément dans un texte}

\begin{methode}[Repérer une phrase prépositionnelle en 4 étapes]
\begin{enumerate}
\item \textbf{Localise le mot de liaison.} Une préposition est un petit mot invariable (\eng{in, on, at, to, from, by, with, between, near, during...}) qui relie deux éléments de la phrase.
\item \textbf{Trouve son complément.} Toute préposition est suivie d'un groupe nominal (article + nom, pronom ou nom propre) : \eng{in the morning}, \eng{with her cousin}, \eng{at Mfoundi market}. L'ensemble \cle{préposition + complément} forme la \cle{prepositional phrase} (phrase prépositionnelle).
\item \textbf{Demande-toi ce qu'elle exprime.} Pose la question : \emph{où ?} (place), \emph{quand ?} (time), \emph{vers où ?} (direction), \emph{comment ?} (manner). Exemple : \eng{She arrived \textbf{at school} \textbf{by bus} \textbf{at 7 o'clock}.} → lieu, manière, temps.
\item \textbf{Vérifie la place dans la phrase.} La phrase prépositionnelle se place en général après le verbe ou en fin de phrase ; en tête de phrase, elle est suivie d'une virgule : \eng{In the evening, we visit our grandparents.}
\end{enumerate}
\end{methode}

\begin{exoresolu}[Repérage dans un texte]
Read the text and underline the prepositional phrases. Then say what each one expresses (place, time, manner or direction).

\eng{Every Saturday, my aunt travels from Bamenda to Yaoundé by bus. She leaves home at 5 a.m. and arrives at the station in the afternoon. During the journey, she talks with the other passengers.}
\tcblower
\textbf{Étape 1 — Repérage des prépositions} : \eng{from, to, by, at, at, in, during, with}.

\textbf{Étape 2 — Phrases prépositionnelles complètes} (préposition + complément) :
\begin{itemize}
\item \eng{from Bamenda} → \textbf{direction} (point de départ : \emph{d'où ?})
\item \eng{to Yaoundé} → \textbf{direction} (destination : \emph{vers où ?})
\item \eng{by bus} → \textbf{manner} (moyen de transport : \emph{comment ?})
\item \eng{at 5 a.m.} → \textbf{time} (heure précise : \emph{à quelle heure ?})
\item \eng{at the station} → \textbf{place} (lieu : \emph{où ?})
\item \eng{in the afternoon} → \textbf{time} (moment de la journée)
\item \eng{During the journey} → \textbf{time} (durée : \emph{pendant quoi ?})
\item \eng{with the other passengers} → \textbf{manner/accompagnement} (\emph{avec qui ?})
\end{itemize}

\textbf{Conclusion} : on retient le réflexe \eng{at} + heure, \eng{in} + moment de la journée, \eng{by} + moyen de transport sans article, \eng{from... to...} pour un trajet.
\end{exoresolu}

\courssub{Choisir la bonne préposition de lieu et de temps}

\begin{methode}[IN, ON, AT : la pyramide de la précision]
\begin{enumerate}
\item \textbf{Pour le lieu}, va du plus grand au plus précis : \eng{in} = espace fermé ou vaste (pays, ville, pièce, voiture) ; \eng{on} = surface ou ligne (rue, table, étage, bus) ; \eng{at} = point précis (adresse, arrêt, événement).
\item \textbf{Pour le temps}, même logique : \eng{in} = mois, années, saisons, moments de la journée (\eng{in July, in 2024, in the morning}) ; \eng{on} = jours et dates (\eng{on Monday, on 12th May, on Christmas Day}) ; \eng{at} = heures et moments ponctuels (\eng{at 6 o'clock, at noon, at night, at the weekend}).
\item \textbf{Pas de préposition} devant \eng{today, yesterday, tomorrow, tonight, last week, next month, this morning} : \eng{She came last Friday.} (jamais \eng{on last Friday}).
\item \textbf{Vérifie les cas fixes} : \eng{at home, at school, at work, in bed, on foot, on the radio, in the newspaper}. Ils s'apprennent par cœur.
\end{enumerate}
\end{methode}

\begin{propriete}[Tableau récapitulatif IN / ON / AT]
\begin{center}
\begin{tabularx}{\linewidth}{|X|X|X|}
\hline
\rowcolor{popblueL} \textbf{Français} & \textbf{English} & \textbf{Remarque} \\
\hline
En 2024, en juillet, en été, le matin, l'après-midi, le soir & \eng{in 2024, in July, in summer, in the morning, in the afternoon, in the evening} & Périodes longues / moments de la journée \\
\hline
Lundi, le 12 mai, le jour de Noël, le week-end (US) & \eng{on Monday, on 12th May, on Christmas Day, on the weekend} & Jours, dates précises \\
\hline
À 6 heures, à midi, à minuit, la nuit, le week-end (UK) & \eng{at 6 o'clock, at noon, at midnight, at night, at the weekend} & Heures exactes / moments ponctuels \\
\hline
À la maison, à l'école, au travail, au lit & \eng{at home, at school, at work, in bed} & Expressions figées (lieu) \\
\hline
À pied, à la radio, dans le journal & \eng{on foot, on the radio, in the newspaper} & Expressions figées (moyen/média) \\
\hline
Aujourd'hui, hier, demain, ce matin, la semaine dernière, le mois prochain & \eng{today, yesterday, tomorrow, this morning, last week, next month} & \textbf{Zéro préposition} (\eng{Ø}) \\
\hline
\end{tabularx}
\end{center}
\end{propriete}

\begin{exoresolu}[Complétion à trous type examen]
Complete the sentences with \eng{in}, \eng{on}, \eng{at} or nothing (Ø).

1. \eng{My cousin was born \_\_\_ 2010 \_\_\_ Douala.}
2. \eng{The wedding takes place \_\_\_ Saturday \_\_\_ 3 p.m.}
3. \eng{We always have lunch together \_\_\_ Christmas Day.}
4. \eng{There is a big market \_\_\_ Mokolo, \_\_\_ the centre of Yaoundé.}
5. \eng{She goes to the farm \_\_\_ foot every morning.}
6. \eng{The family meeting is \_\_\_ next Sunday.}
\tcblower
\textbf{Raisonnement phrase par phrase} :

1. \eng{in 2010} (année → \eng{in}) ; \eng{in Douala} (ville = espace vaste → \eng{in}). Réponse : \eng{My cousin was born \textbf{in} 2010 \textbf{in} Douala.}

2. \eng{on Saturday} (jour de la semaine → \eng{on}) ; \eng{at 3 p.m.} (heure précise → \eng{at}). Réponse : \eng{The wedding takes place \textbf{on} Saturday \textbf{at} 3 p.m.}

3. \eng{on Christmas Day} (date fêtée = jour précis → \eng{on}). Attention : on dit \eng{at Christmas} pour la période, mais \eng{on Christmas Day} pour le jour exact.

4. \eng{at Mokolo} ou \eng{in Mokolo} : un quartier précis fonctionne comme un point → \eng{at} est possible ; \eng{in the centre of Yaoundé} (espace délimité → \eng{in}). Réponse : \eng{There is a big market \textbf{at} Mokolo, \textbf{in} the centre of Yaoundé.}

5. \eng{on foot} : expression figée, c'est la seule exception à la règle \eng{by} + transport. Réponse : \eng{She goes to the farm \textbf{on} foot every morning.}

6. \eng{next Sunday} : devant \eng{next}, aucune préposition. Réponse : \eng{The family meeting is \textbf{Ø} next Sunday.}

\textbf{Conclusion} : la pyramide \eng{in > on > at} (du général au précis) résout 90 \% des cas ; le reste est à mémoriser.
\end{exoresolu}

\courssub{Exprimer la direction et le mouvement}

\begin{methode}[Prépositions de mouvement : le trajet en 3 points]
\begin{enumerate}
\item \textbf{Point de départ} : \eng{from} (\eng{She comes from Buea.}), \eng{out of} (sortir de), \eng{off} (descendre de, quitter une surface : \eng{He fell off his bike.}).
\item \textbf{Trajet} : \eng{across} (d'un côté à l'autre d'une surface : \eng{across the road}), \eng{through} (à travers un espace fermé : \eng{through the forest}), \eng{along} (le long de : \eng{along the Wouri river}), \eng{past} (devant, en passant : \eng{She walked past the school.}), \eng{over} (par-dessus), \eng{under} (en dessous).
\item \textbf{Destination} : \eng{to} (\eng{We are going to the stadium.}), \eng{into} (entrer dans : \eng{Come into the house.}), \eng{onto} (monter sur : \eng{The children climbed onto the lorry.}), \eng{towards} (en direction de, sans arriver forcément).
\item \textbf{Piège majeur} : \eng{to} marque la destination, \eng{at/in} marquent la position. Compare : \eng{She went \textbf{to} school} (mouvement) / \eng{She is \textbf{at} school} (position). De même : \eng{He jumped \textbf{into} the pool} / \eng{He is swimming \textbf{in} the pool}.
\item \textbf{Arriver} : toujours \eng{arrive \textbf{at}} + lieu précis, \eng{arrive \textbf{in}} + ville/pays. Jamais \eng{arrive to}.
\end{enumerate}
\end{methode}

\begin{exoresolu}[Choix entre position et mouvement]
Choose the correct preposition in brackets, then translate the sentence into French.

1. \eng{The bendskin rider drove (through / across) the tunnel.}
2. \eng{The ball flew (over / along) the fence into the neighbour's yard.}
3. \eng{My uncle arrived (to / in) Garoua last night.}
4. \eng{The cat jumped (onto / on) the table. Now it is sleeping (onto / on) it.}
5. \eng{We walked (along / through) the main road to the market.}
\tcblower
1. \eng{through} : un tunnel est un espace fermé que l'on traverse de part en part. \eng{The bendskin rider drove \textbf{through} the tunnel.} « Le conducteur de bendskin a traversé le tunnel. » (\eng{across} conviendrait pour une surface : \eng{across the bridge}.)

2. \eng{over} : le ballon passe par-dessus la clôture. \eng{The ball flew \textbf{over} the fence...} « Le ballon est passé par-dessus la clôture jusque dans la cour du voisin. »

3. \eng{in} : Garoua est une ville → \eng{arrive in}. \eng{My uncle arrived \textbf{in} Garoua last night.} « Mon oncle est arrivé à Garoua hier soir. » (\eng{arrive to} n'existe pas en anglais.)

4. \eng{onto} puis \eng{on} : le saut est un mouvement vers une surface (\eng{onto}) ; le sommeil est une position (\eng{on}). \eng{The cat jumped \textbf{onto} the table. Now it is sleeping \textbf{on} it.} « Le chat a sauté sur la table. Maintenant, il dort dessus. »

5. \eng{along} : on suit la route dans sa longueur. \eng{We walked \textbf{along} the main road to the market.} « Nous avons marché le long de la route principale jusqu'au marché. »

\textbf{Conclusion} : toujours se demander : \emph{mouvement ou position ?} — \eng{to/into/onto} pour le mouvement, \eng{at/in/on} pour la position.
\end{exoresolu}

\courssub{Exprimer la manière et éviter les calques du français}

\begin{methode}[La manière : BY, WITH, IN, LIKE]
\begin{enumerate}
\item \textbf{Moyen de transport ou de communication} : \eng{by} + nom sans article : \eng{by bus, by car, by plane, by phone, by email}. Exception unique : \eng{on foot}.
\item \textbf{Instrument ou accompagnement} : \eng{with} + article : \eng{She cut the cake with a knife.} / \eng{I went there with my brother.}
\item \textbf{Manière, style, langue} : \eng{in} : \eng{in a loud voice, in English, in a hurry, in a friendly way}.
\item \textbf{Comparaison} : \eng{like} + nom : \eng{He sings like a professional.} (jamais \eng{as} devant un simple nom pour « comme » de comparaison familière).
\item \textbf{Adverbe en -ly} : pour dire « comment », l'anglais préfère souvent un adverbe plutôt qu'une phrase prépositionnelle : \eng{She spoke politely} plutôt que \eng{in a polite way}. Formation : adjectif + \emph{-ly} (\eng{slowly, carefully}) ; attention aux adjectifs en \emph{-y} → \emph{-ily} (\eng{happy → happily}) et aux irréguliers (\eng{good → well}, \eng{fast → fast}, \eng{hard → hard}).
\end{enumerate}
\end{methode}

\begin{exoresolu}[Transformation : exprimer la manière autrement]
Rewrite each sentence using the word in brackets, without changing the meaning.

1. \eng{She answered in a polite way. (politely)}
2. \eng{He travelled to Douala using the train. (by)}
3. \eng{The old man walks with the help of a stick. (with)}
4. \eng{They welcomed us in a warm way. (warmly)}
5. \eng{My grandmother speaks Ewondo like a native. (fluently)}
\tcblower
\textbf{Principe} : l'anglais remplace volontiers la phrase prépositionnelle de manière (\eng{in a ... way}) par un adverbe en \emph{-ly}, placé après le verbe ou en fin de phrase.

1. \eng{polite → politely} (adjectif + \emph{-ly}). Réponse : \eng{She answered \textbf{politely}.} « Elle a répondu poliment. »

2. Moyen de transport → \eng{by} + nom sans article. Réponse : \eng{He travelled to Douala \textbf{by train}.} « Il a voyagé à Douala en train. » (Jamais \eng{by the train} pour le moyen général.)

3. Instrument → \eng{with} + article. Réponse : \eng{The old man walks \textbf{with a stick}.} « Le vieil homme marche avec une canne. »

4. \eng{warm → warmly}. Réponse : \eng{They welcomed us \textbf{warmly}.} « Ils nous ont accueillis chaleureusement. »

5. \eng{fluent → fluently} (adjectif en \emph{-ent} → \emph{-ently}). Réponse : \eng{My grandmother speaks Ewondo \textbf{fluently}.} « Ma grand-mère parle ewondo couramment. »

\textbf{Conclusion} : deux stratégies correctes — l'adverbe en \emph{-ly} (plus naturel) ou la phrase prépositionnelle (\eng{by / with / in / like}). Les deux sont acceptées à l'examen, à condition que la préposition soit exacte.
\end{exoresolu}

\courssub{Utiliser les prépositions dans un court texte de production}

\begin{methode}[Rédiger un paragraphe riche en phrases prépositionnelles]
\begin{enumerate}
\item \textbf{Planifie le squelette} : qui ? fait quoi ? (\eng{subject + verb}). Puis enrichis chaque phrase avec au moins deux phrases prépositionnelles : une de temps, une de lieu ou de direction.
\item \textbf{Ordre des compléments} : en anglais, l'ordre habituel est \cle{manière → lieu → temps} : \eng{We travelled \textbf{by bus} (manner) \textbf{to the village} (direction) \textbf{in the morning} (time).} C'est l'inverse du français (« le matin, nous sommes allés au village en car »).
\item \textbf{Varie les prépositions} : ne répète pas \eng{in} partout ; puise dans les quatre familles (lieu, temps, direction, manière).
\item \textbf{Relis avec la grille de contrôle} : \eng{at} + heure ? \eng{on} + jour ? \eng{in} + mois/année ? \eng{to} après \eng{go/come/travel} ? \eng{by} + transport sans article ? \eng{arrive at/in}, jamais \eng{arrive to} ?
\end{enumerate}
\end{methode}

\begin{exoresolu}[Production guidée type composition]
Write a short paragraph (60–80 words) about a family journey. Use at least six prepositional phrases: two of time, two of place or direction, one of manner and one of accompaniment.

\tcblower
\textbf{Étape 1 — Plan} : départ (heure, lieu) → trajet (manière, direction) → arrivée (lieu, temps) → activité (accompagnement).

\textbf{Étape 2 — Rédaction} :

\eng{Last December, my family left Yaoundé at 5 o'clock in the morning. We travelled by bus to our village in the West Region. During the journey, we sang songs with our cousins. We arrived at my grandmother's house in the afternoon, and in the evening we all sat around the fire and listened to her stories.}

« En décembre dernier, ma famille a quitté Yaoundé à 5 heures du matin. Nous avons voyagé en car jusqu'à notre village dans la région de l'Ouest. Pendant le trajet, nous avons chanté des chansons avec nos cousins. Nous sommes arrivés chez ma grand-mère dans l'après-midi, et le soir nous nous sommes tous assis autour du feu pour écouter ses histoires. »

\textbf{Étape 3 — Vérification des phrases prépositionnelles} :
\begin{itemize}
\item Temps : \eng{Last December} (sans préposition devant \eng{last} — correct), \eng{at 5 o'clock} (\eng{at} + heure), \eng{in the morning / in the afternoon / in the evening} (\eng{in} + moment de la journée), \eng{During the journey}.
\item Lieu/direction : \eng{to our village} (\eng{to} + destination), \eng{in the West Region}, \eng{at my grandmother's house} (\eng{at} + point précis), \eng{around the fire}.
\item Manière : \eng{by bus} (sans article).
\item Accompagnement : \eng{with our cousins}.
\end{itemize}

\textbf{Conclusion} : le paragraphe respecte l'ordre anglais manière → lieu → temps, compte 75 mots et contient onze phrases prépositionnelles correctes : objectif largement dépassé.
\end{exoresolu}

\begin{attention}[Les 5 erreurs qui coûtent des points au Bac]
\begin{enumerate}
\item \textbf{\eng{arrive to}} : calque du français « arriver à ». On dit toujours \eng{arrive \textbf{at} school}, \eng{arrive \textbf{in} Yaoundé}. Même piège avec \eng{listen} : \eng{listen \textbf{to} the radio} (la préposition est obligatoire en anglais, contrairement à « écouter la radio »).
\item \textbf{Préposition devant \eng{last, next, yesterday, tomorrow}} : \eng{on last Monday}, \eng{in next week} sont fautifs. On écrit directement \eng{last Monday}, \eng{next week}, \eng{yesterday evening}.
\item \textbf{Confusion \eng{in/on/at} de temps} : \eng{in Monday}, \eng{at July}, \eng{on 6 o'clock} sont des calques. Règle : \eng{at} + heure, \eng{on} + jour/date, \eng{in} + mois/année/saison. Et attention : \eng{at night} mais \eng{in the morning}.
\item \textbf{\eng{by foot} et \eng{with the car}} : le moyen de transport prend \eng{by} sans article (\eng{by car, by bus}), sauf l'expression figée \eng{\textbf{on} foot}. \eng{With} sert pour l'instrument ou l'accompagnement, pas pour le transport.
\item \textbf{Ordre des mots calqué sur le français} : « Nous sommes allés au village le matin en car » ne se traduit pas mot à mot. L'anglais place les compléments dans l'ordre \cle{manière → lieu → temps} : \eng{We went \textbf{by bus} \textbf{to the village} \textbf{in the morning}.} Un temps en tête de phrase exige une virgule : \eng{In the morning, we went to the village by bus.}
\end{enumerate}
\end{attention}

\newpage

\coursec{Exercices}

\courssub{Je vérifie mes connaissances}

\begin{exercice}[QCM : la bonne préposition]
Entoure la bonne réponse. Une seule réponse correcte par question.
\begin{enumerate}
\item The meeting is \_\_\_ Monday \_\_\_ 10 a.m.\\
a) at / in \hspace{1cm} b) on / at \hspace{1cm} c) in / on \hspace{1cm} d) on / in
\item My uncle has lived \_\_\_ Bamenda \_\_\_ 2015.\\
a) at / since \hspace{1cm} b) in / for \hspace{1cm} c) in / since \hspace{1cm} d) on / during
\item The children ran \_\_\_ the field and hid \_\_\_ a big tree.\\
a) across / behind \hspace{1cm} b) through / between \hspace{1cm} c) across / between \hspace{1cm} d) along / under
\item She goes to school \_\_\_ foot, but her brother goes \_\_\_ bendskin.\\
a) by / on \hspace{1cm} b) on / by \hspace{1cm} c) with / by \hspace{1cm} d) on / with
\item There is a pharmacy \_\_\_ the post office and the bank, just \_\_\_ the corner.\\
a) among / on \hspace{1cm} b) between / at \hspace{1cm} c) between / in \hspace{1cm} d) among / at
\end{enumerate}
\end{exercice}

\begin{exercice}[Vrai ou faux — justifie ta réponse]
Indique si chaque phrase est correcte ou incorrecte. Si elle est incorrecte, corrige-la et justifie brièvement.
\begin{enumerate}
\item \eng{I arrived to school late this morning.}
\item \eng{We have been waiting for the bus since twenty minutes.}
\item \eng{She entered the room without saying a word.}
\item \eng{The market is in front of the church, near to the main road.}
\item \eng{My grandmother was born on 1960.}
\item \eng{He walked through the crowd and sat down next to me.}
\end{enumerate}
\end{exercice}

\begin{exercice}[Vocabulaire : prépositions et sens]
Associe chaque préposition (1--8) à sa catégorie de sens (a--d). Certaines catégories contiennent plusieurs prépositions.
\begin{center}
\begin{tabular}{|l|l|}
\hline
\rowcolor{popblueL} \textbf{Prépositions} & \textbf{Catégories de sens} \\
\hline
1. \eng{towards} & a) lieu (place) \\
\hline
2. \eng{since} & b) temps (time) \\
\hline
3. \eng{behind} & c) direction / mouvement (direction) \\
\hline
4. \eng{without} & d) manière (manner) \\
\hline
5. \eng{into} & \\
\hline
6. \eng{until} & \\
\hline
7. \eng{with} & \\
\hline
8. \eng{underneath} & \\
\hline
\end{tabular}
\end{center}
\end{exercice}

\begin{exercice}[Complète avec la préposition qui convient]
Complète chaque phrase avec une préposition de ton choix, puis traduis la phrase en français.
\begin{enumerate}
\item My father works \_\_\_ a bank \_\_\_ the centre of Yaoundé.
\item We usually have lunch \_\_\_ noon, \_\_\_ the dining room.
\item The bus leaves \_\_\_ 6 a.m., so we must be \_\_\_ the station early.
\item She has been studying English \_\_\_ three years.
\item The cat jumped \_\_\_ the wall and ran \_\_\_ the garden.
\end{enumerate}
\end{exercice}

\courssub{Je m'entraîne}

\begin{exercice}[Texte à trous : une famille camerounaise]
Complète le paragraphe suivant avec les douze prépositions de la banque. Chaque préposition ne peut être utilisée qu'une seule fois.
\begin{center}
\textbf{Banque :} \eng{in -- on -- at -- to -- by -- with -- for -- since -- between -- across -- through -- without}
\end{center}
\begin{textebox}[The Ndi family's daily routine]
The Ndi family has lived \_\_\_ Douala \_\_\_ 2010. Their house is \_\_\_ a quiet street, \_\_\_ a school and a small market. Every morning, Mr Ndi leaves home \_\_\_ 6.30 and goes \_\_\_ work \_\_\_ taxi, because he cannot go there \_\_\_ his car: it is being repaired. His wife walks \_\_\_ the market \_\_\_ her basket and buys vegetables \_\_\_ the week. The children cross the road carefully and run \_\_\_ the school gate just before the bell rings. \_\_\_ the evening, the whole family sits together and talks about the day.
\end{textebox}
\end{exercice}

\begin{exercice}[Corrige les erreurs de francophones]
Chaque phrase contient une erreur de préposition typique d'un élève francophone. Relève l'erreur et réécris la phrase correctement.
\begin{enumerate}
\item \eng{I arrived to school late.}
\item \eng{She is waiting the bus at the junction.}
\item \eng{He entered into the room quietly.}
\item \eng{We go to the stadium by foot every Saturday.}
\item \eng{They are married since ten years.}
\item \eng{The teacher explained us the lesson.}
\item \eng{My house is near of the hospital.}
\item \eng{He listened the radio during the night.}
\end{enumerate}
\end{exercice}

\begin{exercice}[Traduction français $\rightarrow$ anglais]
Traduis le texte suivant en anglais. Fais attention aux prépositions de lieu, de temps et de direction.
\begin{quote}
Ma sœur habite à Douala depuis 2018. Elle travaille dans un hôpital, près du marché. Elle y va à pied tous les matins à 7 heures. Le soir, elle traverse le grand pont et rentre à la maison avant 6 heures. Le dimanche, elle rend visite à nos parents à Bafoussam et elle voyage en car.
\end{quote}
\end{exercice}

\begin{exercice}[Transformation : adverbe $\rightarrow$ syntagme prépositionnel]
Réécris chaque phrase en remplaçant l'adverbe souligné par un syntagme prépositionnel de manière, comme dans l'exemple.\\
\textbf{Exemple :} \eng{He left \textbf{hurriedly}.} $\rightarrow$ \eng{He left \textbf{in a hurry}.}
\begin{enumerate}
\item \eng{She spoke \textbf{angrily} to the driver.}
\item \eng{The students worked \textbf{silently} during the test.}
\item \eng{He answered the question \textbf{politely}.}
\item \eng{They welcomed the visitors \textbf{warmly}.}
\item \eng{The boy did his homework \textbf{carelessly}.}
\end{enumerate}
\end{exercice}

\courssub{Je me prépare au Bac}

\begin{exercice}[Compréhension et langue — type Bac]
Lis le texte suivant, puis réponds aux questions.
\begin{textebox}[Moving to the city]
When Amina's family moved from their village to Yaoundé in 2019, everything changed. In the village, their house stood between two mango trees, across the river from the primary school. Every morning, Amina walked through the fields with her cousins and arrived at school before seven o'clock.

In Yaoundé, life follows a different rhythm. The family now lives on the third floor of a block of flats near the Mokolo market. At half past six, Amina leaves home and waits at the bus stop with dozens of other students. The bus moves slowly through the traffic, and she sometimes arrives late. “In the village, I never saw a traffic jam,” she says with a smile. “Here, I spend an hour on the road every day.”

Yet Amina has adapted well. Since January, she has been attending an English club at her new school, and she has made friends from different regions of the country. On Saturdays, she helps her mother at the market, and on Sundays, the family visits relatives who live on the other side of the city.
\end{textebox}
\textbf{Comprehension questions}
\begin{enumerate}
\item Where did Amina's family live before 2019? (1 mark)
\item Give two differences between Amina's life in the village and her life in Yaoundé. (2 marks)
\item What does Amina do on Saturdays and on Sundays? (2 marks)
\end{enumerate}
\textbf{Language work}
\begin{enumerate}
\item Find in the text three prepositions of place and two prepositions of time, and use each of them in a sentence of your own. (5 marks)
\item Complete with the correct preposition: \eng{Amina has lived in Yaoundé \_\_\_ 2019. She waits \_\_\_ the bus \_\_\_ the bus stop.} (3 marks)
\item Correct the mistakes: \eng{She arrived to the club at Monday. She goes there with foot.} (2 marks)
\end{enumerate}
\end{exercice}

\begin{exercice}[Traduction — type Bac]
Traduis le texte suivant en anglais. Tu seras évalué(e) sur la correction des prépositions, des temps verbaux et du vocabulaire.
\begin{quote}
Depuis deux ans, mon frère vit chez mon oncle à Buea. Sa chambre se trouve au deuxième étage, entre la cuisine et le salon. Chaque matin, il part de la maison à six heures et demie et marche jusqu'à l'école, qui est en face de l'église. Le week-end, il joue au football avec ses amis sur le terrain derrière le marché. Il ne voyage jamais sans son téléphone.
\end{quote}
\end{exercice}

\begin{exercice}[Expression écrite — type Bac]
\textbf{Sujet :} \eng{Describe your home and your daily journey to school.}\\
Rédige un texte de \textbf{80 à 100 mots} en anglais. Tu décriras :
\begin{itemize}
\item l'endroit où tu habites (ville ou village, situation de la maison) ;
\item ton trajet quotidien pour aller à l'école (heure de départ, moyen de transport, itinéraire) ;
\item ce que tu fais en rentrant le soir.
\end{itemize}
\textbf{Critères d'évaluation :}
\begin{itemize}
\item variété et correction des prépositions de lieu, de temps, de direction et de manière (5 points) ;
\item correction des temps verbaux et de la syntaxe (3 points) ;
\item organisation et respect du nombre de mots (2 points).
\end{itemize}
\end{exercice}

\newpage

\coursec{Corrigés des exercices}

\begin{corrige}[Exercice 1 — QCM : la bonne préposition]
\begin{enumerate}
\item \textbf{b) on / at} \\
\eng{The meeting is on Monday at 10 a.m.} \\
Règle : \eng{on} pour les jours de la semaine / dates ; \eng{at} pour l'heure précise.
\item \textbf{c) in / since} \\
\eng{My uncle has lived in Bamenda since 2015.} \\
Règle : \eng{in} pour une ville ; \eng{since} + point de départ (date) avec present perfect.
\item \textbf{a) across / behind} \\
\eng{The children ran across the field and hid behind a big tree.} \\
\eng{across} = traverser une surface ; \eng{behind} = position (derrière).
\item \textbf{b) on / by} \\
\eng{She goes to school on foot, but her brother goes by bendskin.} \\
Idiomes : \eng{on foot} (à pied) ; \eng{by} + moyen de transport (by car, by bus, by bendskin).
\item \textbf{b) between / at} \\
\eng{There is a pharmacy between the post office and the bank, just at the corner.} \\
\eng{between} pour deux repères précis ; \eng{at the corner} (au coin de la rue).
\end{enumerate}
\end{corrige}

\begin{corrige}[Exercice 2 — Vrai ou faux — justifie ta réponse]
\begin{enumerate}
\item \textbf{Incorrect.} \eng{I arrived \textbf{at} school late this morning.} \\
Justification : le verbe \eng{arrive} s'emploie avec \eng{at} (lieu précis, bâtiment) ou \eng{in} (ville, pays), jamais \eng{to}.
\item \textbf{Incorrect.} \eng{We have been waiting for the bus \textbf{for} twenty minutes.} \\
Justification : \eng{for} + durée ; \eng{since} + point de départ (ex. \eng{since 10 a.m.}).
\item \textbf{Correct.} \eng{She entered the room without saying a word.} \\
Justification : \eng{enter} est transitif direct (pas de préposition) ; \eng{without} + \eng{-ing} correct.
\item \textbf{Incorrect.} \eng{The market is in front of the church, near the main road.} \\
Justification : \eng{near} (préposition) ne prend pas \eng{to} (cf. \eng{near to} \rightarrow \textbf{erreur de francophone}).
\item \textbf{Incorrect.} \eng{My grandmother was born \textbf{in} 1960.} \\
Justification : \eng{in} pour l'année (période longue) ; \eng{on} pour la date complète (ex. \eng{on 12 May 1960}).
\item \textbf{Correct.} \eng{He walked through the crowd and sat down next to me.} \\
Justification : \eng{through} (traverser un volume/groupe) ; \eng{next to} (à côté de) correct.
\end{enumerate}
\end{corrige}

\begin{corrige}[Exercice 3 — Vocabulaire : prépositions et sens]
\begin{center}
\begin{tabularx}{\linewidth}{|X|X|X|}
\hline
\rowcolor{popblueL} \textbf{Préposition} & \textbf{Catégorie} & \textbf{Exemple / Note} \\
\hline
\eng{towards} & c) direction / mouvement & \eng{He walked towards the station.} \\
\hline
\eng{since} & b) temps & \eng{I've known him since 2010.} \\
\hline
\eng{behind} & a) lieu (place) & \eng{The car is behind the house.} \\
\hline
\eng{without} & d) manière & \eng{She left without a word.} \\
\hline
\eng{into} & c) direction / mouvement & \eng{He came into the room.} (mouvement + entrée) \\
\hline
\eng{until} & b) temps & \eng{We waited until midnight.} \\
\hline
\eng{with} & d) manière (ou instrument) & \eng{He cut it with a knife.} / \eng{She spoke with confidence.} \\
\hline
\eng{underneath} & a) lieu (place) & \eng{The keys are underneath the mat.} (sous, dessous) \\
\hline
\end{tabularx}
\end{center}
\end{corrige}

\begin{corrige}[Exercice 4 — Complète avec la préposition qui convient]
\begin{enumerate}
\item \eng{My father works \textbf{in} a bank \textbf{in} the centre of Yaoundé.} \\
« Mon père travaille dans une banque au centre de Yaoundé. » (\eng{in} pour bâtiment/ville/quartier).
\item \eng{We usually have lunch \textbf{at} noon, \textbf{in} the dining room.} \\
« Nous déjeunons généralement à midi, dans la salle à manger. » (\eng{at} heure précise ; \eng{in} pièce close).
\item \eng{The bus leaves \textbf{at} 6 a.m., so we must be \textbf{at} the station early.} \\
« Le bus part à 6h, nous devons être à la gare tôt. » (\eng{at} heure / lieu précis).
\item \eng{She has been studying English \textbf{for} three years.} \\
« Elle étudie l'anglais depuis trois ans. » (\eng{for} + durée avec present perfect).
\item \eng{The cat jumped \textbf{over} the wall and ran \textbf{into} (ou \textbf{across} / \textbf{through}) the garden.} \\
« Le chat a sauté par-dessus le mur et a couru dans le jardin. » (\eng{over} = par-dessus ; \eng{into} = mouvement vers l'intérieur).
\end{enumerate}
\end{corrige}

\begin{corrige}[Exercice 5 — Texte à trous : une famille camerounaise]
\textbf{Texte complété :}
\begin{textebox}[The Ndi family's daily routine — Corrigé]
The Ndi family has lived \textbf{in} Douala \textbf{since} 2010. Their house is \textbf{on} a quiet street, \textbf{between} a school and a small market. Every morning, Mr Ndi leaves home \textbf{at} 6.30 and goes \textbf{to} work \textbf{by} taxi, because he cannot go there \textbf{without} his car: it is being repaired. His wife walks \textbf{to} the market \textbf{with} her basket and buys vegetables \textbf{for} the week. The children cross the road carefully and run \textbf{through} the school gate just before the bell rings. \textbf{In} the evening, the whole family sits together and talks about the day.
\end{textebox}

\textbf{Correspondances (banque de 12 prépositions) :}
\begin{itemize}
\item \eng{in} (Douala) — \eng{since} (2010) — \eng{on} (street) — \eng{between} (school and market)
\item \eng{at} (6.30) — \eng{to} (work) — \eng{by} (taxi) — \eng{without} (his car)
\item \eng{to} (market) — \eng{with} (basket) — \eng{for} (the week) — \eng{through} (gate) — \eng{in} (evening)
\end{itemize}
\end{corrige}

\begin{attention}[Contrainte « une seule fois »]
Le texte comporte 13 trous pour 12 prépositions uniques. Les prépositions \eng{in} et \eng{to} sont grammaticalement indispensables deux fois (\eng{in Douala} / \eng{In the evening} ; \eng{to work} / \eng{to the market}). La préposition \eng{across} n'est pas utilisée ici. En examen, on privilégie la correction grammaticale.


\end{attention}

\begin{corrige}[Exercice 6 — Corrige les erreurs de francophones]
\begin{enumerate}
\item \eng{I arrived \textbf{at} school late.} (arrive \textbf{at/in}, jamais \eng{to})
\item \eng{She is waiting \textbf{for} the bus at the junction.} (wait \textbf{for} qn/qch)
\item \eng{He entered the room quietly.} (enter = verbe transitif direct, pas de préposition)
\item \eng{We go to the stadium \textbf{on} foot every Saturday.} (idiome : \textbf{on} foot)
\item \eng{They have been married \textbf{for} ten years.} (married \textbf{for} + durée ; ou \eng{since} + date)
\item \eng{The teacher explained the lesson \textbf{to} us.} (explain sth \textbf{to} sb)
\item \eng{My house is near the hospital.} (near = préposition, pas \eng{near to})
\item \eng{He listened \textbf{to} the radio during the night.} (listen \textbf{to})
\end{enumerate}
\end{corrige}

\begin{corrige}[Exercice 7 — Traduction français $\rightarrow$ anglais]
\eng{My sister has lived in Douala since 2018. She works in a hospital, near the market. She goes there on foot every morning at 7 o'clock. In the evening, she crosses the big bridge and goes home before 6 o'clock. On Sundays, she visits our parents in Bafoussam and travels by coach.}

\textbf{Points de vigilance :}
\begin{itemize}
\item \eng{since 2018} (point de départ) + present perfect \eng{has lived}.
\item \eng{near the market} (sans \eng{to}).
\item \eng{on foot} (à pied), \eng{at 7 o'clock} (heure précise).
\item \eng{In the evening} (dans la soirée), \eng{crosses} (traverse), \eng{goes home} (rentre, sans \eng{to}).
\item \eng{On Sundays} (régularité), \eng{in Bafoussam} (ville), \eng{by coach} (en car).
\end{itemize}
\end{corrige}

\begin{corrige}[Exercice 8 — Transformation : adverbe $\rightarrow$ syntagme prépositionnel]
\begin{enumerate}
\item \eng{She spoke \textbf{in an angry manner} (ou \textbf{with anger}) to the driver.}
\item \eng{The students worked \textbf{in silence} (ou \textbf{without a sound}) during the test.}
\item \eng{He answered the question \textbf{in a polite way} (ou \textbf{with politeness}).}
\item \eng{They welcomed the visitors \textbf{in a warm way} (ou \textbf{with warmth}).}
\item \eng{The boy did his homework \textbf{in a careless way} (ou \textbf{without care}).}
\end{enumerate}
\end{corrige}

\begin{aretenir}[Structure]
Adverbe en \eng{-ly} $\rightarrow$ \eng{in a ... manner / way} ou \eng{with + nom abstrait} (souvent le même radical : \eng{angry} $\rightarrow$ \eng{anger}, \eng{polite} $\rightarrow$ \eng{politeness}, \eng{careless} $\rightarrow$ \eng{care}).


\end{aretenir}

\begin{corrige}[Exercice 9 — Compréhension et langue — type Bac]
\textbf{Comprehension questions (5 marks)}
\begin{enumerate}
\item Amina's family lived in a village (across the river from the primary school). (1 mark)
\item Différences (2 marks, 1 par différence) :
\begin{itemize}
\item Village : maison entre deux manguiers, trajet à pied à travers les champs, arrivée avant 7h, pas d'embouteillages.
\item Yaoundé : appartement au 3e étage près du marché Mokolo, trajet en bus, embouteillages, arrivée parfois en retard, une heure de route par jour.
\end{itemize}
\item Samedi : elle aide sa mère au marché. Dimanche : la famille rend visite à des parents de l'autre côté de la ville. (2 marks)
\end{enumerate}

\textbf{Language work (10 marks)}
\begin{enumerate}
\item \textbf{Prepositions de lieu (3 sur 5) :} \eng{between} (l.2), \eng{across} (l.2), \eng{near} (l.6), \eng{on} (l.6 — floor), \eng{at} (l.7 — bus stop), \eng{through} (l.8 — traffic), \eng{on} (l.9 — the road), \eng{on} (l.12 — the other side of).
\textbf{Prepositions de temps (2 sur 5) :} \eng{in} 2019 (l.1), \eng{before} seven o'clock (l.4), \eng{At} half past six (l.7), \eng{every day} (l.9), \eng{Since} January (l.10), \eng{On} Saturdays (l.11), \eng{on} Sundays (l.12).
Exemples de phrases personnelles (1 mark par phrase correcte) :
\eng{The bank is between the pharmacy and the post office.} / \eng{We arrived in Yaoundé in 2020.} / \eng{She waited at the bus stop for ten minutes.} / \eng{Since Monday, I have been very busy.} / \eng{On Sundays, we visit our grandparents.}
\item \eng{Amina has lived in Yaoundé \textbf{since} 2019. She waits \textbf{for} the bus \textbf{at} the bus stop.} (3 marks)
\item \eng{She arrived \textbf{at} the club \textbf{on} Monday. She goes there \textbf{on} foot.} (2 marks)
\end{enumerate}
\textbf{Barème indicatif :} Total 15 marks. Comprehension 5, Language 10.
\end{corrige}

\begin{corrige}[Exercice 10 — Traduction — type Bac]
\eng{For two years, my brother has been living with my uncle in Buea. His room is on the second floor, between the kitchen and the living room. Every morning, he leaves home at half past six and walks to school, which is opposite the church. At the weekend, he plays football with his friends on the pitch behind the market. He never travels without his phone.}

\textbf{Barème indicatif (10 points) :}
\begin{itemize}
\item Prépositions correctes : \eng{for two years} / \eng{since two years ago} (1), \eng{with} (1), \eng{in Buea} (1), \eng{on the second floor} (1), \eng{between ... and} (1), \eng{at half past six} (1), \eng{to school} (1), \eng{opposite} (1), \eng{at the weekend} / \eng{on weekends} (1), \eng{on the pitch} (1), \eng{behind} (1), \eng{without} (1) — \textit{sur 12 prépositions ciblées, notation sur 5 pts}.
\item Temps verbaux : \eng{has been living} (present perfect continuous, 1), \eng{is} (1), \eng{leaves/walks} (present simple habitude, 1), \eng{plays} (1), \eng{never travels} (1) — \textit{notation sur 3 pts}.
\item Vocabulaire / Syntaxe : \eng{pitch} / \eng{field}, \eng{opposite} / \eng{in front of}, \eng{living room}, \eng{coach} / \eng{bus} — \textit{notation sur 2 pts}.
\end{itemize}
\end{corrige}

\begin{corrige}[Exercice 11 — Expression écrite — type Bac]
\textbf{Proposition de rédaction modèle (environ 90 mots) :}
\begin{quote}
I live in a small house in Buea, at the foot of Mount Cameroon. My house is on a quiet street behind the Presbyterian Church, between a grocery shop and a primary school. Every weekday, I leave home at a quarter to seven and walk to school because it is only ten minutes away. I cross the main road at the traffic lights and go through the school gate just before the bell rings. In the evening, I return home at half past five, do my homework and help my mother prepare dinner before going to bed at ten o'clock.
\end{quote}

\textbf{Grille d'évaluation (10 points) :}
\begin{itemize}
\item \textbf{Prépositions (5 pts) :} Variété (lieu : \eng{in, at, on, behind, between, through, near/opposite} ; temps : \eng{at, in, before, after} ; direction : \eng{to, towards, across} ; manière : \eng{on foot, by...}). Correction : pas de \eng{to} après \eng{arrive/return/home}, \eng{on foot}, \eng{at} l'heure, \eng{in} la rue/ville.
\item \textbf{Grammaire / Syntaxe (3 pts) :} Present simple (habitudes), préterit si récit passé, word order (SVOMPT), accords, negation.
\item \textbf{Organisation / Nombre de mots (2 pts) :} Paragraphe unique ou deux, connecteurs (\eng{Every weekday, In the evening, because, just before}), 80–100 mots.
\end{itemize}
\end{corrige}

\newpage

\coursec{Fiche bilan}

\begin{popfigure}
\begin{tikzpicture}[mindmap, grow cyclic, every node/.style={concept, align=center, text=white, font=\small},
root concept/.append style={concept color=popink, font=\bfseries},
level 1/.append style={level distance=3.4cm, sibling angle=60},
level 2/.append style={level distance=2.2cm, font=\scriptsize}]
\node[root concept]{Prepositions and prepositional phrases}
child[concept color=popblue]{ node[align=center]{Place\\ \eng{in, on, at}}
  child { node[concept color=popblue!70] {\eng{in} = espace clos} }
  child { node[concept color=popblue!70] {\eng{on} = surface} }
  child { node[concept color=popblue!70] {\eng{at} = point} }
}
child[concept color=poporange]{ node[align=center]{Time\\ \eng{in, on, at}}
  child { node[concept color=poporange!70] {\eng{in} + mois, année} }
  child { node[concept color=poporange!70] {\eng{on} + jour, date} }
  child { node[concept color=poporange!70] {\eng{at} + heure} }
}
child[concept color=popgreen]{ node[align=center]{Direction\\ \eng{to, into, across, through}} }
child[concept color=poppurple]{ node[align=center]{Manner\\ \eng{by, with, in, like}} }
child[concept color=popgold]{ node[align=center]{No preposition\\ \eng{home, here, last week}} }
child[concept color=poppink]{ node[align=center]{Pièges\\ des francophones} };
\end{tikzpicture}
\legende{Carte mentale de la leçon : les prépositions de lieu, de temps, de direction et de manière.}
\end{popfigure}

\begin{bilan}
\textbf{1. Lexique clé (English / Français)}
\begin{itemize}
\item \eng{preposition} : préposition ; \eng{prepositional phrase} : syntagme prépositionnel (préposition + complément).
\item \eng{in front of} : devant ; \eng{behind} : derrière ; \eng{between} : entre (deux) ; \eng{among} : parmi (plusieurs) ; \eng{next to / beside} : à côté de ; \eng{opposite} : en face de ; \eng{near} : près de.
\item \eng{towards} : vers ; \eng{through} : à travers ; \eng{across} : en traversant (une surface) ; \eng{along} : le long de ; \eng{past} : devant (en passant) ; \eng{over / under} : au-dessus de / sous.
\item \eng{by bus} : en bus ; \eng{on foot} : à pied ; \eng{in a hurry} : pressé ; \eng{with care} : avec soin ; \eng{like a child} : comme un enfant.
\end{itemize}

\textbf{2. Règles à connaître par cœur}
\begin{center}
\begin{tabularx}{\linewidth}{|l|X|X|}
\hline
\rowcolor{popblueL} Préposition & Emploi & Exemple \\
\hline
\eng{in} (lieu) & pays, villes, pièces, espaces clos & \eng{in Cameroon, in the kitchen} \\
\hline
\eng{on} (lieu) & surfaces, rues, étages & \eng{on the table, on the first floor} \\
\hline
\eng{at} (lieu) & point précis, adresse, événement & \eng{at the bus stop, at school} \\
\hline
\eng{in} (temps) & mois, années, saisons, parties du jour & \eng{in July, in 2024, in the morning} \\
\hline
\eng{on} (temps) & jours et dates & \eng{on Monday, on 20 May} \\
\hline
\eng{at} (temps) & heures, \eng{night}, \eng{weekend} (GB) & \eng{at 7 o'clock, at night} \\
\hline
\eng{to / into / out of} & mouvement vers / entrée / sortie & \eng{go to Buea, get into the car} \\
\hline
\eng{by / with / like} & moyen, instrument, comparaison & \eng{by taxi, cut with a knife} \\
\hline
\end{tabularx}
\end{center}

\textbf{3. Méthodes express}
\begin{itemize}
\item Lieu : poser la question « espace clos, surface ou point ? » $\rightarrow$ \eng{in / on / at}.
\item Temps : retenir la pyramide \eng{in} (grand) $>$ \eng{on} (jour) $>$ \eng{at} (heure précise).
\item Zéro préposition devant \eng{home, here, there, abroad} et devant \eng{last, next, this, every} + nom de temps : \eng{She came home late.} / \eng{We met last week.}
\item Une préposition anglaise est toujours suivie d'un nom, d'un pronom (\eng{with me}) ou d'un verbe en \eng{-ing} : \eng{Thank you for coming.}
\end{itemize}
\end{bilan}

\begin{aretenir}[Checklist avant le Bac]
$\square$ Je distingue \eng{in / on / at} pour le lieu (espace clos, surface, point précis).\\
$\square$ Je distingue \eng{in / on / at} pour le temps (mois-année, jour-date, heure).\\
$\square$ Je sais qu'on dit \eng{at night} mais \eng{in the morning / afternoon / evening}.\\
$\square$ J'utilise \eng{to} pour la destination (\eng{go to Douala}) mais jamais devant \eng{home, here, there}.\\
$\square$ Je dis \eng{by bus / by car} (moyen de transport) mais \eng{on foot}.\\
$\square$ Je dis \eng{listen to}, \eng{wait for}, \eng{look at} : je connais les verbes à préposition obligatoire.\\
$\square$ Je ne mets pas de préposition devant \eng{last, next, this, every} (\eng{next Monday}, pas \eng{on next Monday}).\\
$\square$ Je sais que \eng{married to}, \eng{angry with}, \eng{good at}, \eng{interested in} exigent une préposition fixe.\\
$\square$ Je mets le verbe en \eng{-ing} après une préposition (\eng{before leaving}).\\
$\square$ Je sais placer la préposition en fin de question : \eng{Who are you talking to?}\\
$\square$ Je sais construire un syntagme prépositionnel : préposition + déterminant + nom (\eng{in the morning}).\\
$\square$ Je relis mes productions pour traquer les prépositions calquées du français (\eng{depend of} $\rightarrow$ \eng{depend on}).
\end{aretenir}

\findecours

\end{document}
